\documentclass[12pt,a4paper,openany]{book}
\usepackage[utf8]{inputenc}
\usepackage[T1]{fontenc}
\usepackage[french]{babel}
\usepackage[margin=2.5cm]{geometry}
\usepackage{amsmath,amssymb,amsthm,mathrsfs,mathtools}
\usepackage{tikz}
\usetikzlibrary{arrows.meta,calc,positioning,patterns,decorations.markings,decorations.pathmorphing}
\usepackage{pgfplots}
\pgfplotsset{compat=1.18}
\usepackage[most]{tcolorbox}
\tcbuselibrary{theorems,skins,breakable}
\usepackage[colorlinks=true,linkcolor=blue!60!black,citecolor=green!50!black]{hyperref}
\usepackage{makeidx}
\makeindex
\usepackage{enumitem}
\usepackage{fancyhdr}
\usepackage{caption}
\usepackage{booktabs}
\usepackage{array}
\usepackage{longtable}
\usepackage{minitoc}
\usepackage{microtype}
\usepackage{xcolor}

\pagestyle{fancy}
\fancyhf{}
\fancyhead[LE]{\leftmark}
\fancyhead[RO]{\rightmark}
\fancyfoot[C]{\thepage}

% Theorem environments — \newtheorem (amsthm), NOT \newtcbtheorem
\theoremstyle{definition}
\newtheorem{definition}{Définition}[chapter]
\newtheorem{exemple}[definition]{Exemple}
\newtheorem{exercice}{Exercice}[chapter]

\theoremstyle{plain}
\newtheorem{theoreme}[definition]{Théorème}
\newtheorem{proposition}[definition]{Proposition}
\newtheorem{lemme}[definition]{Lemme}
\newtheorem{corollaire}[definition]{Corollaire}

\theoremstyle{remark}
\newtheorem{remarque}[definition]{Remarque}
\newtheorem{notation}[definition]{Notation}

% tcolorbox wrapping (visual only)
\tcolorboxenvironment{definition}{enhanced,breakable,colback=blue!5,colframe=blue!70!black,fonttitle=\bfseries,before skip=10pt,after skip=10pt,left=8pt,right=8pt}
\tcolorboxenvironment{theoreme}{enhanced,breakable,colback=red!5,colframe=red!70!black,fonttitle=\bfseries,before skip=10pt,after skip=10pt,left=8pt,right=8pt}
\tcolorboxenvironment{proposition}{enhanced,breakable,colback=orange!5,colframe=orange!70!black,fonttitle=\bfseries,before skip=10pt,after skip=10pt,left=8pt,right=8pt}
\tcolorboxenvironment{lemme}{enhanced,breakable,colback=yellow!5,colframe=yellow!50!black,fonttitle=\bfseries,before skip=10pt,after skip=10pt,left=8pt,right=8pt}
\tcolorboxenvironment{corollaire}{enhanced,breakable,colback=green!5,colframe=green!50!black,fonttitle=\bfseries,before skip=10pt,after skip=10pt,left=8pt,right=8pt}
\tcolorboxenvironment{remarque}{enhanced,breakable,colback=gray!5,colframe=gray!50!black,before skip=8pt,after skip=8pt,left=6pt,right=6pt}
\tcolorboxenvironment{exemple}{enhanced,breakable,colback=purple!5,colframe=purple!50!black,before skip=8pt,after skip=8pt,left=6pt,right=6pt}

% Custom commands
\newcommand{\N}{\mathbb{N}}
\newcommand{\Z}{\mathbb{Z}}
\newcommand{\Q}{\mathbb{Q}}
\newcommand{\R}{\mathbb{R}}
\newcommand{\C}{\mathbb{C}}
\newcommand{\abs}[1]{\left|#1\right|}
\newcommand{\set}[1]{\left\{#1\right\}}
\newcommand{\powerset}{\mathcal{P}}
\DeclareMathOperator{\card}{card}
\DeclareMathOperator{\pgcd}{pgcd}
\newcommand{\barvee}{\overline{\vee}}

% ============================================================
\begin{document}
\dominitoc
\begin{titlepage}
\begin{tikzpicture}[remember picture, overlay]
  \fill[left color=blue!15!white, right color=blue!5!white]
    (current page.south west) rectangle (current page.north east);
  \fill[color=orange!70!yellow]
    ([yshift=-2cm]current page.north west) rectangle ([yshift=-2.3cm]current page.north east);
  \fill[color=blue!80!black, rounded corners=8pt]
    ([xshift=3cm, yshift=-4cm]current page.north west) rectangle
    ([xshift=-3cm, yshift=-10cm]current page.north east);
  \node[anchor=center, text=white, font=\Huge\bfseries]
    at ([yshift=-6cm]current page.north) {Math\'ematiques Discr\`etes};
  \node[anchor=center, text=orange!80!yellow, font=\Large]
    at ([yshift=-7.5cm]current page.north) {Notes de Cours};
  \node[anchor=center, text=white!80, font=\large]
    at ([yshift=-8.8cm]current page.north) {Licence L2 --- 2025--2026};
  \node[anchor=center, text=blue!80!black, font=\Large\itshape]
    at ([yshift=-12cm]current page.north) {Ya\'e Ulrich Gaba};
  \draw[orange!70!yellow, line width=2pt]
    ([xshift=5cm, yshift=-13.5cm]current page.north west) --
    ([xshift=-5cm, yshift=-13.5cm]current page.north east);
  \node[anchor=center, text width=12cm, align=center, text=blue!60!black, font=\itshape]
    at ([yshift=-16cm]current page.north) {<<\,Dieu a fait les nombres entiers ; tout le reste est l'\oe{}uvre de l'homme.\,>>\\[4pt]--- Leopold Kronecker};
  \node[anchor=south, text=gray, font=\small]
    at ([yshift=1.5cm]current page.south) {\today};
\end{tikzpicture}
\end{titlepage}
% ============================================================
%  PREFACE
% ============================================================
\chapter*{Pr\'eface}
\addcontentsline{toc}{chapter}{Pr\'eface}
\label{ch:preface}

Les \textbf{math\'ematiques discr\`etes}\index{mathematiques discretes@math\'ematiques discr\`etes} constituent le socle th\'eorique sur lequel repose une grande partie de l'informatique moderne. Contrairement \`a l'analyse, qui \'etudie des ph\'enom\`enes continus, les math\'ematiques discr\`etes s'int\'eressent aux structures finies ou d\'enombrables~: entiers, graphes, arbres, langages formels, automates et circuits logiques.

Ce cours a \'et\'e con\c{c}u pour les \'etudiants de premi\`ere et deuxi\`eme ann\'ee de licence (L1/L2) en math\'ematiques, en informatique ou dans toute formation scientifique o\`u ces notions interviennent. Aucun pr\'erequis n'est exig\'e au-del\`a d'une familiarit\'e raisonnable avec le calcul alg\'ebrique et les notions \'el\'ementaires d'ensembles.

\medskip
\noindent\textbf{Objectifs p\'edagogiques.}
\begin{itemize}[leftmargin=2em]
  \item Ma\^itriser le raisonnement logique et les principales techniques de d\'emonstration.
  \item Conna\^itre les fondements de la th\'eorie des ensembles, des relations et des fonctions.
  \item \'Etudier les structures combinatoires fondamentales~: d\'enombrements, permutations, combinaisons.
  \item Comprendre les concepts cl\'es de la th\'eorie des graphes.
  \item S'initier \`a l'arithm\'etique modulaire et \`a ses applications.
  \item D\'evelopper la rigueur math\'ematique n\'ecessaire \`a l'informatique th\'eorique.
\end{itemize}

\medskip
\noindent\textbf{Organisation du cours.}
Chaque chapitre suit une progression qui va des d\'efinitions et des th\'eor\`emes fondamentaux vers des exemples d\'etaill\'es et des exercices de difficult\'e croissante, signal\'ee par des \'etoiles~:
\begin{center}
\begin{tabular}{cl}
  $\bigstar$ & Exercice d'application directe \\
  $\bigstar\bigstar$ & Exercice de synth\`ese \\
  $\bigstar\bigstar\bigstar$ & Exercice d'approfondissement ou de recherche \\
\end{tabular}
\end{center}

Les encadr\'es color\'es distinguent visuellement les d\'efinitions (bleu), les th\'eor\`emes (rouge), les propositions (orange), les lemmes (jaune), les corollaires (vert), les remarques (gris) et les exemples (violet).

\bigskip
\noindent Nous esp\'erons que cet ouvrage accompagnera le lecteur dans sa d\'ecouverte d'un domaine aussi riche que stimulant.

\vfill
\hfill\textit{L'auteur}

% ============================================================
%  TABLE DES NOTATIONS
% ============================================================
\chapter*{Table des notations}
\addcontentsline{toc}{chapter}{Table des notations}
\label{ch:notations}

\begin{longtable}{@{} >{\raggedright\arraybackslash}p{4cm} p{9cm} @{}}
\toprule
\textbf{Notation} & \textbf{Signification} \\
\midrule
\endhead
$\N$ & Ensemble des entiers naturels $\set{0,1,2,\dots}$ \\[4pt]
$\N^*$ & $\N \setminus \set{0}$ \\[4pt]
$\Z$ & Ensemble des entiers relatifs \\[4pt]
$\Q$ & Ensemble des nombres rationnels \\[4pt]
$\R$ & Ensemble des nombres r\'eels \\[4pt]
$\C$ & Ensemble des nombres complexes \\[4pt]
$\powerset(E)$ & Ensemble des parties (ensemble puissance) de $E$ \\[4pt]
$\card(E)$, $\abs{E}$ & Cardinalit\'e de l'ensemble $E$ \\[4pt]
$\emptyset$ & Ensemble vide \\[4pt]
$\in$, $\notin$ & Appartenance, non-appartenance \\[4pt]
$\subset$, $\subseteq$ & Inclusion stricte, inclusion large \\[4pt]
$\cup$, $\cap$ & R\'eunion, intersection \\[4pt]
$\setminus$ & Diff\'erence ensembliste \\[4pt]
$A^c$, $\overline{A}$ & Compl\'ementaire de $A$ \\[4pt]
$A \times B$ & Produit cart\'esien \\[4pt]
$\neg p$, $\lnot p$ & N\'egation de $p$ \\[4pt]
$p \land q$ & Conjonction (<<~et~>>) \\[4pt]
$p \lor q$ & Disjonction (<<~ou~>>) \\[4pt]
$p \Rightarrow q$, $p \to q$ & Implication \\[4pt]
$p \Leftrightarrow q$, $p \leftrightarrow q$ & \'Equivalence logique \\[4pt]
$\forall$ & Quantificateur universel (<<~pour tout~>>) \\[4pt]
$\exists$ & Quantificateur existentiel (<<~il existe~>>) \\[4pt]
$\exists!$ & Quantificateur d'unicit\'e \\[4pt]
$n!$ & Factorielle de $n$ \\[4pt]
$\binom{n}{k}$ & Coefficient binomial <<~$n$ parmi $k$~>> \\[4pt]
$a \mid b$ & $a$ divise $b$ \\[4pt]
$a \equiv b \pmod{n}$ & Congruence modulo $n$ \\[4pt]
$\gcd(a,b)$, $\text{pgcd}(a,b)$ & Plus grand commun diviseur \\[4pt]
$\text{lcm}(a,b)$, $\text{ppcm}(a,b)$ & Plus petit commun multiple \\[4pt]
$\lfloor x \rfloor$ & Partie enti\`ere inf\'erieure \\[4pt]
$\lceil x \rceil$ & Partie enti\`ere sup\'erieure \\[4pt]
$\sum$, $\prod$ & Somme, produit \\[4pt]
$f \circ g$ & Composition de fonctions \\[4pt]
$G = (V,E)$ & Graphe de sommets $V$ et d'ar\^etes $E$ \\
\bottomrule
\end{longtable}

% ============================================================
%  TABLE OF CONTENTS
% ============================================================
\tableofcontents

% ############################################################
% ############################################################
%  CHAPITRE 1 : LOGIQUE PROPOSITIONNELLE ET PREDICATS
% ############################################################
% ############################################################
\chapter{Logique propositionnelle et pr\'edicats}
\label{ch:logique}
\index{logique propositionnelle}
\index{logique des predicats@logique des pr\'edicats}

\section*{Introduction}
\addcontentsline{toc}{section}{Introduction}

La logique math\'ematique\index{logique mathematique@logique math\'ematique} est le langage fondamental dans lequel s'expriment toutes les math\'ematiques. Elle fournit les r\`egles de raisonnement qui garantissent la validit\'e des d\'emonstrations. En informatique, la logique intervient aussi bien dans la conception de circuits num\'eriques que dans la v\'erification formelle de programmes.

Ce chapitre pr\'esente deux niveaux de logique~:
\begin{enumerate}[label=(\roman*)]
  \item la \textbf{logique propositionnelle}\index{logique propositionnelle}, qui manipule des propositions atomiques reli\'ees par des connecteurs~;
  \item la \textbf{logique des pr\'edicats}\index{logique des predicats@logique des pr\'edicats}, qui enrichit le formalisme pr\'ec\'edent avec des variables et des quantificateurs.
\end{enumerate}

% ============================================================
\section{Propositions et valeurs de v\'erit\'e}
\label{sec:propositions}
\index{proposition}

\begin{definition}[Proposition]\label{def:proposition}
  Une \textbf{proposition}\index{proposition} (ou \emph{\'enonc\'e}) est une phrase d\'eclarative qui poss\`ede exactement une \textbf{valeur de v\'erit\'e}\index{valeur de verite@valeur de v\'erit\'e}~: soit \textbf{vrai} (not\'e $V$ ou $1$), soit \textbf{faux} (not\'e $F$ ou $0$).
\end{definition}

\begin{exemple}[Propositions]\label{ex:propositions}
  Les phrases suivantes sont des propositions~:
  \begin{enumerate}[label=(\alph*)]
    \item <<~$2 + 3 = 5$~>> est une proposition vraie.
    \item <<~Tout nombre premier est impair~>> est une proposition fausse ($2$ est premier et pair).
    \item <<~Il existe une infinit\'e de nombres premiers~>> est une proposition vraie.
  \end{enumerate}
  En revanche, les phrases <<~Ferme la porte~!~>> et <<~$x + 1 = 3$~>> (sans pr\'eciser $x$) ne sont \emph{pas} des propositions.
\end{exemple}

\begin{notation}[Variables propositionnelles]\label{not:var-prop}
  On d\'esigne les propositions par des lettres minuscules $p, q, r, s, \dots$ appel\'ees \textbf{variables propositionnelles}\index{variable propositionnelle}.
\end{notation}

% ============================================================
\section{Connecteurs logiques}
\label{sec:connecteurs}
\index{connecteur logique}

\`A partir de propositions simples, on construit des propositions compos\'ees \`a l'aide de \textbf{connecteurs logiques}\index{connecteur logique}.

\begin{definition}[N\'egation]\label{def:negation}
  Soit $p$ une proposition. La \textbf{n\'egation}\index{negation@n\'egation} de $p$, not\'ee $\neg p$\index{$\neg$}, est la proposition dont la valeur de v\'erit\'e est l'oppos\'ee de celle de $p$.

  \medskip
  \begin{center}
  \begin{tabular}{c|c}
    \toprule
    $p$ & $\neg p$ \\ \midrule
    $V$ & $F$ \\
    $F$ & $V$ \\
    \bottomrule
  \end{tabular}
  \end{center}
\end{definition}

\begin{definition}[Conjonction]\label{def:conjonction}
  Soient $p$ et $q$ deux propositions. La \textbf{conjonction}\index{conjonction} de $p$ et $q$, not\'ee $p \land q$\index{$\land$}, est vraie si et seulement si $p$ et $q$ sont toutes deux vraies.

  \medskip
  \begin{center}
  \begin{tabular}{cc|c}
    \toprule
    $p$ & $q$ & $p \land q$ \\ \midrule
    $V$ & $V$ & $V$ \\
    $V$ & $F$ & $F$ \\
    $F$ & $V$ & $F$ \\
    $F$ & $F$ & $F$ \\
    \bottomrule
  \end{tabular}
  \end{center}
\end{definition}

\begin{definition}[Disjonction]\label{def:disjonction}
  Soient $p$ et $q$ deux propositions. La \textbf{disjonction}\index{disjonction} (ou inclusif) de $p$ et $q$, not\'ee $p \lor q$\index{$\lor$}, est vraie si et seulement si au moins l'une des deux propositions est vraie.

  \medskip
  \begin{center}
  \begin{tabular}{cc|c}
    \toprule
    $p$ & $q$ & $p \lor q$ \\ \midrule
    $V$ & $V$ & $V$ \\
    $V$ & $F$ & $V$ \\
    $F$ & $V$ & $V$ \\
    $F$ & $F$ & $F$ \\
    \bottomrule
  \end{tabular}
  \end{center}
\end{definition}

\begin{definition}[Implication]\label{def:implication}
  Soient $p$ et $q$ deux propositions. L'\textbf{implication}\index{implication} $p \to q$\index{$\to$} (lue <<~$p$ implique $q$~>> ou <<~si $p$ alors $q$~>>) est fausse uniquement lorsque $p$ est vraie et $q$ est fausse.

  \medskip
  \begin{center}
  \begin{tabular}{cc|c}
    \toprule
    $p$ & $q$ & $p \to q$ \\ \midrule
    $V$ & $V$ & $V$ \\
    $V$ & $F$ & $F$ \\
    $F$ & $V$ & $V$ \\
    $F$ & $F$ & $V$ \\
    \bottomrule
  \end{tabular}
  \end{center}

  On appelle $p$ l'\textbf{hypoth\`ese}\index{hypothese@hypoth\`ese} (ou ant\'ec\'edent) et $q$ la \textbf{conclusion}\index{conclusion} (ou cons\'equent).
\end{definition}

\begin{remarque}[Implication mat\'erielle]\label{rem:implication-materielle}
  En math\'ematiques, $p \to q$ est automatiquement vraie d\`es que $p$ est fausse. Cette convention, dite de l'\emph{implication mat\'erielle}\index{implication materielle@implication mat\'erielle}, peut sembler contre-intuitive mais est essentielle \`a la coh\'erence du syst\`eme logique. Par exemple, <<~si $1 = 2$ alors la Lune est un fromage~>> est une proposition vraie.
\end{remarque}

\begin{definition}[\'Equivalence logique (biconditionnelle)]\label{def:equivalence}
  Soient $p$ et $q$ deux propositions. La \textbf{biconditionnelle}\index{equivalence logique@\'equivalence logique}\index{biconditionnelle} $p \leftrightarrow q$\index{$\leftrightarrow$} est vraie si et seulement si $p$ et $q$ ont la m\^eme valeur de v\'erit\'e.

  \medskip
  \begin{center}
  \begin{tabular}{cc|c}
    \toprule
    $p$ & $q$ & $p \leftrightarrow q$ \\ \midrule
    $V$ & $V$ & $V$ \\
    $V$ & $F$ & $F$ \\
    $F$ & $V$ & $F$ \\
    $F$ & $F$ & $V$ \\
    \bottomrule
  \end{tabular}
  \end{center}

  On dit aussi que $p$ et $q$ sont \textbf{logiquement \'equivalentes} et on \'ecrit $p \equiv q$\index{$\equiv$} ou $p \Leftrightarrow q$.
\end{definition}

\begin{exemple}[Priorit\'e des connecteurs]\label{ex:priorite}
  Par convention, l'ordre de priorit\'e d\'ecroissante des connecteurs est~:
  \[
    \neg \quad > \quad \land \quad > \quad \lor \quad > \quad \to \quad > \quad \leftrightarrow.
  \]
  Ainsi, $\neg p \land q \to r$ se lit $((\neg p) \land q) \to r$.
\end{exemple}

% ============================================================
\section{Tables de v\'erit\'e}
\label{sec:tables-verite}
\index{table de verite@table de v\'erit\'e}

Une \textbf{table de v\'erit\'e}\index{table de verite@table de v\'erit\'e} \'enum\`ere syst\'ematiquement toutes les combinaisons possibles de valeurs de v\'erit\'e des variables propositionnelles et donne la valeur de v\'erit\'e de la proposition compos\'ee pour chacune d'elles.

\begin{exemple}[Table de v\'erit\'e compl\`ete]\label{ex:table-complete}
  Construisons la table de v\'erit\'e de $\varphi = (p \to q) \land (\neg q \to \neg p)$.
  \begin{center}
  \begin{tabular}{cc|c|c|c|c}
    \toprule
    $p$ & $q$ & $\neg p$ & $\neg q$ & $p \to q$ & $\neg q \to \neg p$ \\
    \midrule
    $V$ & $V$ & $F$ & $F$ & $V$ & $V$ \\
    $V$ & $F$ & $F$ & $V$ & $F$ & $F$ \\
    $F$ & $V$ & $V$ & $F$ & $V$ & $V$ \\
    $F$ & $F$ & $V$ & $V$ & $V$ & $V$ \\
    \bottomrule
  \end{tabular}
  \end{center}
  On observe que les colonnes de $p \to q$ et de $\neg q \to \neg p$ sont identiques~: il s'agit en fait de la \emph{contrapos\'ee}\index{contraposee@contrapos\'ee}, dont nous reparlerons au 
ef{ch:demonstration}.
\end{exemple}

% ============================================================
\section{Tautologies, contradictions et \'equivalences}
\label{sec:tautologies}

\begin{definition}[Tautologie et contradiction]\label{def:tautologie}
  \begin{enumerate}[label=(\roman*)]
    \item Une proposition compos\'ee qui est \textbf{toujours vraie}, quelle que soit la valeur de v\'erit\'e de ses variables, est appel\'ee une \textbf{tautologie}\index{tautologie}.
    \item Une proposition compos\'ee qui est \textbf{toujours fausse} est appel\'ee une \textbf{contradiction}\index{contradiction}.
    \item Une proposition qui n'est ni une tautologie ni une contradiction est dite \textbf{contingente}\index{contingence}.
  \end{enumerate}
\end{definition}

\begin{exemple}[Tautologie classique]\label{ex:tautologie}
  La proposition $p \lor \neg p$ est une tautologie (principe du tiers exclu\index{tiers exclu}). La proposition $p \land \neg p$ est une contradiction.
\end{exemple}

\begin{definition}[\'Equivalence logique]\label{def:equiv-logique}
  Deux propositions $\varphi$ et $\psi$ sont \textbf{logiquement \'equivalentes}\index{equivalence logique@\'equivalence logique}, not\'e $\varphi \equiv \psi$, si et seulement si $\varphi \leftrightarrow \psi$ est une tautologie, c'est-\`a-dire si $\varphi$ et $\psi$ ont la m\^eme valeur de v\'erit\'e pour toute affectation des variables.
\end{definition}

\begin{proposition}[\'Equivalences fondamentales]\label{prop:equiv-fondamentales}
  Pour toutes propositions $p$, $q$ et $r$~:
  \begin{enumerate}[label=(\roman*)]
    \item \textbf{Double n\'egation}\index{double negation@double n\'egation}~: $\neg(\neg p) \equiv p$.
    \item \textbf{Commutativit\'e}\index{commutativite@commutativit\'e}~: $p \land q \equiv q \land p$ \quad et \quad $p \lor q \equiv q \lor p$.
    \item \textbf{Associativit\'e}\index{associativite@associativit\'e}~: $(p \land q) \land r \equiv p \land (q \land r)$ \quad et \quad $(p \lor q) \lor r \equiv p \lor (q \lor r)$.
    \item \textbf{Distributivit\'e}\index{distributivite@distributivit\'e}~: $p \land (q \lor r) \equiv (p \land q) \lor (p \land r)$ \quad et \quad $p \lor (q \land r) \equiv (p \lor q) \land (p \lor r)$.
    \item \textbf{Idempotence}\index{idempotence}~: $p \land p \equiv p$ \quad et \quad $p \lor p \equiv p$.
    \item \textbf{Absorption}\index{absorption}~: $p \land (p \lor q) \equiv p$ \quad et \quad $p \lor (p \land q) \equiv p$.
    \item \textbf{\'El\'ements neutres}~: $p \land V \equiv p$, \quad $p \lor F \equiv p$.
    \item \textbf{\'El\'ements absorbants}~: $p \land F \equiv F$, \quad $p \lor V \equiv V$.
    \item \textbf{Implication}\index{implication}~: $p \to q \equiv \neg p \lor q$.
    \item \textbf{Biconditionnelle}~: $p \leftrightarrow q \equiv (p \to q) \land (q \to p)$.
  \end{enumerate}
\end{proposition}

\begin{proof}
  Chaque \'equivalence se v\'erifie par construction de la table de v\'erit\'e correspondante. \`A titre d'illustration, d\'emontrons (ix). La table de v\'erit\'e de $p \to q$ et de $\neg p \lor q$ est~:
  \begin{center}
  \begin{tabular}{cc|c|c|c}
    \toprule
    $p$ & $q$ & $\neg p$ & $p \to q$ & $\neg p \lor q$ \\ \midrule
    $V$ & $V$ & $F$ & $V$ & $V$ \\
    $V$ & $F$ & $F$ & $F$ & $F$ \\
    $F$ & $V$ & $V$ & $V$ & $V$ \\
    $F$ & $F$ & $V$ & $V$ & $V$ \\
    \bottomrule
  \end{tabular}
  \end{center}
  Les deux derni\`eres colonnes \'etant identiques, on conclut $p \to q \equiv \neg p \lor q$.
\end{proof}

% ============================================================
\section{Lois de De Morgan}
\label{sec:de-morgan}
\index{lois de De Morgan}

Les lois de De Morgan\index{De Morgan} sont parmi les \'equivalences les plus utiles en logique et en th\'eorie des ensembles.

\begin{theoreme}[Lois de De Morgan]\label{thm:de-morgan}
  Pour toutes propositions $p$ et $q$~:
  \begin{enumerate}[label=(\roman*)]
    \item $\neg(p \land q) \equiv \neg p \lor \neg q$,
    \item $\neg(p \lor q) \equiv \neg p \land \neg q$.
  \end{enumerate}
\end{theoreme}

\begin{proof}
  D\'emontrons (i) par table de v\'erit\'e.
  \begin{center}
  \begin{tabular}{cc|c|c|c|c|c}
    \toprule
    $p$ & $q$ & $p \land q$ & $\neg(p \land q)$ & $\neg p$ & $\neg q$ & $\neg p \lor \neg q$ \\
    \midrule
    $V$ & $V$ & $V$ & $F$ & $F$ & $F$ & $F$ \\
    $V$ & $F$ & $F$ & $V$ & $F$ & $V$ & $V$ \\
    $F$ & $V$ & $F$ & $V$ & $V$ & $F$ & $V$ \\
    $F$ & $F$ & $F$ & $V$ & $V$ & $V$ & $V$ \\
    \bottomrule
  \end{tabular}
  \end{center}
  Les colonnes $\neg(p \land q)$ et $\neg p \lor \neg q$ sont identiques, ce qui \'etablit (i).

  \medskip
  D\'emontrons (ii) par table de v\'erit\'e.
  \begin{center}
  \begin{tabular}{cc|c|c|c|c|c}
    \toprule
    $p$ & $q$ & $p \lor q$ & $\neg(p \lor q)$ & $\neg p$ & $\neg q$ & $\neg p \land \neg q$ \\
    \midrule
    $V$ & $V$ & $V$ & $F$ & $F$ & $F$ & $F$ \\
    $V$ & $F$ & $V$ & $F$ & $F$ & $V$ & $F$ \\
    $F$ & $V$ & $V$ & $F$ & $V$ & $F$ & $F$ \\
    $F$ & $F$ & $F$ & $V$ & $V$ & $V$ & $V$ \\
    \bottomrule
  \end{tabular}
  \end{center}
  Les colonnes $\neg(p \lor q)$ et $\neg p \land \neg q$ co\"incident, ce qui \'etablit (ii).
\end{proof}

\begin{remarque}[Extension aux connecteurs g\'en\'eralis\'es]\label{rem:de-morgan-gen}
  Les lois de De Morgan se g\'en\'eralisent \`a un nombre fini quelconque de propositions~:
  \[
    \neg\bigl(p_1 \land p_2 \land \cdots \land p_n\bigr)
    \equiv \neg p_1 \lor \neg p_2 \lor \cdots \lor \neg p_n,
  \]
  \[
    \neg\bigl(p_1 \lor p_2 \lor \cdots \lor p_n\bigr)
    \equiv \neg p_1 \land \neg p_2 \land \cdots \land \neg p_n.
  \]
\end{remarque}

% ============================================================
\section{Portes logiques --- une perspective informatique}
\label{sec:portes-logiques}
\index{porte logique}

Les connecteurs logiques sont mat\'erialis\'es dans les circuits \'electroniques sous forme de \textbf{portes logiques}\index{porte logique}. Voici les principales portes avec leurs symboles standard~:

\begin{center}
\begin{tikzpicture}[
    gate/.style={draw, minimum width=1.4cm, minimum height=1cm, thick},
    >=Stealth,
    scale=0.95, transform shape
  ]

  % --- NOT gate ---
  \node[gate, label={[font=\footnotesize]above:NON}] (not) at (0,0) {$\neg$};
  \draw[->] (-1.5,0) node[left]{$p$} -- (not.west);
  \draw[->] (not.east) -- (1.5,0) node[right]{$\neg p$};

  % --- AND gate ---
  \node[gate, label={[font=\footnotesize]above:ET}] (and) at (5,0) {$\land$};
  \draw[->] (3.5,0.25) node[left]{$p$} -- (and.west |- 0,0.25);
  \draw[->] (3.5,-0.25) node[left]{$q$} -- (and.west |- 0,-0.25);
  \draw[->] (and.east) -- (6.5,0) node[right]{$p \land q$};

  % --- OR gate ---
  \node[gate, label={[font=\footnotesize]above:OU}] (or) at (10,0) {$\lor$};
  \draw[->] (8.5,0.25) node[left]{$p$} -- (or.west |- 0,0.25);
  \draw[->] (8.5,-0.25) node[left]{$q$} -- (or.west |- 0,-0.25);
  \draw[->] (or.east) -- (11.5,0) node[right]{$p \lor q$};

  % --- XOR gate ---
  \node[gate, label={[font=\footnotesize]above:OU-X}] (xor) at (0,-3) {$\oplus$};
  \draw[->] (-1.5,-2.75) node[left]{$p$} -- (xor.west |- 0,-2.75);
  \draw[->] (-1.5,-3.25) node[left]{$q$} -- (xor.west |- 0,-3.25);
  \draw[->] (xor.east) -- (1.5,-3) node[right]{$p \oplus q$};

  % --- NAND gate ---
  \node[gate, label={[font=\footnotesize]above:NON-ET}] (nand) at (5,-3) {$\barwedge$};
  \draw[->] (3.5,-2.75) node[left]{$p$} -- (nand.west |- 0,-2.75);
  \draw[->] (3.5,-3.25) node[left]{$q$} -- (nand.west |- 0,-3.25);
  \draw[->] (nand.east) -- (6.5,-3) node[right]{$\neg(p \land q)$};

  % --- NOR gate ---
  \node[gate, label={[font=\footnotesize]above:NON-OU}] (nor) at (10,-3) {$\barvee$};
  \draw[->] (8.5,-2.75) node[left]{$p$} -- (nor.west |- 0,-2.75);
  \draw[->] (8.5,-3.25) node[left]{$q$} -- (nor.west |- 0,-3.25);
  \draw[->] (nor.east) -- (11.5,-3) node[right]{$\neg(p \lor q)$};
\end{tikzpicture}
\end{center}

\begin{remarque}[Universalit\'e de NAND]\label{rem:nand-universel}
  La porte NAND\index{NAND} est dite \emph{universelle}\index{porte universelle}~: tout circuit logique peut \^etre r\'ealis\'e en utilisant uniquement des portes NAND. De m\^eme pour la porte NOR\index{NOR}.
\end{remarque}

% ============================================================
\section{Logique des pr\'edicats}
\label{sec:predicats}
\index{predicat@pr\'edicat}

La logique propositionnelle ne permet pas d'exprimer des \'enonc\'es portant sur des variables, comme <<~pour tout entier $n$, $n^2 \geq 0$~>>. La \textbf{logique des pr\'edicats} enrichit le formalisme en introduisant des variables et des quantificateurs.

\begin{definition}[Pr\'edicat]\label{def:predicat}
  Un \textbf{pr\'edicat}\index{predicat@pr\'edicat} (ou \emph{fonction propositionnelle}) est une expression $P(x_1, x_2, \dots, x_n)$ contenant des variables libres, qui devient une proposition d\`es que l'on affecte des valeurs aux variables, \`a condition de pr\'eciser un \textbf{domaine de discours}\index{domaine de discours} (ou \textbf{univers}).
\end{definition}

\begin{exemple}[Pr\'edicats]\label{ex:predicats}
  \begin{enumerate}[label=(\alph*)]
    \item $P(x) :$ <<~$x$ est un nombre premier~>> est un pr\'edicat sur $\N$.
    Alors $P(7)$ est vrai et $P(10)$ est faux.
    \item $Q(x,y) :$ <<~$x + y = 10$~>> est un pr\'edicat \`a deux variables sur $\Z$.
    Alors $Q(3,7)$ est vrai et $Q(2,5)$ est faux.
  \end{enumerate}
\end{exemple}

% ============================================================
\subsection{Quantificateurs}
\label{subsec:quantificateurs}
\index{quantificateur}

\begin{definition}[Quantificateur universel]\label{def:quant-universel}
  Soit $P(x)$ un pr\'edicat de domaine $D$. L'\'enonc\'e
  \[
    \forall x \in D, \; P(x)
  \]
  (lu <<~pour tout $x$ dans $D$, $P(x)$~>>) est vrai si et seulement si $P(a)$ est vrai pour \textbf{chaque} \'el\'ement $a \in D$.
\end{definition}

\begin{definition}[Quantificateur existentiel]\label{def:quant-existentiel}
  Soit $P(x)$ un pr\'edicat de domaine $D$. L'\'enonc\'e
  \[
    \exists x \in D, \; P(x)
  \]
  (lu <<~il existe $x$ dans $D$ tel que $P(x)$~>>) est vrai si et seulement si $P(a)$ est vrai pour \textbf{au moins un} \'el\'ement $a \in D$.
\end{definition}

\begin{exemple}[Quantificateurs sur $\N$]\label{ex:quantificateurs}
  \begin{enumerate}[label=(\alph*)]
    \item $\forall n \in \N, \; n + 0 = n$ est vrai.
    \item $\exists n \in \N, \; n^2 = 2$ est faux (car $\sqrt{2} \notin \N$).
    \item $\exists n \in \N, \; n^2 = 4$ est vrai ($n = 2$).
  \end{enumerate}
\end{exemple}

\begin{definition}[Quantificateur d'unicit\'e]\label{def:quant-unicite}
  L'\'enonc\'e $\exists! \, x \in D, \; P(x)$\index{quantificateur d'unicite@quantificateur d'unicit\'e} (lu <<~il existe un unique $x$ dans $D$ tel que $P(x)$~>>) signifie qu'il existe exactement un \'el\'ement de $D$ satisfaisant $P$. Formellement~:
  \[
    \exists x \in D, \; P(x) \;\land\; \bigl(\forall y \in D, \; P(y) \to y = x\bigr).
  \]
\end{definition}

% ============================================================
\subsection{N\'egation des \'enonc\'es quantifi\'es}
\label{subsec:negation-quantificateurs}
\index{negation des quantificateurs@n\'egation des quantificateurs}

\begin{theoreme}[N\'egation des quantificateurs]\label{thm:neg-quantif}
  Soit $P(x)$ un pr\'edicat de domaine $D$. Alors~:
  \begin{enumerate}[label=(\roman*)]
    \item $\neg\bigl(\forall x \in D,\; P(x)\bigr) \;\equiv\; \exists x \in D, \; \neg P(x)$.
    \item $\neg\bigl(\exists x \in D,\; P(x)\bigr) \;\equiv\; \forall x \in D, \; \neg P(x)$.
  \end{enumerate}
\end{theoreme}

\begin{proof}
  (i) L'\'enonc\'e $\forall x \in D,\; P(x)$ est faux si et seulement s'il existe au moins un \'el\'ement $a \in D$ pour lequel $P(a)$ est faux, c'est-\`a-dire $\neg P(a)$ est vrai. Cela revient exactement \`a dire $\exists x \in D,\; \neg P(x)$.

  \smallskip
  (ii) L'\'enonc\'e $\exists x \in D,\; P(x)$ est faux si et seulement si aucun \'el\'ement de $D$ ne satisfait $P$, c'est-\`a-dire si $\neg P(a)$ est vrai pour tout $a \in D$. Cela donne $\forall x \in D,\; \neg P(x)$.
\end{proof}

\begin{exemple}[N\'egation d'un \'enonc\'e]\label{ex:negation-quantif}
  N\'egations la proposition <<~tout nombre r\'eel positif admet une racine carr\'ee~>>~:
  \[
    \neg\bigl(\forall x \in \R,\; x \geq 0 \to \exists y \in \R,\; y^2 = x\bigr)
    \;\equiv\; \exists x \in \R,\; x \geq 0 \;\land\; \forall y \in \R,\; y^2 \neq x.
  \]
  La n\'egation dit qu'il existe un r\'eel positif qui n'a pas de racine carr\'ee r\'eelle. (Cet \'enonc\'e est faux, car la proposition initiale est vraie.)
\end{exemple}

% ============================================================
\subsection{Quantificateurs embo\^it\'es}
\label{subsec:quantificateurs-emboites}
\index{quantificateurs emboites@quantificateurs embo\^it\'es}

Lorsque plusieurs quantificateurs interviennent, l'\textbf{ordre} dans lequel ils apparaissent est crucial.

\begin{proposition}[Commutation de quantificateurs de m\^eme type]\label{prop:commut-quantif}
  Pour tout pr\'edicat $P(x,y)$~:
  \begin{enumerate}[label=(\roman*)]
    \item $\forall x,\; \forall y,\; P(x,y) \;\equiv\; \forall y,\; \forall x,\; P(x,y)$.
    \item $\exists x,\; \exists y,\; P(x,y) \;\equiv\; \exists y,\; \exists x,\; P(x,y)$.
  \end{enumerate}
  En revanche, en g\'en\'eral~:
  \[
    \forall x,\; \exists y,\; P(x,y) \;\not\equiv\; \exists y,\; \forall x,\; P(x,y).
  \]
\end{proposition}

\begin{exemple}[Importance de l'ordre]\label{ex:ordre-quantif}
  Sur $\R$~:
  \begin{enumerate}[label=(\alph*)]
    \item $\forall x \in \R,\; \exists y \in \R,\; x + y = 0$ est \textbf{vrai} (prendre $y = -x$).
    \item $\exists y \in \R,\; \forall x \in \R,\; x + y = 0$ est \textbf{faux} (aucun $y$ fix\'e ne convient pour tout $x$).
  \end{enumerate}
\end{exemple}

\begin{remarque}[N\'egation de quantificateurs embo\^it\'es]\label{rem:neg-emboites}
  Pour nier un \'enonc\'e \`a quantificateurs multiples, on inverse chaque quantificateur de gauche \`a droite et on nie le pr\'edicat final. Par exemple~:
  \[
    \neg\bigl(\forall x,\; \exists y,\; P(x,y)\bigr)
    \;\equiv\; \exists x,\; \forall y,\; \neg P(x,y).
  \]
\end{remarque}

% ============================================================
\section{Exercices}
\label{sec:exo-logique}

\begin{exercice}[$\bigstar$]\label{exo:table-verite-1}
  Construire la table de v\'erit\'e de $(p \to q) \to (\neg q \to \neg p)$. Que constatez-vous~?
\end{exercice}

\begin{exercice}[$\bigstar$]\label{exo:equiv-simple}
  Montrer par table de v\'erit\'e que $p \to (q \to r) \equiv (p \land q) \to r$.
\end{exercice}

\begin{exercice}[$\bigstar$]\label{exo:negation-simple}
  \'Ecrire la n\'egation de chacune des propositions suivantes~:
  \begin{enumerate}[label=(\alph*)]
    \item $\forall n \in \N, \; n^2 \geq n$.
    \item $\exists x \in \R, \; x^2 + x + 1 = 0$.
    \item $\forall \varepsilon > 0, \; \exists N \in \N, \; \forall n \geq N, \; \abs{u_n - \ell} < \varepsilon$.
  \end{enumerate}
\end{exercice}

\begin{exercice}[$\bigstar\bigstar$]\label{exo:de-morgan-3}
  \'Etendre les lois de De Morgan \`a trois propositions. D\'emontrer que~:
  \[
    \neg(p \land q \land r) \equiv \neg p \lor \neg q \lor \neg r.
  \]
\end{exercice}

\begin{exercice}[$\bigstar\bigstar$]\label{exo:equivalences}
  Simplifier la proposition suivante en utilisant les \'equivalences logiques de la 
ef{prop:equiv-fondamentales}~:
  \[
    \bigl[(p \lor q) \land (p \lor \neg q)\bigr] \lor q.
  \]
\end{exercice}

\begin{exercice}[$\bigstar\bigstar$]\label{exo:quantif-ordre}
  On d\'efinit $P(x,y)$ comme <<~$x \leq y$~>> sur $\N$. D\'eterminer la valeur de v\'erit\'e de~:
  \begin{enumerate}[label=(\alph*)]
    \item $\forall x \in \N,\; \exists y \in \N,\; P(x,y)$.
    \item $\exists y \in \N,\; \forall x \in \N,\; P(x,y)$.
    \item $\forall x \in \N,\; \forall y \in \N,\; P(x,y) \lor P(y,x)$.
  \end{enumerate}
\end{exercice}

\begin{exercice}[$\bigstar\bigstar\bigstar$]\label{exo:systeme-complet}
  Un ensemble de connecteurs est dit \textbf{fonctionnellement complet} si toute fonction bool\'eenne peut s'exprimer uniquement \`a l'aide de ces connecteurs. Montrer que $\set{\neg, \land}$ est fonctionnellement complet.
  \textit{Indication~: exprimer $p \lor q$ en fonction de $\neg$ et $\land$.}
\end{exercice}

\begin{exercice}[$\bigstar\bigstar\bigstar$]\label{exo:nand-universel}
  Montrer que toute fonction bool\'eenne peut \^etre r\'ealis\'ee uniquement avec la porte NAND (not\'ee $\mid$, d\'efinie par $p \mid q \equiv \neg(p \land q)$).
  \textit{Indication~: exprimer $\neg p$ et $p \land q$ \`a l'aide de $\mid$.}
\end{exercice}

% ============================================================
\section*{R\'esum\'e du chapitre}
\addcontentsline{toc}{section}{R\'esum\'e du chapitre}

\begin{itemize}[leftmargin=2em]
  \item Une \textbf{proposition} est un \'enonc\'e vrai ou faux. On combine les propositions avec les connecteurs $\neg$, $\land$, $\lor$, $\to$, $\leftrightarrow$.
  \item Les \textbf{tables de v\'erit\'e} permettent de d\'eterminer la valeur de v\'erit\'e de toute proposition compos\'ee.
  \item Une \textbf{tautologie} est toujours vraie~; une \textbf{contradiction} est toujours fausse.
  \item Les \textbf{lois de De Morgan} relient n\'egation, conjonction et disjonction~: $\neg(p \land q) \equiv \neg p \lor \neg q$ et $\neg(p \lor q) \equiv \neg p \land \neg q$.
  \item Un \textbf{pr\'edicat} est un \'enonc\'e contenant des variables libres. Les \textbf{quantificateurs} $\forall$ et $\exists$ transforment un pr\'edicat en proposition.
  \item Pour nier un \'enonc\'e quantifi\'e, on \'echange $\forall$ et $\exists$ puis on nie le pr\'edicat.
  \item L'ordre des quantificateurs de types diff\'erents est crucial~: $\forall x\, \exists y$ ne signifie pas la m\^eme chose que $\exists y\, \forall x$.
\end{itemize}


% ############################################################
% ############################################################
%  CHAPITRE 2 : TECHNIQUES DE DEMONSTRATION
% ############################################################
% ############################################################
\chapter{Techniques de d\'emonstration}
\label{ch:demonstration}
\index{demonstration@d\'emonstration}
\index{preuve@preuve|see{d\'emonstration}}

\section*{Introduction}
\addcontentsline{toc}{section}{Introduction}

Savoir \'enoncer un th\'eor\`eme ne suffit pas~: il faut aussi \^etre capable de le \emph{d\'emontrer}. La d\'emonstration est le processus par lequel on \'etablit la v\'erit\'e d'un \'enonc\'e math\'ematique \`a partir d'axiomes, de d\'efinitions et de th\'eor\`emes d\'ej\`a \'etablis, en utilisant des r\`egles de d\'eduction logiquement valides.

Ce chapitre pr\'esente les techniques de d\'emonstration les plus courantes et les plus puissantes. La ma\^itrise de ces techniques est indispensable non seulement en math\'ematiques mais aussi en informatique, o\`u elles servent \`a prouver la correction d'algorithmes, la terminaison de programmes et la validit\'e de protocoles.

% ============================================================
\section{D\'emonstration directe}
\label{sec:preuve-directe}
\index{demonstration directe@d\'emonstration directe}

\begin{definition}[D\'emonstration directe]\label{def:preuve-directe}
  Une \textbf{d\'emonstration directe}\index{demonstration directe@d\'emonstration directe} d'un \'enonc\'e $p \to q$ consiste \`a supposer que $p$ est vrai, puis \`a en d\'eduire que $q$ est vrai par une cha\^ine d'implications logiques.
\end{definition}

\begin{theoreme}[Somme de deux nombres pairs]\label{thm:somme-pairs}
  La somme de deux nombres entiers pairs est paire.
\end{theoreme}

\begin{proof}
  Soient $a$ et $b$ deux entiers pairs\index{nombre pair}. Par d\'efinition, il existe $k, \ell \in \Z$ tels que $a = 2k$ et $b = 2\ell$. Alors~:
  \[
    a + b = 2k + 2\ell = 2(k + \ell).
  \]
  Comme $k + \ell \in \Z$, on conclut que $a + b$ est pair.
\end{proof}

\begin{theoreme}[Produit d'un pair et d'un entier]\label{thm:produit-pair}
  Si $a$ est un entier pair, alors pour tout entier $b$, le produit $ab$ est pair.
\end{theoreme}

\begin{proof}
  Soit $a$ un entier pair~: il existe $k \in \Z$ tel que $a = 2k$. Pour tout $b \in \Z$~:
  \[
    ab = (2k)b = 2(kb).
  \]
  Comme $kb \in \Z$, le produit $ab$ est pair.
\end{proof}

\begin{theoreme}[Carr\'e d'un impair]\label{thm:carre-impair}
  Le carr\'e de tout nombre impair est impair.
\end{theoreme}

\begin{proof}
  Soit $n$ un entier impair\index{nombre impair}. Il existe $k \in \Z$ tel que $n = 2k + 1$. Alors~:
  \[
    n^2 = (2k+1)^2 = 4k^2 + 4k + 1 = 2(2k^2 + 2k) + 1.
  \]
  Comme $2k^2 + 2k \in \Z$, l'entier $n^2$ est de la forme $2m + 1$ avec $m = 2k^2 + 2k$, donc $n^2$ est impair.
\end{proof}

\begin{exemple}[Divisibilit\'e]\label{ex:div-directe}
  Montrons directement que pour tout entier $n$, si $6 \mid n$ alors $3 \mid n$.

  \smallskip
  \noindent\textit{D\'emonstration.} Supposons $6 \mid n$\index{divisibilite@divisibilit\'e}. Il existe $k \in \Z$ tel que $n = 6k = 3(2k)$. Comme $2k \in \Z$, on a bien $3 \mid n$. \qed
\end{exemple}

% ============================================================
\section{D\'emonstration par l'absurde}
\label{sec:preuve-absurde}
\index{demonstration par l'absurde@d\'emonstration par l'absurde}
\index{reductio ad absurdum@\textit{reductio ad absurdum}|see{d\'emonstration par l'absurde}}

\begin{definition}[D\'emonstration par l'absurde]\label{def:preuve-absurde}
  Pour d\'emontrer un \'enonc\'e $\varphi$, une \textbf{d\'emonstration par l'absurde}\index{demonstration par l'absurde@d\'emonstration par l'absurde} consiste \`a supposer $\neg\varphi$, puis \`a en d\'eduire une contradiction (une proposition de la forme $\psi \land \neg\psi$). On conclut alors que $\varphi$ est vrai.

  \smallskip
  Ce raisonnement repose sur l'\'equivalence~: $\varphi \equiv \neg(\neg\varphi)$ et sur le fait qu'un \'enonc\'e dont la n\'egation m\`ene \`a une contradiction doit \^etre vrai.
\end{definition}

\begin{theoreme}[Irrationalit\'e de $\sqrt{2}$]\label{thm:sqrt2-irr}
  Le nombre $\sqrt{2}$\index{$\sqrt{2}$, irrationalite@$\sqrt{2}$, irrationalit\'e} est irrationnel.
\end{theoreme}

\begin{proof}
  Supposons, par l'absurde, que $\sqrt{2}$ est rationnel. Alors il existe $p, q \in \Z$ avec $q \neq 0$ et $\gcd(p,q) = 1$ (fraction irr\'eductible) tels que $\sqrt{2} = \dfrac{p}{q}$.

  En \'elevant au carr\'e~: $2 = \dfrac{p^2}{q^2}$, donc $p^2 = 2q^2$.

  Ainsi $p^2$ est pair, ce qui implique que $p$ est pair (car si $p$ \'etait impair, $p^2$ serait impair d'apr\`es le 
ef{thm:carre-impair}). On \'ecrit $p = 2k$ pour un certain $k \in \Z$.

  Substituant~: $(2k)^2 = 2q^2$, soit $4k^2 = 2q^2$, d'o\`u $q^2 = 2k^2$. Donc $q^2$ est pair, et par le m\^eme argument, $q$ est pair.

  Mais alors $p$ et $q$ sont tous deux pairs, ce qui contredit l'hypoth\`ese $\gcd(p,q) = 1$. Cette contradiction montre que $\sqrt{2}$ est irrationnel.
\end{proof}

\begin{theoreme}[Infinit\'e des nombres premiers]\label{thm:inf-premiers}
  Il existe une infinit\'e de nombres premiers\index{nombre premier}\index{nombres premiers, infinitude}.
\end{theoreme}

\begin{proof}
  Supposons, par l'absurde, qu'il n'existe qu'un nombre fini de nombres premiers~: $p_1, p_2, \dots, p_n$. Consid\'erons l'entier~:
  \[
    N = p_1 \cdot p_2 \cdots p_n + 1.
  \]
  Puisque $N > 1$, par le th\'eor\`eme fondamental de l'arithm\'etique, $N$ admet au moins un diviseur premier $p$. Or, pour tout $i \in \set{1, \dots, n}$, on a $N \equiv 1 \pmod{p_i}$, donc $p_i \nmid N$. Ainsi $p$ n'est aucun des $p_i$, ce qui contredit l'hypoth\`ese que $p_1, \dots, p_n$ sont \emph{tous} les nombres premiers.
\end{proof}

\begin{exemple}[Absence de plus petit rationnel strictement positif]\label{ex:pas-plus-petit-rationnel}
  Montrons par l'absurde qu'il n'existe pas de plus petit nombre rationnel strictement positif.

  \smallskip
  \noindent\textit{D\'emonstration.} Supposons qu'un tel nombre $r \in \Q$ existe, avec $r > 0$ et $r \leq s$ pour tout $s \in \Q$ avec $s > 0$. Consid\'erons $r' = \dfrac{r}{2}$. Alors $r' \in \Q$, $r' > 0$ et $r' < r$, ce qui contredit la minimalit\'e de $r$. \qed
\end{exemple}

% ============================================================
\section{D\'emonstration par contrapos\'ee}
\label{sec:preuve-contraposee}
\index{demonstration par contraposee@d\'emonstration par contrapos\'ee}
\index{contraposee@contrapos\'ee}

\begin{definition}[Contrapos\'ee]\label{def:contraposee}
  La \textbf{contrapos\'ee}\index{contraposee@contrapos\'ee} de l'implication $p \to q$ est l'implication $\neg q \to \neg p$. On a l'\'equivalence logique~:
  \[
    (p \to q) \;\equiv\; (\neg q \to \neg p).
  \]
  Une \textbf{d\'emonstration par contrapos\'ee} de $p \to q$ consiste donc \`a supposer $\neg q$ et \`a en d\'eduire $\neg p$.
\end{definition}

\begin{remarque}[Contrapos\'ee vs.\ absurde]\label{rem:contrap-vs-absurde}
  La preuve par contrapos\'ee est un cas particulier de la preuve par l'absurde (on suppose $p$ et $\neg q$, et on obtient une contradiction avec $p$). Cependant, la contrapos\'ee a l'avantage d'\^etre plus directe~: on part de $\neg q$ et on arrive \`a $\neg p$ sans chercher explicitement une contradiction.
\end{remarque}

\begin{theoreme}[Crit\`ere de parit\'e du carr\'e]\label{thm:carre-pair}
  Pour tout entier $n$, si $n^2$ est pair alors $n$ est pair.
\end{theoreme}

\begin{proof}
  Nous d\'emontrons la contrapos\'ee~: si $n$ est impair, alors $n^2$ est impair. C'est exactement le 
ef{thm:carre-impair}, qui a d\'ej\`a \'et\'e \'etabli par preuve directe.
\end{proof}

\begin{theoreme}[Condition de divisibilit\'e]\label{thm:div-contrap}
  Pour tous entiers $a$ et $b$, si $ab$ est impair alors $a$ est impair et $b$ est impair.
\end{theoreme}

\begin{proof}
  Par contrapos\'ee, supposons que $a$ est pair ou $b$ est pair (c'est-\`a-dire $\neg(a \text{ impair} \land b \text{ impair})$). Supposons, sans perte de g\'en\'eralit\'e, que $a$ est pair~: $a = 2k$ pour un certain $k \in \Z$. Alors $ab = 2kb = 2(kb)$, donc $ab$ est pair. On a bien montr\'e que $\neg(a \text{ impair} \land b \text{ impair}) \to ab \text{ pair}$.
\end{proof}

\begin{exemple}[Irrationnalit\'e par contrapos\'ee]\label{ex:contrap-irr}
  Montrons que pour tout entier positif $n$, si $\sqrt{n}$ est rationnel alors $\sqrt{n}$ est un entier.

  \smallskip
  \noindent\textit{D\'emonstration par contrapos\'ee.} On montre~: si $\sqrt{n}$ n'est pas un entier, alors $\sqrt{n}$ est irrationnel. Supposons $\sqrt{n} \notin \Z$ et supposons par l'absurde que $\sqrt{n} = \frac{p}{q}$ avec $q > 1$ et $\gcd(p,q) = 1$. Alors $n = \frac{p^2}{q^2}$, donc $nq^2 = p^2$. Soit $r$ un facteur premier de $q$~; alors $r \mid q^2$ donc $r \mid p^2$, ce qui implique $r \mid p$ (par le lemme d'Euclide\index{lemme d'Euclide}). Mais cela contredit $\gcd(p,q)=1$. \qed
\end{exemple}

% ============================================================
\section{R\'ecurrence math\'ematique}
\label{sec:recurrence}
\index{recurrence@r\'ecurrence}
\index{induction|see{r\'ecurrence}}

La r\'ecurrence est l'une des m\'ethodes de d\'emonstration les plus importantes en math\'ematiques discr\`etes. Elle exploite la structure bien ordonn\'ee de $\N$.

\subsection{R\'ecurrence simple (faible)}
\label{subsec:recurrence-simple}
\index{recurrence simple@r\'ecurrence simple}

\begin{theoreme}[Principe de r\'ecurrence simple]\label{thm:recurrence-simple}
  Soit $P(n)$ une propri\'et\'e d\'efinie pour tout $n \geq n_0$ ($n_0 \in \N$). Si~:
  \begin{enumerate}[label=(\roman*)]
    \item \textbf{Base}\index{base de recurrence@base de r\'ecurrence}~: $P(n_0)$ est vraie, et
    \item \textbf{H\'er\'edit\'e}\index{heredite@h\'er\'edit\'e}~: pour tout $n \geq n_0$, $P(n) \to P(n+1)$,
  \end{enumerate}
  alors $P(n)$ est vraie pour tout $n \geq n_0$.
\end{theoreme}

\begin{remarque}[Analogie du domino]\label{rem:domino}
  La r\'ecurrence s'apparente \`a une file infinie de dominos~: si le premier domino tombe (\textit{base}), et si chaque domino en tombant fait tomber le suivant (\textit{h\'er\'edit\'e}), alors tous les dominos tombent.
\end{remarque}

\begin{theoreme}[Somme des $n$ premiers entiers]\label{thm:somme-n}
  Pour tout $n \in \N^*$~:
  \[
    \sum_{k=1}^{n} k = \frac{n(n+1)}{2}.
  \]
\end{theoreme}

\begin{proof}
  D\'emonstration par r\'ecurrence sur $n$.

  \textbf{Base} ($n = 1$)~: $\displaystyle\sum_{k=1}^{1} k = 1 = \frac{1 \cdot 2}{2}$. \checkmark

  \textbf{H\'er\'edit\'e.} Supposons que la formule est vraie pour un certain $n \geq 1$ (\textbf{hypoth\`ese de r\'ecurrence}\index{hypothese de recurrence@hypoth\`ese de r\'ecurrence}). Montrons qu'elle est vraie pour $n+1$~:
  \begin{align*}
    \sum_{k=1}^{n+1} k &= \left(\sum_{k=1}^{n} k\right) + (n+1) \\
    &= \frac{n(n+1)}{2} + (n+1) \quad\text{(par hypoth\`ese de r\'ecurrence)} \\
    &= \frac{n(n+1) + 2(n+1)}{2} \\
    &= \frac{(n+1)(n+2)}{2}.
  \end{align*}
  C'est bien la formule au rang $n+1$. Par le principe de r\'ecurrence, la propri\'et\'e est vraie pour tout $n \geq 1$.
\end{proof}

\begin{theoreme}[Somme g\'eom\'etrique]\label{thm:somme-geo}
  Pour tout $r \neq 1$ et tout $n \in \N$~:
  \[
    \sum_{k=0}^{n} r^k = \frac{1 - r^{n+1}}{1 - r}.
  \]
\end{theoreme}

\begin{proof}
  Par r\'ecurrence sur $n$.

  \textbf{Base} ($n = 0$)~: $\displaystyle\sum_{k=0}^{0} r^k = 1 = \frac{1 - r}{1 - r}$. \checkmark

  \textbf{H\'er\'edit\'e.} Supposons la formule vraie au rang $n$. Alors~:
  \begin{align*}
    \sum_{k=0}^{n+1} r^k &= \left(\sum_{k=0}^{n} r^k\right) + r^{n+1} = \frac{1 - r^{n+1}}{1 - r} + r^{n+1} \\
    &= \frac{1 - r^{n+1} + r^{n+1}(1 - r)}{1 - r} = \frac{1 - r^{n+1} + r^{n+1} - r^{n+2}}{1 - r} \\
    &= \frac{1 - r^{n+2}}{1 - r}.
  \end{align*}
  C'est la formule au rang $n+1$, ce qui ach\`eve la r\'ecurrence.
\end{proof}

\begin{theoreme}[In\'egalit\'e de Bernoulli]\label{thm:bernoulli}
  Pour tout $x \geq -1$ et tout $n \in \N$~:
  \[
    (1 + x)^n \geq 1 + nx.
  \]
\end{theoreme}

\begin{proof}
  Par r\'ecurrence sur $n$.

  \textbf{Base} ($n = 0$)~: $(1+x)^0 = 1 \geq 1 + 0 = 1$. \checkmark

  \textbf{H\'er\'edit\'e.} Supposons $(1 + x)^n \geq 1 + nx$ pour un certain $n \geq 0$. Puisque $x \geq -1$, on a $1 + x \geq 0$, donc~:
  \begin{align*}
    (1 + x)^{n+1} &= (1 + x)^n \cdot (1 + x) \\
    &\geq (1 + nx)(1 + x) \quad\text{(par hypoth\`ese de r\'ecurrence et $1+x \geq 0$)} \\
    &= 1 + nx + x + nx^2 \\
    &= 1 + (n+1)x + nx^2 \\
    &\geq 1 + (n+1)x \quad\text{(car $nx^2 \geq 0$).}
  \end{align*}
  La propri\'et\'e est donc h\'er\'editaire.
\end{proof}

\subsection{R\'ecurrence forte}
\label{subsec:recurrence-forte}
\index{recurrence forte@r\'ecurrence forte}

\begin{theoreme}[Principe de r\'ecurrence forte]\label{thm:recurrence-forte}
  Soit $P(n)$ une propri\'et\'e d\'efinie pour tout $n \geq n_0$. Si~:
  \begin{enumerate}[label=(\roman*)]
    \item \textbf{Base}~: $P(n_0)$ est vraie (\'eventuellement plusieurs cas de base), et
    \item \textbf{H\'er\'edit\'e forte}~: pour tout $n \geq n_0$, si $P(k)$ est vraie pour tout $k$ avec $n_0 \leq k \leq n$, alors $P(n+1)$ est vraie,
  \end{enumerate}
  alors $P(n)$ est vraie pour tout $n \geq n_0$.
\end{theoreme}

\begin{remarque}[R\'ecurrence simple vs.\ forte]\label{rem:rec-simple-forte}
  En r\'ecurrence simple, l'hypoth\`ese au rang $n$ porte uniquement sur $P(n)$. En r\'ecurrence forte, elle porte sur $P(n_0), P(n_0+1), \dots, P(n)$ simultan\'ement. La r\'ecurrence forte est particuli\`erement utile lorsque $P(n+1)$ d\'epend de plusieurs valeurs ant\'erieures, et pas seulement de $P(n)$.
\end{remarque}

\begin{theoreme}[Th\'eor\`eme fondamental de l'arithm\'etique --- existence]\label{thm:tfa-existence}
  Tout entier $n \geq 2$ peut s'\'ecrire comme produit de nombres premiers\index{decomposition en facteurs premiers@d\'ecomposition en facteurs premiers}.
\end{theoreme}

\begin{proof}
  Par r\'ecurrence forte sur $n \geq 2$.

  \textbf{Base} ($n = 2$)~: $2$ est premier, donc $2$ est un produit de nombres premiers (un seul facteur). \checkmark

  \textbf{H\'er\'edit\'e forte.} Soit $n \geq 2$ et supposons que tout entier $k$ avec $2 \leq k \leq n$ s'\'ecrit comme produit de nombres premiers. Consid\'erons $n + 1$.
  \begin{itemize}
    \item Si $n + 1$ est premier, c'est un produit de nombres premiers. \checkmark
    \item Si $n + 1$ n'est pas premier, alors $n + 1 = ab$ avec $2 \leq a, b \leq n$. Par hypoth\`ese de r\'ecurrence forte, $a$ et $b$ s'\'ecrivent chacun comme produit de nombres premiers. Donc $n + 1 = ab$ est un produit de nombres premiers. \checkmark
  \end{itemize}
  Par le principe de r\'ecurrence forte, tout entier $n \geq 2$ est un produit de nombres premiers.
\end{proof}

\begin{exemple}[Suite de Fibonacci]\label{ex:fibonacci-rec-forte}
  Soit $(F_n)_{n \geq 0}$ la suite de Fibonacci\index{Fibonacci, suite de} d\'efinie par $F_0 = 0$, $F_1 = 1$ et $F_{n} = F_{n-1} + F_{n-2}$ pour $n \geq 2$. Montrons par r\'ecurrence forte que $F_n < 2^n$ pour tout $n \geq 0$.

  \smallskip
  \noindent\textit{D\'emonstration.}
  \textbf{Bases}~: $F_0 = 0 < 1 = 2^0$ et $F_1 = 1 < 2 = 2^1$. \checkmark

  \textbf{H\'er\'edit\'e forte.} Soit $n \geq 1$ et supposons $F_k < 2^k$ pour tout $0 \leq k \leq n$. Alors~:
  \[
    F_{n+1} = F_n + F_{n-1} < 2^n + 2^{n-1} = 2^{n-1}(2 + 1) = 3 \cdot 2^{n-1} < 4 \cdot 2^{n-1} = 2^{n+1}.
  \]
  Donc $F_{n+1} < 2^{n+1}$. \qed
\end{exemple}

% ============================================================
\subsection{R\'ecurrence structurelle (mention)}
\label{subsec:recurrence-structurelle}
\index{recurrence structurelle@r\'ecurrence structurelle}

\begin{remarque}[R\'ecurrence structurelle]\label{rem:rec-structurelle}
  La r\'ecurrence math\'ematique s'applique aux entiers naturels. En informatique, on utilise souvent la \textbf{r\'ecurrence structurelle}\index{recurrence structurelle@r\'ecurrence structurelle}, qui g\'en\'eralise le principe \`a des structures d\'efinies inductivement (listes, arbres, expressions, etc.). Le principe est le m\^eme~:
  \begin{enumerate}[label=(\roman*)]
    \item on montre la propri\'et\'e pour les cas de base de la structure~;
    \item on montre que si la propri\'et\'e est vraie pour les sous-structures, elle l'est aussi pour la structure construite \`a partir d'elles.
  \end{enumerate}
  Par exemple, pour prouver une propri\'et\'e sur les arbres binaires, on la d\'emontre pour l'arbre vide (base), puis on montre que si elle est vraie pour les sous-arbres gauche et droit, elle l'est pour l'arbre tout entier (h\'er\'edit\'e).
\end{remarque}

% ============================================================
\section{Principe des tiroirs de Dirichlet}
\label{sec:pigeonhole}
\index{principe des tiroirs}
\index{pigeonhole principle|see{principe des tiroirs}}
\index{Dirichlet, principe de}

\begin{theoreme}[Principe des tiroirs --- forme simple]\label{thm:pigeonhole-simple}
  Si $n + 1$ objets (ou plus) sont r\'epartis dans $n$ tiroirs, alors au moins un tiroir contient au moins deux objets.
\end{theoreme}

\begin{proof}
  Par l'absurde. Supposons que chacun des $n$ tiroirs contient au plus un objet. Alors le nombre total d'objets est au plus $n \cdot 1 = n$, ce qui contredit l'hypoth\`ese de disposer de $n+1$ objets (ou plus). Donc au moins un tiroir contient au moins deux objets.
\end{proof}

\begin{theoreme}[Principe des tiroirs --- forme g\'en\'eralis\'ee]\label{thm:pigeonhole-gen}
  Si $N$ objets sont r\'epartis dans $k$ tiroirs, alors au moins un tiroir contient au moins $\left\lceil \dfrac{N}{k} \right\rceil$ objets.
\end{theoreme}

\begin{proof}
  Par l'absurde. Supposons que chaque tiroir contient au plus $\left\lceil \frac{N}{k} \right\rceil - 1$ objets. Le nombre total d'objets serait au plus~:
  \[
    k \cdot \left(\left\lceil \frac{N}{k} \right\rceil - 1\right)
    < k \cdot \left(\frac{N}{k} + 1 - 1\right)
    = N,
  \]
  o\`u l'on a utilis\'e $\lceil x \rceil < x + 1$. Cela contredit l'hypoth\`ese de disposer de $N$ objets.
\end{proof}

\begin{exemple}[Anniversaires]\label{ex:anniversaires}
  Dans un groupe de $367$ personnes, au moins deux personnes partagent le m\^eme jour d'anniversaire. En effet, il y a au plus $366$ jours possibles dans une ann\'ee (en comptant le 29 f\'evrier), et $367 > 366$. Par le principe des tiroirs, au moins un <<~tiroir~>> (jour de l'ann\'ee) contient au moins deux personnes.
\end{exemple}

\begin{exemple}[Chaussettes]\label{ex:chaussettes}
  Un tiroir contient $10$ chaussettes rouges et $10$ chaussettes bleues m\'elang\'ees. Combien faut-il en tirer (dans l'obscurit\'e) pour \^etre s\^ur d'avoir au moins deux chaussettes de la m\^eme couleur~?

  \smallskip
  Les <<~tiroirs~>> sont les $2$ couleurs. Par le principe des tiroirs, il suffit de tirer $2 + 1 = 3$ chaussettes.
\end{exemple}

\begin{theoreme}[Application classique --- sous-ensemble de somme divisible]\label{thm:somme-divisible}
  Soit $n \geq 1$. Parmi $n$ entiers quelconques $a_1, a_2, \dots, a_n$, il existe un sous-ensemble contigu $\set{a_{i+1}, a_{i+2}, \dots, a_j}$ dont la somme est divisible par $n$.
\end{theoreme}

\begin{proof}
  Consid\'erons les $n$ sommes partielles\index{sommes partielles}~:
  \[
    S_k = a_1 + a_2 + \cdots + a_k, \quad k = 1, 2, \dots, n.
  \]
  Examinons les restes de $S_1, S_2, \dots, S_n$ modulo $n$. Il y a $n$ restes possibles~: $0, 1, \dots, n-1$.
  \begin{itemize}
    \item Si l'un des $S_k$ satisfait $S_k \equiv 0 \pmod{n}$, alors $a_1 + \cdots + a_k$ est divisible par $n$ et on a termin\'e.
    \item Sinon, les $n$ restes $S_1 \bmod n, \dots, S_n \bmod n$ prennent des valeurs dans $\set{1, 2, \dots, n-1}$ (seulement $n-1$ valeurs possibles). Par le principe des tiroirs, il existe $i < j$ tels que $S_i \equiv S_j \pmod{n}$. Alors~:
    \[
      S_j - S_i = a_{i+1} + a_{i+2} + \cdots + a_j \equiv 0 \pmod{n}. \qedhere
    \]
  \end{itemize}
\end{proof}

\begin{exemple}[Points dans un carr\'e]\label{ex:points-carre}
  On place $5$ points dans un carr\'e de c\^ot\'e $1$. Montrons que deux d'entre eux sont \`a distance au plus $\frac{\sqrt{2}}{2}$.

  \smallskip
  \noindent\textit{D\'emonstration.} D\'ecoupons le carr\'e en $4$ petits carr\'es de c\^ot\'e $\frac{1}{2}$.

  \begin{center}
  \begin{tikzpicture}[scale=2]
    \draw[thick] (0,0) rectangle (1,1);
    \draw[dashed] (0.5,0) -- (0.5,1);
    \draw[dashed] (0,0.5) -- (1,0.5);
    \node at (0.25,0.75) {\small I};
    \node at (0.75,0.75) {\small II};
    \node at (0.25,0.25) {\small III};
    \node at (0.75,0.25) {\small IV};
    % Points
    \fill[blue] (0.15,0.85) circle (1pt);
    \fill[blue] (0.35,0.6) circle (1pt);
    \fill[blue] (0.8,0.9) circle (1pt);
    \fill[blue] (0.6,0.3) circle (1pt);
    \fill[blue] (0.1,0.2) circle (1pt);
  \end{tikzpicture}
  \end{center}

  Nous avons $5$ points et $4$ sous-carr\'es. Par le principe des tiroirs, au moins un sous-carr\'e contient au moins $2$ points. La diagonale d'un carr\'e de c\^ot\'e $\frac{1}{2}$ est $\frac{\sqrt{2}}{2}$, ce qui est la distance maximale entre deux points d'un tel sous-carr\'e. Donc deux points du m\^eme sous-carr\'e sont \`a distance au plus $\frac{\sqrt{2}}{2}$. \qed
\end{exemple}

% ============================================================
\section{Principe du bon ordre}
\label{sec:bon-ordre}
\index{principe du bon ordre}
\index{bon ordre}

\begin{theoreme}[Principe du bon ordre de $\N$]\label{thm:bon-ordre}
  Toute partie non vide de $\N$ poss\`ede un plus petit \'el\'ement.
\end{theoreme}

\begin{remarque}[Lien avec la r\'ecurrence]\label{rem:bon-ordre-recurrence}
  Le principe du bon ordre et le principe de r\'ecurrence sont logiquement \'equivalents~: on peut d\'emontrer chacun \`a partir de l'autre. Le principe du bon ordre fournit une m\'ethode de d\'emonstration alternative~: pour montrer qu'une propri\'et\'e $P(n)$ est vraie pour tout $n \in \N$, on peut supposer que l'ensemble $S = \set{n \in \N : \neg P(n)}$ est non vide et consid\'erer son plus petit \'el\'ement, pour obtenir une contradiction.
\end{remarque}

\begin{theoreme}[Division euclidienne]\label{thm:div-euclidienne}
  Pour tous entiers $a \in \Z$ et $b \in \N^*$, il existe un unique couple $(q,r) \in \Z \times \N$ tel que~:
  \[
    a = bq + r \quad \text{avec} \quad 0 \leq r < b.
  \]
\end{theoreme}

\begin{proof}
  \textbf{Existence.} Consid\'erons l'ensemble~:
  \[
    S = \set{a - bk : k \in \Z} \cap \N = \set{a - bk : k \in \Z, \; a - bk \geq 0}.
  \]
  $S$ est non vide~: si $a \geq 0$, on prend $k = 0$ et $a - b \cdot 0 = a \geq 0$~; si $a < 0$, on prend $k = a$ (par exemple) et $a - ba = a(1 - b) \geq 0$ car $a < 0$ et $1 - b \leq 0$. Par le principe du bon ordre\index{principe du bon ordre}, $S$ poss\`ede un plus petit \'el\'ement $r = a - bq$ pour un certain $q \in \Z$.

  On a $r \geq 0$ par construction. Montrons $r < b$. Si $r \geq b$, alors $r - b = a - b(q+1) \geq 0$, donc $r - b \in S$ et $r - b < r$, ce qui contredit la minimalit\'e de $r$. Donc $0 \leq r < b$.

  \textbf{Unicit\'e.} Supposons $a = bq_1 + r_1 = bq_2 + r_2$ avec $0 \leq r_1, r_2 < b$. Alors $b(q_1 - q_2) = r_2 - r_1$. Ainsi $b \mid (r_2 - r_1)$. Or $\abs{r_2 - r_1} < b$ (car $0 \leq r_1, r_2 < b$), donc $r_2 - r_1 = 0$, d'o\`u $r_1 = r_2$ et $q_1 = q_2$.
\end{proof}

\begin{exemple}[Application du bon ordre]\label{ex:bon-ordre}
  Montrons que $\sqrt{3}$ est irrationnel en utilisant le principe du bon ordre.

  \smallskip
  \noindent\textit{D\'emonstration.} Supposons $\sqrt{3} \in \Q$. Alors l'ensemble~:
  \[
    S = \set{n \in \N^* : n\sqrt{3} \in \N^*}
  \]
  est non vide (il contient le d\'enominateur d'une \'ecriture fractionnaire de $\sqrt{3}$). Par le principe du bon ordre, $S$ poss\`ede un plus petit \'el\'ement $m$. Posons $m' = m\sqrt{3} - m = m(\sqrt{3} - 1)$. Puisque $1 < \sqrt{3} < 2$, on a $0 < \sqrt{3} - 1 < 1$, donc $0 < m' < m$. De plus~:
  \[
    m' \sqrt{3} = m(\sqrt{3} - 1)\sqrt{3} = m(3 - \sqrt{3}) = 3m - m\sqrt{3}.
  \]
  Or $m\sqrt{3} \in \N^*$ (puisque $m \in S$), donc $m'\sqrt{3} = 3m - m\sqrt{3} \in \Z$. Comme $m' > 0$ et $\sqrt{3} > 0$, on a $m'\sqrt{3} > 0$, donc $m'\sqrt{3} \in \N^*$. Ainsi $m' \in S$ avec $m' < m$, ce qui contredit la minimalit\'e de $m$. \qed
\end{exemple}

% ============================================================
\section{Exercices}
\label{sec:exo-demonstration}

\subsection*{D\'emonstration directe}

\begin{exercice}[$\bigstar$]\label{exo:direct-1}
  D\'emontrer directement que pour tout entier $n$, si $n$ est pair, alors $n^2$ est pair.
\end{exercice}

\begin{exercice}[$\bigstar$]\label{exo:direct-2}
  Montrer que la somme de trois entiers cons\'ecutifs est toujours divisible par $3$.
\end{exercice}

\begin{exercice}[$\bigstar\bigstar$]\label{exo:direct-3}
  Montrer que pour tous r\'eels $a, b \geq 0$, on a $\dfrac{a+b}{2} \geq \sqrt{ab}$ (in\'egalit\'e arithm\'etico-g\'eom\'etrique\index{inegalite arithmetico-geometrique@in\'egalit\'e arithm\'etico-g\'eom\'etrique}).
\end{exercice}

\subsection*{D\'emonstration par l'absurde}

\begin{exercice}[$\bigstar$]\label{exo:absurde-1}
  Montrer par l'absurde que $\sqrt{3}$ est irrationnel.
\end{exercice}

\begin{exercice}[$\bigstar\bigstar$]\label{exo:absurde-2}
  Montrer qu'il n'existe pas d'entiers $a$ et $b$ tels que $a^2 = 4b + 2$.
\end{exercice}

\begin{exercice}[$\bigstar\bigstar$]\label{exo:absurde-3}
  Montrer par l'absurde que si $a$ et $b$ sont des r\'eels tels que $a + b$ est irrationnel, alors $a$ est irrationnel ou $b$ est irrationnel.
\end{exercice}

\subsection*{D\'emonstration par contrapos\'ee}

\begin{exercice}[$\bigstar$]\label{exo:contrap-1}
  Montrer par contrapos\'ee~: pour tout entier $n$, si $n^2$ est impair alors $n$ est impair.
\end{exercice}

\begin{exercice}[$\bigstar\bigstar$]\label{exo:contrap-2}
  Montrer par contrapos\'ee~: pour tous r\'eels $x$ et $y$, si $xy$ est irrationnel, alors $x$ est irrationnel ou $y$ est irrationnel.
\end{exercice}

\subsection*{R\'ecurrence}

\begin{exercice}[$\bigstar$]\label{exo:rec-1}
  Montrer par r\'ecurrence que pour tout $n \geq 1$~:
  \[
    \sum_{k=1}^{n} k^2 = \frac{n(n+1)(2n+1)}{6}.
  \]
\end{exercice}

\begin{exercice}[$\bigstar$]\label{exo:rec-2}
  Montrer par r\'ecurrence que pour tout $n \geq 1$~:
  \[
    \sum_{k=1}^{n} (2k - 1) = n^2.
  \]
\end{exercice}

\begin{exercice}[$\bigstar\bigstar$]\label{exo:rec-3}
  Montrer par r\'ecurrence que $n! \geq 2^{n-1}$ pour tout $n \geq 1$.
\end{exercice}

\begin{exercice}[$\bigstar\bigstar$]\label{exo:rec-4}
  Montrer par r\'ecurrence forte que tout entier $n \geq 2$ est divisible par un nombre premier.
\end{exercice}

\begin{exercice}[$\bigstar\bigstar\bigstar$]\label{exo:rec-5}
  Soit $(F_n)$ la suite de Fibonacci. Montrer par r\'ecurrence forte que~:
  \[
    F_n = \frac{\varphi^n - \psi^n}{\sqrt{5}}, \quad \text{o\`u}\quad \varphi = \frac{1+\sqrt{5}}{2},\quad \psi = \frac{1-\sqrt{5}}{2}.
  \]
  \textit{Indication~: v\'erifier que $\varphi^2 = \varphi + 1$ et $\psi^2 = \psi + 1$.}
\end{exercice}

\subsection*{Principe des tiroirs}

\begin{exercice}[$\bigstar$]\label{exo:tiroirs-1}
  Montrer que dans un groupe de $13$ personnes, au moins deux sont n\'ees le m\^eme mois.
\end{exercice}

\begin{exercice}[$\bigstar\bigstar$]\label{exo:tiroirs-2}
  Soit $A$ un sous-ensemble de $\set{1, 2, \dots, 2n}$ contenant $n+1$ \'el\'ements. Montrer que $A$ contient deux \'el\'ements dont l'un divise l'autre.
\end{exercice}

\begin{exercice}[$\bigstar\bigstar\bigstar$]\label{exo:tiroirs-3}
  On choisit $n+1$ entiers dans l'ensemble $\set{1, 2, \dots, 2n}$. Montrer qu'il en existe deux dont la diff\'erence vaut exactement $1$.
  \textit{Indication~: consid\'erer les <<~tiroirs~>> $\set{1,2}, \set{3,4}, \dots, \set{2n-1, 2n}$.}
\end{exercice}

\begin{exercice}[$\bigstar\bigstar\bigstar$]\label{exo:tiroirs-4}
  Soit $(a_1, a_2, \dots, a_{n^2+1})$ une suite de $n^2+1$ r\'eels distincts. Montrer qu'elle contient soit une sous-suite croissante de longueur $n+1$, soit une sous-suite d\'ecroissante de longueur $n+1$.
  \textit{Indication~: associer \`a chaque $a_i$ la longueur de la plus longue sous-suite croissante commen\c{c}ant en $a_i$.}
\end{exercice}

% ============================================================
\section*{R\'esum\'e du chapitre}
\addcontentsline{toc}{section}{R\'esum\'e du chapitre}

\begin{itemize}[leftmargin=2em]
  \item La \textbf{d\'emonstration directe} de $p \to q$ suppose $p$ vrai et en d\'eduit $q$ par une cha\^ine logique.
  \item La \textbf{d\'emonstration par l'absurde} suppose $\neg\varphi$ et aboutit \`a une contradiction, prouvant ainsi $\varphi$.
  \item La \textbf{d\'emonstration par contrapos\'ee} de $p \to q$ consiste \`a d\'emontrer $\neg q \to \neg p$.
  \item La \textbf{r\'ecurrence simple}~: on montre $P(n_0)$ (base) et $P(n) \Rightarrow P(n+1)$ (h\'er\'edit\'e).
  \item La \textbf{r\'ecurrence forte}~: l'hypoth\`ese porte sur $P(n_0), \dots, P(n)$ pour d\'emontrer $P(n+1)$.
  \item La \textbf{r\'ecurrence structurelle} g\'en\'eralise la r\'ecurrence aux structures d\'efinies inductivement.
  \item Le \textbf{principe des tiroirs} (Dirichlet)~: $n+1$ objets dans $n$ tiroirs impliquent un tiroir avec au moins deux objets. Forme g\'en\'eralis\'ee~: au moins $\lceil N/k \rceil$ objets dans un tiroir.
  \item Le \textbf{principe du bon ordre}~: toute partie non vide de $\N$ a un plus petit \'el\'ement. Ce principe est \'equivalent au principe de r\'ecurrence.
\end{itemize}

% === FIN DU CHUNK preamble+ch1-2 ===
% === DÉBUT DU CHUNK ch3-4 ===

%%%%%%%%%%%%%%%%%%%%%%%%%%%%%%%%%%%%%%%%%%%%%%%%%%%%%%%%%%%%%
%  CHAPITRE 3 — Ensembles, Relations, Fonctions
%%%%%%%%%%%%%%%%%%%%%%%%%%%%%%%%%%%%%%%%%%%%%%%%%%%%%%%%%%%%%
\chapter{Ensembles, Relations, Fonctions}\label{ch:ensembles-relations-fonctions}
\index{ensemble}\index{relation}\index{fonction}

\begin{center}
\textit{«~La théorie des ensembles est le langage universel des mathématiques.~»}

--- Georg Cantor
\end{center}

\bigskip

Ce chapitre présente les fondements ensemblistes sur lesquels repose
l'ensemble des mathématiques discrètes. Nous y étudions les ensembles et
leurs opérations, les relations et leurs propriétés, puis les fonctions et la
notion de cardinalité.

%=====================================================
\section{Ensembles : notations et opérations fondamentales}\label{sec:ensembles-notations}
%=====================================================

\subsection{Notions de base}

\begin{definition}[Ensemble]\label{def:ensemble}
\index{ensemble!définition}
Un \textbf{ensemble} est une collection d'objets distincts, appelés
\textbf{éléments}. On note $x \in A$ pour signifier que $x$ est un
élément de l'ensemble $A$, et $x \notin A$ dans le cas contraire.
\end{definition}

\begin{notation}\label{not:ensembles-classiques}
Les ensembles classiques sont notés :
\[
  \N \subset \Z \subset \Q \subset \R \subset \C.
\]
L'ensemble vide est noté $\varnothing$ ou $\{\}$.
\end{notation}

On décrit un ensemble soit en \emph{extension}, soit en \emph{compréhension} :
\begin{itemize}
  \item \textbf{Extension :} $A = \{2, 3, 5, 7, 11\}$.
  \item \textbf{Compréhension :} $A = \set{p \in \N \mid p \text{ est premier et } p \leq 11}$.
\end{itemize}

\begin{definition}[Inclusion, égalité]\label{def:inclusion}
\index{inclusion}\index{sous-ensemble}
Soient $A$ et $B$ deux ensembles.
\begin{enumerate}
  \item $A \subseteq B$ (\textbf{$A$ est inclus dans $B$}) si tout élément de $A$ est élément de $B$ :
    \[
      A \subseteq B \iff \bigl(\forall x,\; x \in A \Rightarrow x \in B\bigr).
    \]
  \item $A = B$ si et seulement si $A \subseteq B$ et $B \subseteq A$.
  \item $A \subsetneq B$ (\textbf{inclusion stricte}) si $A \subseteq B$ et $A \neq B$.
\end{enumerate}
\end{definition}

\begin{definition}[Ensemble des parties]\label{def:powerset}
\index{ensemble des parties}\index{powerset@$\powerset$}
L'\textbf{ensemble des parties} de $A$, noté $\powerset(A)$, est
l'ensemble de tous les sous-ensembles de $A$ :
\[
  \powerset(A) = \set{B \mid B \subseteq A}.
\]
\end{definition}

\begin{exemple}[Parties d'un petit ensemble]\label{ex:parties}
Si $A = \{1, 2, 3\}$, alors
\[
  \powerset(A) = \bigl\{\varnothing, \{1\}, \{2\}, \{3\}, \{1,2\}, \{1,3\}, \{2,3\}, \{1,2,3\}\bigr\}.
\]
On observe que $\card{\powerset(A)} = 2^{\card{A}} = 2^3 = 8$.
\end{exemple}

\begin{definition}[Produit cartésien]\label{def:produit-cartesien}
\index{produit cartésien}
Le \textbf{produit cartésien} de $A$ et $B$ est
\[
  A \times B = \set{(a, b) \mid a \in A,\; b \in B}.
\]
Plus généralement, $A_1 \times A_2 \times \cdots \times A_n
= \set{(a_1, a_2, \ldots, a_n) \mid a_i \in A_i \text{ pour tout } i}$.
\end{definition}

%-----------------------------------------------------
\subsection{Opérations ensemblistes}
%-----------------------------------------------------

Soit $\Omega$ un ensemble de référence (univers).

\begin{definition}[Opérations sur les ensembles]\label{def:operations-ensembles}
\index{union}\index{intersection}\index{complément}\index{différence!ensembliste}\index{différence symétrique}
Soient $A, B \subseteq \Omega$.
\begin{enumerate}
  \item \textbf{Union :} $A \cup B = \set{x \in \Omega \mid x \in A \text{ ou } x \in B}$.
  \item \textbf{Intersection :} $A \cap B = \set{x \in \Omega \mid x \in A \text{ et } x \in B}$.
  \item \textbf{Complément :} $\overline{A} = A^c = \set{x \in \Omega \mid x \notin A}$.
  \item \textbf{Différence :} $A \setminus B = \set{x \in A \mid x \notin B}$.
  \item \textbf{Différence symétrique :} $A \bigtriangleup B = (A \setminus B) \cup (B \setminus A)$.
\end{enumerate}
\end{definition}

% --- Diagrammes de Venn ---
\begin{figure}[ht]
\centering
\begin{tikzpicture}[scale=0.7]
  % Union
  \begin{scope}[shift={(-5,0)}]
    \fill[blue!20] (-0.6,0) circle (1.2);
    \fill[blue!20] (0.6,0) circle (1.2);
    \draw (-0.6,0) circle (1.2) node[left=0.8cm] {$A$};
    \draw (0.6,0) circle (1.2) node[right=0.8cm] {$B$};
    \node at (0,-2) {$A \cup B$};
  \end{scope}
  % Intersection
  \begin{scope}
    \begin{scope}
      \clip (-0.6,0) circle (1.2);
      \fill[red!25] (0.6,0) circle (1.2);
    \end{scope}
    \draw (-0.6,0) circle (1.2) node[left=0.8cm] {$A$};
    \draw (0.6,0) circle (1.2) node[right=0.8cm] {$B$};
    \node at (0,-2) {$A \cap B$};
  \end{scope}
  % Difference
  \begin{scope}[shift={(5,0)}]
    \fill[green!25] (-0.6,0) circle (1.2);
    \fill[white] (0.6,0) circle (1.2);
    \draw (-0.6,0) circle (1.2) node[left=0.8cm] {$A$};
    \draw (0.6,0) circle (1.2) node[right=0.8cm] {$B$};
    \node at (0,-2) {$A \setminus B$};
  \end{scope}
\end{tikzpicture}
\caption{Diagrammes de Venn : union, intersection, différence.}
\label{fig:venn-operations}
\end{figure}

\begin{theoreme}[Lois de De Morgan pour les ensembles]\label{thm:de-morgan-ensembles}
\index{De Morgan!lois (ensembles)}
Soient $A, B \subseteq \Omega$. Alors :
\begin{enumerate}
  \item $\overline{A \cup B} = \overline{A} \cap \overline{B}$,
  \item $\overline{A \cap B} = \overline{A} \cup \overline{B}$.
\end{enumerate}
\end{theoreme}

\begin{proof}
Montrons (1) par double inclusion.

$(\subseteq)$ Soit $x \in \overline{A \cup B}$. Alors $x \notin A \cup B$,
donc $x \notin A$ \emph{et} $x \notin B$. Par conséquent
$x \in \overline{A}$ et $x \in \overline{B}$, d'où
$x \in \overline{A} \cap \overline{B}$.

$(\supseteq)$ Soit $x \in \overline{A} \cap \overline{B}$. Alors
$x \notin A$ et $x \notin B$, donc $x \notin A \cup B$, c'est-à-dire
$x \in \overline{A \cup B}$.

La preuve de (2) est analogue (ou s'obtient en appliquant (1) à
$\overline{A}$ et $\overline{B}$, puis en simplifiant).
\end{proof}

\begin{remarque}[Paradoxe de Russell]\label{rem:russell}
\index{Russell!paradoxe de}
La notion naïve d'ensemble conduit à des contradictions. Considérons
l'«~ensemble~» $R = \set{X \mid X \notin X}$. Si $R \in R$, alors par
définition $R \notin R$ ; si $R \notin R$, alors $R \in R$. Cette contradiction,
découverte par Bertrand Russell en 1901, a motivé l'axiomatisation de la
théorie des ensembles (par exemple les axiomes de Zermelo--Fraenkel).
\end{remarque}

%=====================================================
\section{Relations}\label{sec:relations}
%=====================================================

\begin{definition}[Relation binaire]\label{def:relation-binaire}
\index{relation!binaire}
Une \textbf{relation binaire} de $A$ vers $B$ est un sous-ensemble
$\mathcal{R} \subseteq A \times B$. Si $A = B$, on dit que $\mathcal{R}$
est une relation \emph{sur} $A$. On écrit $a\,\mathcal{R}\,b$ au lieu de
$(a,b) \in \mathcal{R}$.
\end{definition}

\begin{definition}[Propriétés d'une relation sur $A$]\label{def:proprietes-relation}
\index{relation!réflexive}\index{relation!symétrique}\index{relation!antisymétrique}\index{relation!transitive}
Soit $\mathcal{R}$ une relation sur $A$.
\begin{enumerate}
  \item \textbf{Réflexive :} $\forall a \in A,\; a\,\mathcal{R}\,a$.
  \item \textbf{Symétrique :} $\forall a,b \in A,\; a\,\mathcal{R}\,b \Rightarrow b\,\mathcal{R}\,a$.
  \item \textbf{Antisymétrique :} $\forall a,b \in A,\;
    \bigl(a\,\mathcal{R}\,b \text{ et } b\,\mathcal{R}\,a\bigr) \Rightarrow a = b$.
  \item \textbf{Transitive :} $\forall a,b,c \in A,\;
    \bigl(a\,\mathcal{R}\,b \text{ et } b\,\mathcal{R}\,c\bigr) \Rightarrow a\,\mathcal{R}\,c$.
\end{enumerate}
\end{definition}

%-----------------------------------------------------
\subsection{Relations d'équivalence}\label{subsec:equiv}
%-----------------------------------------------------

\begin{definition}[Relation d'équivalence]\label{def:equivalence}
\index{relation!d'équivalence}\index{équivalence}
Une relation $\sim$ sur $A$ est une \textbf{relation d'équivalence} si elle
est réflexive, symétrique et transitive.
\end{definition}

\begin{definition}[Classe d'équivalence]\label{def:classe-equiv}
\index{classe d'équivalence}
Si $\sim$ est une relation d'équivalence sur $A$ et $a \in A$, la
\textbf{classe d'équivalence} de $a$ est
\[
  [a] = \set{x \in A \mid x \sim a}.
\]
L'ensemble des classes d'équivalence, noté $A/{\sim}$, est appelé
\textbf{ensemble quotient}\index{ensemble quotient}.
\end{definition}

\begin{exemple}[Congruence modulo $n$]\label{ex:congruence-equiv}
Sur $\Z$, la relation $a \equiv b \pmod{n}$ est une relation d'équivalence.
Les classes sont $[0], [1], \ldots, [n-1]$ et l'ensemble quotient est
$\Z/n\Z$.
\end{exemple}

\begin{definition}[Partition]\label{def:partition}
\index{partition}
Une \textbf{partition} d'un ensemble $A$ est une famille
$\{A_i\}_{i \in I}$ de sous-ensembles non vides de $A$ tels que :
\begin{enumerate}
  \item $\displaystyle\bigcup_{i \in I} A_i = A$,
  \item $A_i \cap A_j = \varnothing$ pour tous $i \neq j$.
\end{enumerate}
\end{definition}

\begin{theoreme}[Correspondance équivalence--partition]\label{thm:equiv-partition}
\index{partition!et équivalence}
Soit $A$ un ensemble non vide. Il y a une correspondance bijective entre
les relations d'équivalence sur $A$ et les partitions de $A$ : les classes
d'équivalence d'une relation d'équivalence forment une partition, et
réciproquement toute partition induit une relation d'équivalence.
\end{theoreme}

\begin{proof}
\textbf{(Équivalence $\Rightarrow$ Partition).}
Soit $\sim$ une relation d'équivalence sur $A$. Montrons que
$\{[a] \mid a \in A\}$ est une partition.
\begin{itemize}
  \item \emph{Non-vacuité.} Pour tout $a \in A$, on a $a \in [a]$ (réflexivité), donc $[a] \neq \varnothing$.
  \item \emph{Recouvrement.} Comme $a \in [a]$ pour tout $a$, on a
    $A = \bigcup_{a \in A} [a]$.
  \item \emph{Disjonction.} Supposons $[a] \cap [b] \neq \varnothing$ et
    montrons $[a] = [b]$. Soit $c \in [a] \cap [b]$, donc $c \sim a$ et
    $c \sim b$. Soit $x \in [a]$; alors $x \sim a$. Par symétrie, $a \sim c$;
    par transitivité, $x \sim c$; puis $c \sim b$ donne $x \sim b$, donc
    $x \in [b]$. Ainsi $[a] \subseteq [b]$. Par symétrie, $[b] \subseteq [a]$.
\end{itemize}

\textbf{(Partition $\Rightarrow$ Équivalence).}
Soit $\{A_i\}_{i \in I}$ une partition de $A$. Définissons $a \sim b$ si et
seulement s'il existe $i \in I$ tel que $a, b \in A_i$. On vérifie aisément
que $\sim$ est réflexive (car $a$ appartient à un certain $A_i$), symétrique
(par symétrie de la définition) et transitive (si $a, b \in A_i$ et
$b, c \in A_j$, alors $b \in A_i \cap A_j$, d'où $i = j$ par disjonction,
et donc $a, c \in A_i$).
\end{proof}

%-----------------------------------------------------
\subsection{Ordres partiels et diagrammes de Hasse}\label{subsec:ordres}
%-----------------------------------------------------

\begin{definition}[Ordre partiel]\label{def:ordre-partiel}
\index{ordre partiel}
Une relation $\preceq$ sur $A$ est un \textbf{ordre partiel} si elle est
réflexive, antisymétrique et transitive. Le couple $(A, \preceq)$ est un
\textbf{ensemble partiellement ordonné} (ou \textbf{poset}).
\end{definition}

\begin{definition}[Chaîne, antichaîne]\label{def:chaine-antichaine}
\index{chaîne}\index{antichaîne}
Dans un poset $(A, \preceq)$ :
\begin{itemize}
  \item une \textbf{chaîne} est un sous-ensemble $C \subseteq A$ tel que
    deux éléments quelconques de $C$ sont comparables ;
  \item une \textbf{antichaîne} est un sous-ensemble $S \subseteq A$ tel que
    deux éléments distincts quelconques de $S$ sont incomparables.
\end{itemize}
\end{definition}

\begin{exemple}[Divisibilité]\label{ex:divisibilite-poset}
Sur $A = \{1, 2, 3, 4, 6, 12\}$, la relation $a \mid b$ est un ordre
partiel. Le diagramme de Hasse correspondant est donné à la
figure~\ref{fig:hasse-div12}.
\end{exemple}

\begin{figure}[ht]
\centering
\begin{tikzpicture}[scale=1.1,
    every node/.style={draw, circle, minimum size=7mm, inner sep=0pt, font=\small}]
  \node (1)  at (0,0)   {$1$};
  \node (2)  at (-1.5,1) {$2$};
  \node (3)  at (1.5,1) {$3$};
  \node (4)  at (-1.5,2) {$4$};
  \node (6)  at (1.5,2) {$6$};
  \node (12) at (0,3)   {$12$};
  \draw (1)--(2) (1)--(3) (2)--(4) (2)--(6) (3)--(6) (4)--(12) (6)--(12);
\end{tikzpicture}
\caption{Diagramme de Hasse de $(\{1,2,3,4,6,12\}, \mid)$.}
\label{fig:hasse-div12}
\end{figure}

\begin{remarque}\label{rem:hasse-convention}
Dans un diagramme de Hasse, on ne trace que les relations de
\emph{couverture} : $a$ couvre $b$ si $b \prec a$ et il n'existe aucun $c$
avec $b \prec c \prec a$. Les relations déduites par réflexivité et
transitivité sont omises.
\end{remarque}

\begin{definition}[Éléments remarquables d'un poset]\label{def:elements-poset}
\index{minimum}\index{maximum}\index{borne supérieure}\index{borne inférieure}
Soit $(A, \preceq)$ un poset et $S \subseteq A$.
\begin{itemize}
  \item Un \textbf{élément minimal} de $S$ est un $a \in S$ tel qu'il
    n'existe aucun $b \in S$ avec $b \prec a$.
  \item Un \textbf{minimum} de $S$ est un $a \in S$ tel que $a \preceq b$
    pour tout $b \in S$ (s'il existe, il est unique).
  \item La \textbf{borne supérieure} (ou \emph{supremum}) de $S$, notée
    $\sup S$, est le plus petit majorant de $S$.
  \item La \textbf{borne inférieure} (ou \emph{infimum}) de $S$, notée
    $\inf S$, est le plus grand minorant de $S$.
\end{itemize}
\end{definition}

\begin{exemple}[Poset $(\powerset(\{a,b\}), \subseteq)$]\label{ex:hasse-powerset}
\begin{center}
\begin{tikzpicture}[scale=1.0,
    every node/.style={draw, rounded corners, minimum size=6mm, inner sep=2pt, font=\small}]
  \node (vide) at (0,0) {$\varnothing$};
  \node (a) at (-1.2,1.5) {$\{a\}$};
  \node (b) at (1.2,1.5) {$\{b\}$};
  \node (ab) at (0,3) {$\{a,b\}$};
  \draw (vide)--(a) (vide)--(b) (a)--(ab) (b)--(ab);
\end{tikzpicture}
\captionof{figure}{Diagramme de Hasse de $(\powerset(\{a,b\}), \subseteq)$.}
\label{fig:hasse-powerset-ab}
\end{center}
Ce poset possède un minimum ($\varnothing$) et un maximum ($\{a,b\}$).
Les éléments $\{a\}$ et $\{b\}$ forment une antichaîne.
\end{exemple}

\begin{theoreme}[Dilworth]\label{thm:dilworth}
\index{Dilworth!théorème de}
Dans un poset fini, la taille maximale d'une antichaîne est égale au nombre
minimal de chaînes nécessaires pour partitionner le poset.
\end{theoreme}

%=====================================================
\section{Fonctions}\label{sec:fonctions}
%=====================================================

\begin{definition}[Fonction]\label{def:fonction}
\index{fonction!définition}
Une \textbf{fonction} $f \colon A \to B$ est une relation
$f \subseteq A \times B$ telle que pour tout $a \in A$, il existe un unique
$b \in B$ avec $(a, b) \in f$. L'ensemble $A$ est le \textbf{domaine},
$B$ le \textbf{codomaine}.
\end{definition}

\begin{definition}[Injection, surjection, bijection]\label{def:inj-surj-bij}
\index{injection}\index{surjection}\index{bijection}
Soit $f \colon A \to B$.
\begin{enumerate}
  \item $f$ est \textbf{injective} si
    $\forall a_1, a_2 \in A,\; f(a_1) = f(a_2) \Rightarrow a_1 = a_2$.
  \item $f$ est \textbf{surjective} si
    $\forall b \in B,\; \exists\, a \in A,\; f(a) = b$.
  \item $f$ est \textbf{bijective} si elle est à la fois injective et surjective.
\end{enumerate}
\end{definition}

% --- Diagramme fonctions ---
\begin{figure}[ht]
\centering
\begin{tikzpicture}[scale=0.75,
    dot/.style={circle, fill, inner sep=1.5pt},
    >=stealth]
  % Injection
  \begin{scope}[shift={(-5,0)}]
    \draw[rounded corners] (-0.7,-1.8) rectangle (0.7,1.8);
    \draw[rounded corners] (2.3,-1.8) rectangle (3.7,1.8);
    \node[dot,label=left:{$a$}] (a1) at (0,1) {};
    \node[dot,label=left:{$b$}] (a2) at (0,0) {};
    \node[dot,label=left:{$c$}] (a3) at (0,-1) {};
    \node[dot,label=right:{$1$}] (b1) at (3,1.2) {};
    \node[dot,label=right:{$2$}] (b2) at (3,0.4) {};
    \node[dot,label=right:{$3$}] (b3) at (3,-0.4) {};
    \node[dot,label=right:{$4$}] (b4) at (3,-1.2) {};
    \draw[->] (a1)--(b1); \draw[->] (a2)--(b3); \draw[->] (a3)--(b4);
    \node at (1.5,-2.5) {Injection};
  \end{scope}
  % Surjection
  \begin{scope}
    \draw[rounded corners] (-0.7,-1.8) rectangle (0.7,1.8);
    \draw[rounded corners] (2.3,-1.8) rectangle (3.7,1.8);
    \node[dot,label=left:{$a$}] (c1) at (0,1.2) {};
    \node[dot,label=left:{$b$}] (c2) at (0,0.4) {};
    \node[dot,label=left:{$c$}] (c3) at (0,-0.4) {};
    \node[dot,label=left:{$d$}] (c4) at (0,-1.2) {};
    \node[dot,label=right:{$1$}] (d1) at (3,1) {};
    \node[dot,label=right:{$2$}] (d2) at (3,0) {};
    \node[dot,label=right:{$3$}] (d3) at (3,-1) {};
    \draw[->] (c1)--(d1); \draw[->] (c2)--(d1); \draw[->] (c3)--(d2); \draw[->] (c4)--(d3);
    \node at (1.5,-2.5) {Surjection};
  \end{scope}
  % Bijection
  \begin{scope}[shift={(5,0)}]
    \draw[rounded corners] (-0.7,-1.8) rectangle (0.7,1.8);
    \draw[rounded corners] (2.3,-1.8) rectangle (3.7,1.8);
    \node[dot,label=left:{$a$}] (e1) at (0,1) {};
    \node[dot,label=left:{$b$}] (e2) at (0,0) {};
    \node[dot,label=left:{$c$}] (e3) at (0,-1) {};
    \node[dot,label=right:{$1$}] (f1) at (3,1) {};
    \node[dot,label=right:{$2$}] (f2) at (3,0) {};
    \node[dot,label=right:{$3$}] (f3) at (3,-1) {};
    \draw[->] (e1)--(f2); \draw[->] (e2)--(f3); \draw[->] (e3)--(f1);
    \node at (1.5,-2.5) {Bijection};
  \end{scope}
\end{tikzpicture}
\caption{Illustration de l'injection, la surjection et la bijection.}
\label{fig:inj-surj-bij}
\end{figure}

\begin{definition}[Composition]\label{def:composition}
\index{composition de fonctions}
Si $f \colon A \to B$ et $g \colon B \to C$, la \textbf{composée}
$g \circ f \colon A \to C$ est définie par $(g \circ f)(a) = g\bigl(f(a)\bigr)$.
\end{definition}

\begin{proposition}[Composition et injectivité/surjectivité]\label{prop:comp-inj-surj}
\index{composition!et injectivité}
\begin{enumerate}
  \item Si $f$ et $g$ sont injectives, alors $g \circ f$ est injective.
  \item Si $f$ et $g$ sont surjectives, alors $g \circ f$ est surjective.
  \item Si $g \circ f$ est injective, alors $f$ est injective.
  \item Si $g \circ f$ est surjective, alors $g$ est surjective.
\end{enumerate}
\end{proposition}

\begin{proof}
(1) Supposons $g(f(a_1)) = g(f(a_2))$. L'injectivité de $g$ donne
$f(a_1) = f(a_2)$, puis celle de $f$ donne $a_1 = a_2$.

(3) Supposons $f(a_1) = f(a_2)$. Alors $g(f(a_1)) = g(f(a_2))$, et
l'injectivité de $g \circ f$ donne $a_1 = a_2$. Les preuves de (2) et (4)
sont laissées en exercice.
\end{proof}

\begin{definition}[Fonction inverse]\label{def:fonction-inverse}
\index{fonction!inverse}
Si $f \colon A \to B$ est bijective, il existe une unique fonction
$f^{-1} \colon B \to A$ telle que $f^{-1} \circ f = \mathrm{id}_A$ et
$f \circ f^{-1} = \mathrm{id}_B$.
\end{definition}

%=====================================================
\section{Cardinalité}\label{sec:cardinalite}
%=====================================================
\index{cardinalité}

\begin{definition}[Équipotence]\label{def:equipotence}
\index{équipotence}
Deux ensembles $A$ et $B$ sont \textbf{équipotents} (ont même cardinalité),
noté $\card{A} = \card{B}$, s'il existe une bijection $f \colon A \to B$.
\end{definition}

\begin{definition}[Ensemble dénombrable]\label{def:denombrable}
\index{ensemble!dénombrable}\index{dénombrable}
Un ensemble $A$ est \textbf{dénombrable} s'il est fini ou s'il existe une
bijection $f \colon \N \to A$. Un ensemble qui n'est pas dénombrable est dit
\textbf{indénombrable}\index{indénombrable}.
\end{definition}

\begin{exemple}\label{ex:Z-denombrable}
L'ensemble $\Z$ est dénombrable. La bijection $f \colon \N \to \Z$ peut être
définie par :
\[
  f(n) = \begin{cases} n/2 & \text{si } n \text{ est pair,} \\ -(n+1)/2 & \text{si } n \text{ est impair.} \end{cases}
\]
Ainsi $f(0)=0,\; f(1)=-1,\; f(2)=1,\; f(3)=-2,\; f(4)=2,\; \ldots$
\end{exemple}

\begin{theoreme}[Cantor : indénombrabilité de $\R$]\label{thm:cantor-diagonal}
\index{Cantor!argument diagonal}\index{argument diagonal}
L'ensemble $\R$ est indénombrable. Plus précisément, l'intervalle
$]0,1[$ n'est pas dénombrable.
\end{theoreme}

\begin{proof}[Preuve par l'argument diagonal de Cantor]
Supposons par l'absurde qu'il existe une bijection $f \colon \N \to \,]0,1[\,$
et écrivons les développements décimaux :
\[
  f(0) = 0{,}d_{00}\,d_{01}\,d_{02}\,\ldots, \quad
  f(1) = 0{,}d_{10}\,d_{11}\,d_{12}\,\ldots, \quad
  f(2) = 0{,}d_{20}\,d_{21}\,d_{22}\,\ldots, \quad \ldots
\]
Construisons $x = 0{,}x_0 x_1 x_2 \ldots$ avec
\[
  x_n = \begin{cases} 3 & \text{si } d_{nn} \neq 3, \\ 7 & \text{si } d_{nn} = 3. \end{cases}
\]
On a $x \in \,]0,1[\,$ (on évite les chiffres $0$ et $9$ pour exclure les
ambiguïtés de représentation décimale). Or $x \neq f(n)$ pour tout
$n \in \N$ car $x$ diffère de $f(n)$ au rang $n$. Ceci contredit la
surjectivité de $f$.
\end{proof}

\begin{corollaire}\label{cor:powerset-N}
$\card{\N} < \card{\powerset(\N)} = \card{\R}$.
\end{corollaire}

\begin{theoreme}[Cantor : pas de surjection vers l'ensemble des parties]\label{thm:cantor-general}
\index{Cantor!théorème général}
Pour tout ensemble $A$, il n'existe pas de surjection $f \colon A \to \powerset(A)$.
En particulier, $\card{A} < \card{\powerset(A)}$.
\end{theoreme}

\begin{proof}
Supposons par l'absurde qu'il existe une surjection $f \colon A \to \powerset(A)$.
Définissons $B = \set{a \in A \mid a \notin f(a)}$. Comme $B \subseteq A$, on
a $B \in \powerset(A)$, donc par surjectivité il existe $b \in A$ tel que
$f(b) = B$.

Si $b \in B$, alors par définition de $B$, on a $b \notin f(b) = B$ :
contradiction. Si $b \notin B$, alors $b \notin f(b)$, donc par définition de
$B$, on a $b \in B$ : contradiction. L'hypothèse est donc absurde.
\end{proof}

\begin{proposition}[Propriétés de cardinalité]\label{prop:card-props}
\index{cardinalité!propriétés}
\begin{enumerate}
  \item Si $A \subseteq B$ et $B$ est dénombrable, alors $A$ est dénombrable.
  \item L'union dénombrable d'ensembles dénombrables est dénombrable.
  \item Le produit cartésien de deux ensembles dénombrables est dénombrable.
  \item L'ensemble $\{0,1\}^{\N}$ (suites binaires infinies) est
    indénombrable et a la même cardinalité que $\R$.
\end{enumerate}
\end{proposition}

\begin{exemple}[Nombres algébriques vs transcendants]\label{ex:algebriques}
\index{nombre algébrique}\index{nombre transcendant}
L'ensemble des nombres algébriques (racines de polynômes à coefficients
entiers) est dénombrable, car il est réunion dénombrable d'ensembles finis.
Par conséquent, l'ensemble des nombres transcendants est indénombrable ---
un argument d'existence non constructif remarquable.
\end{exemple}

%=====================================================
\section{Exercices}\label{sec:exo-ch3}
%=====================================================

\begin{exercice}[Opérations ensemblistes]\label{exo:ops-ensembles}
Soient $A = \{1,2,3,4,5\}$, $B = \{3,4,5,6,7\}$ et
$\Omega = \{1,2,\ldots,10\}$. Calculer $A \cup B$, $A \cap B$,
$A \setminus B$, $\overline{A}$, $A \bigtriangleup B$ et vérifier les lois
de De Morgan.
\end{exercice}

\begin{exercice}[Ensemble des parties]\label{exo:powerset-cardinal}
Montrer par récurrence que $\card{\powerset(A)} = 2^{\card{A}}$ pour tout
ensemble fini $A$.
\end{exercice}

\begin{exercice}[Relation d'équivalence]\label{exo:equiv-relation}
Sur $\Z \times (\Z \setminus \{0\})$, on définit $(a,b) \sim (c,d)$ si et
seulement si $ad = bc$. Montrer que $\sim$ est une relation d'équivalence et
interpréter les classes.
\end{exercice}

\begin{exercice}[Composition de fonctions]\label{exo:comp-surj}
Montrer les points (2) et (4) de la proposition~\ref{prop:comp-inj-surj}.
\end{exercice}

\begin{exercice}[Dénombrabilité de $\Q$]\label{exo:Q-denombrable}
Montrer que $\Q$ est dénombrable en utilisant un argument de diagonale de
Cantor (parcours en zigzag de $\N \times \N^*$).
\end{exercice}

\begin{exercice}[Théorème de Cantor]\label{exo:cantor-powerset}
Montrer que pour tout ensemble $A$ (fini ou infini), il n'existe pas de
surjection $f \colon A \to \powerset(A)$.
\textit{Indication : considérer $B = \set{a \in A \mid a \notin f(a)}$.}
\end{exercice}

\begin{exercice}[Diagramme de Hasse]\label{exo:hasse}
Dessiner le diagramme de Hasse de $(\powerset(\{a,b,c\}), \subseteq)$.
Identifier les chaînes maximales et les antichaînes maximales.
\end{exercice}

\begin{exercice}[Schröder--Bernstein]\label{exo:schroder-bernstein}
\index{Schröder--Bernstein}
En admettant le théorème de Cantor--Bernstein (s'il existe une injection
de $A$ dans $B$ et une injection de $B$ dans $A$, alors il existe une
bijection entre $A$ et $B$), montrer que $\card{]0,1[} = \card{\R}$.
\end{exercice}

%=====================================================
\section*{Résumé du chapitre~\thechapter}\label{sec:resume-ch3}
%=====================================================
\addcontentsline{toc}{section}{Résumé du chapitre~\thechapter}

\begin{itemize}
  \item Un \textbf{ensemble} est caractérisé par ses éléments ; l'ensemble des
    parties $\powerset(A)$ a $2^{\card{A}}$ éléments.
  \item Les lois de \textbf{De Morgan} relient union, intersection et complémentaire.
  \item Une \textbf{relation d'équivalence} (réflexive, symétrique, transitive)
    partitionne son ensemble en classes d'équivalence.
  \item Un \textbf{ordre partiel} (réflexif, antisymétrique, transitif) se
    visualise par un diagramme de Hasse.
  \item Une fonction est \textbf{bijective} si et seulement si elle admet un
    inverse.
  \item L'argument diagonal de \textbf{Cantor} montre que $\R$ est
    indénombrable : $\card{\N} < \card{\R}$.
\end{itemize}

\bigskip
\hrule
\bigskip


%%%%%%%%%%%%%%%%%%%%%%%%%%%%%%%%%%%%%%%%%%%%%%%%%%%%%%%%%%%%%
%  CHAPITRE 4 — Combinatoire : Dénombrement
%%%%%%%%%%%%%%%%%%%%%%%%%%%%%%%%%%%%%%%%%%%%%%%%%%%%%%%%%%%%%
\chapter{Combinatoire --- Dénombrement}\label{ch:combinatoire}
\index{combinatoire}\index{dénombrement}

\begin{center}
\textit{«~L'art de compter est l'art de ne compter qu'une fois.~»}
\end{center}

\bigskip

Le dénombrement --- l'art de compter le nombre d'éléments d'un ensemble fini
--- est une compétence fondamentale en mathématiques discrètes, en
informatique et en probabilités. Ce chapitre développe les techniques
classiques de combinatoire.

%=====================================================
\section{Principes fondamentaux du dénombrement}\label{sec:principes-denombrement}
%=====================================================

\begin{theoreme}[Principe d'addition]\label{thm:principe-addition}
\index{principe d'addition}
Si $A_1, A_2, \ldots, A_k$ sont des ensembles finis \emph{deux à deux
disjoints}, alors
\[
  \card{A_1 \cup A_2 \cup \cdots \cup A_k}
  = \card{A_1} + \card{A_2} + \cdots + \card{A_k}.
\]
\end{theoreme}

\begin{proof}
Par récurrence sur $k$. Pour $k=1$, c'est trivial. Si le résultat vaut
pour $k-1$ ensembles, alors
$\card{A_1 \cup \cdots \cup A_k} = \card{(A_1 \cup \cdots \cup A_{k-1}) \cup A_k}$.
Comme $A_k$ est disjoint de $A_1 \cup \cdots \cup A_{k-1}$, on obtient
$\card{A_1 \cup \cdots \cup A_{k-1}} + \card{A_k}$, et on conclut par
hypothèse de récurrence.
\end{proof}

\begin{theoreme}[Principe de multiplication]\label{thm:principe-multiplication}
\index{principe de multiplication}
Si une procédure se décompose en $k$ étapes successives, la $i$-ème étape
pouvant être réalisée de $n_i$ façons (indépendamment des choix précédents),
alors le nombre total de réalisations est $n_1 \cdot n_2 \cdots n_k$.
Formellement :
\[
  \card{A_1 \times A_2 \times \cdots \times A_k}
  = \card{A_1} \cdot \card{A_2} \cdots \card{A_k}.
\]
\end{theoreme}

\begin{proof}
Par récurrence sur $k$. Pour $k=1$, c'est trivial. Supposons le résultat
vrai pour $k-1$ ensembles. On peut écrire
$A_1 \times \cdots \times A_k = (A_1 \times \cdots \times A_{k-1}) \times A_k$.
Pour chacun des $\card{A_1} \cdots \card{A_{k-1}}$ éléments du premier
facteur (par hypothèse de récurrence), il y a $\card{A_k}$ choix pour la
dernière composante, d'où le résultat.
\end{proof}

\begin{exemple}[Plaques d'immatriculation]\label{ex:plaques}
Une plaque est composée de 2~lettres suivies de 3~chiffres puis de 2~lettres
(format AA-000-AA). Le nombre de plaques possibles est
$26^2 \cdot 10^3 \cdot 26^2 = 26^4 \cdot 10^3 = 456\,976 \cdot 1000
= 456\,976\,000$.
\end{exemple}

%=====================================================
\section{Permutations et arrangements}\label{sec:permutations}
%=====================================================

\begin{definition}[Factorielle]\label{def:factorielle}
\index{factorielle}
Pour $n \in \N$, la \textbf{factorielle} de $n$ est
\[
  n! = \begin{cases} 1 & \text{si } n = 0, \\ n \cdot (n-1)! & \text{si } n \geq 1. \end{cases}
\]
\end{definition}

\begin{definition}[Permutation]\label{def:permutation}
\index{permutation}
Une \textbf{permutation} d'un ensemble à $n$ éléments est un arrangement
ordonné de \emph{tous} ses éléments. Le nombre de permutations est $n!$.
\end{definition}

\begin{definition}[Arrangement (permutation partielle)]\label{def:arrangement}
\index{arrangement}
Le nombre de façons de choisir $k$ éléments \emph{ordonnés} parmi $n$ est
\[
  A(n,k) = \frac{n!}{(n-k)!} = n(n-1)\cdots(n-k+1).
\]
\end{definition}

\begin{exemple}\label{ex:podium}
Lors d'une course de 10~athlètes, le nombre de podiums possibles
(or, argent, bronze) est $A(10,3) = 10 \cdot 9 \cdot 8 = 720$.
\end{exemple}

%=====================================================
\section{Combinaisons et coefficients binomiaux}\label{sec:combinaisons}
%=====================================================

\begin{definition}[Combinaison, coefficient binomial]\label{def:combinaison}
\index{combinaison}\index{coefficient binomial}\index{binomial!coefficient}
Le nombre de sous-ensembles à $k$ éléments d'un ensemble à $n$ éléments est
le \textbf{coefficient binomial}
\[
  \binom{n}{k} = \frac{n!}{k!\,(n-k)!}, \quad 0 \leq k \leq n.
\]
Par convention, $\binom{n}{k} = 0$ si $k < 0$ ou $k > n$.
\end{definition}

\begin{exemple}[Calculs de coefficients binomiaux]\label{ex:binom-calculs}
\[
  \binom{10}{3} = \frac{10!}{3!\,7!} = \frac{10 \cdot 9 \cdot 8}{3 \cdot 2 \cdot 1} = 120, \qquad
  \binom{8}{5} = \binom{8}{3} = \frac{8 \cdot 7 \cdot 6}{6} = 56.
\]
\end{exemple}

\begin{proposition}[Propriétés élémentaires]\label{prop:binom-elem}
\begin{enumerate}
  \item Symétrie : $\displaystyle\binom{n}{k} = \binom{n}{n-k}$.
  \item Bords : $\displaystyle\binom{n}{0} = \binom{n}{n} = 1$.
  \item Formule de Pascal\index{Pascal!formule de} :
    $\displaystyle\binom{n}{k} = \binom{n-1}{k-1} + \binom{n-1}{k}$
    pour $1 \leq k \leq n-1$.
\end{enumerate}
\end{proposition}

\begin{proof}
(3) On interprète combinatoirement. Parmi les $\binom{n}{k}$ sous-ensembles
de $\{1,\ldots,n\}$ à $k$ éléments, ceux contenant $n$ sont au nombre de
$\binom{n-1}{k-1}$ (on choisit les $k-1$ restants parmi $\{1,\ldots,n-1\}$)
et ceux ne contenant pas $n$ sont au nombre de $\binom{n-1}{k}$.
\end{proof}

%-----------------------------------------------------
\subsection{Triangle de Pascal}
%-----------------------------------------------------

\begin{figure}[ht]
\centering
\begin{tikzpicture}[scale=0.85]
  \foreach \n in {0,...,7} {
    \foreach \k in {0,...,\n} {
      \pgfmathtruncatemacro{\coeff}{
        ifthenelse(\k==0 || \k==\n, 1,
        ifthenelse(\n==2 && \k==1, 2,
        ifthenelse(\n==3 && \k==1, 3,
        ifthenelse(\n==3 && \k==2, 3,
        ifthenelse(\n==4 && \k==1, 4,
        ifthenelse(\n==4 && \k==2, 6,
        ifthenelse(\n==4 && \k==3, 4,
        ifthenelse(\n==5 && \k==1, 5,
        ifthenelse(\n==5 && \k==2, 10,
        ifthenelse(\n==5 && \k==3, 10,
        ifthenelse(\n==5 && \k==4, 5,
        ifthenelse(\n==6 && \k==1, 6,
        ifthenelse(\n==6 && \k==2, 15,
        ifthenelse(\n==6 && \k==3, 20,
        ifthenelse(\n==6 && \k==4, 15,
        ifthenelse(\n==6 && \k==5, 6,
        ifthenelse(\n==7 && \k==1, 7,
        ifthenelse(\n==7 && \k==2, 21,
        ifthenelse(\n==7 && \k==3, 35,
        ifthenelse(\n==7 && \k==4, 35,
        ifthenelse(\n==7 && \k==5, 21,
        ifthenelse(\n==7 && \k==6, 7, 0))))))))))))))))))))))
      }
      \node at ({(\k - \n/2)*1.1}, {-\n*0.9}) {\coeff};
    }
  }
\end{tikzpicture}
\caption{Les huit premières lignes du triangle de Pascal
($n = 0, 1, \ldots, 7$).}
\label{fig:triangle-pascal}
\end{figure}

%=====================================================
\section{Théorème binomial}\label{sec:binomial}
%=====================================================

\begin{theoreme}[Théorème binomial (binôme de Newton)]\label{thm:binomial}
\index{binôme de Newton}\index{théorème binomial}
Pour tous $a, b$ et tout $n \in \N$ :
\[
  (a+b)^n = \sum_{k=0}^{n} \binom{n}{k} a^{n-k} b^k.
\]
\end{theoreme}

\begin{proof}[Preuve par récurrence]
\textbf{Base.} Pour $n=0$ : $(a+b)^0 = 1 = \binom{0}{0} a^0 b^0$. \checkmark

\textbf{Hérédité.} Supposons la formule vraie au rang $n$. Alors :
\begin{align*}
  (a+b)^{n+1}
  &= (a+b)\sum_{k=0}^{n}\binom{n}{k}a^{n-k}b^k \\
  &= \sum_{k=0}^{n}\binom{n}{k}a^{n+1-k}b^k
     + \sum_{k=0}^{n}\binom{n}{k}a^{n-k}b^{k+1} \\
  &= \sum_{k=0}^{n}\binom{n}{k}a^{n+1-k}b^k
     + \sum_{j=1}^{n+1}\binom{n}{j-1}a^{n+1-j}b^{j}
     \quad (j = k+1) \\
  &= a^{n+1} + \sum_{k=1}^{n}\left[\binom{n}{k}+\binom{n}{k-1}\right]a^{n+1-k}b^k + b^{n+1} \\
  &= a^{n+1} + \sum_{k=1}^{n}\binom{n+1}{k}a^{n+1-k}b^k + b^{n+1}
     \quad \text{(Pascal)} \\
  &= \sum_{k=0}^{n+1}\binom{n+1}{k}a^{n+1-k}b^k. \qedhere
\end{align*}
\end{proof}

\begin{corollaire}\label{cor:somme-binom}
En posant $a = b = 1$ :
$\displaystyle\sum_{k=0}^{n}\binom{n}{k} = 2^n$.
En posant $a=1$, $b=-1$ :
$\displaystyle\sum_{k=0}^{n}(-1)^k\binom{n}{k} = 0$.
\end{corollaire}

%=====================================================
\section{Identité de Vandermonde}\label{sec:vandermonde}
%=====================================================

\begin{theoreme}[Identité de Vandermonde]\label{thm:vandermonde}
\index{Vandermonde!identité de}
Pour tous entiers $m, n, r \geq 0$ :
\[
  \binom{m+n}{r} = \sum_{k=0}^{r} \binom{m}{k}\binom{n}{r-k}.
\]
\end{theoreme}

\begin{proof}[Preuve combinatoire]
Considérons un ensemble de $m+n$ personnes, divisé en un groupe $A$ de $m$
personnes et un groupe $B$ de $n$ personnes. Le nombre de façons de choisir
$r$ personnes parmi $m+n$ est $\binom{m+n}{r}$.

D'autre part, si l'on choisit exactement $k$ personnes dans $A$ et $r-k$
dans $B$, cela peut se faire de $\binom{m}{k}\binom{n}{r-k}$ façons. En
sommant sur toutes les valeurs possibles de $k$ (de $0$ à $r$), on obtient
le membre de droite.
\end{proof}

\begin{corollaire}\label{cor:vandermonde-cas-particulier}
En posant $m = n = r$ :
$\displaystyle\sum_{k=0}^{n}\binom{n}{k}^2 = \binom{2n}{n}$.
\end{corollaire}

%=====================================================
\section{Combinaisons avec répétition}\label{sec:multiset}
%=====================================================

\begin{definition}[Combinaison avec répétition]\label{def:combinaison-repetition}
\index{combinaison!avec répétition}\index{multiensemble}
Le nombre de façons de choisir $k$ éléments parmi $n$ types \emph{avec
répétition autorisée} (l'ordre ne comptant pas) est le
\textbf{coefficient multiensemble}\index{coefficient multiensemble} :
\[
  \left(\!\!\binom{n}{k}\!\!\right)
  = \binom{n+k-1}{k}
  = \binom{n+k-1}{n-1}.
\]
\end{definition}

\begin{exemple}[Boules de glace]\label{ex:glaces}
Un glacier propose 5~parfums. Le nombre de coupes de 3~boules (avec
répétition possible, sans ordre) est
\[
  \left(\!\!\binom{5}{3}\!\!\right) = \binom{5+3-1}{3} = \binom{7}{3} = 35.
\]
\end{exemple}

%=====================================================
\section{Permutations avec répétition et anagrammes}\label{sec:perm-repetition}
%=====================================================

\begin{theoreme}[Permutations avec répétition]\label{thm:perm-repetition}
\index{permutation!avec répétition}\index{coefficient multinomial}
Le nombre de permutations de $n$ objets où l'objet de type $i$ apparaît
$n_i$ fois ($n_1 + n_2 + \cdots + n_r = n$) est le \textbf{coefficient
multinomial} :
\[
  \binom{n}{n_1, n_2, \ldots, n_r} = \frac{n!}{n_1!\, n_2! \cdots n_r!}.
\]
\end{theoreme}

\begin{proof}
Il y a $n!$ façons d'ordonner $n$ objets distincts. Or les $n_i!$
permutations internes des objets de type $i$ produisent le même arrangement.
On divise donc par $n_1!\, n_2! \cdots n_r!$.
\end{proof}

\begin{exemple}[Anagrammes]\label{ex:anagrammes}
Le nombre d'anagrammes du mot \textsc{MISSISSIPPI} est :
\[
  \frac{11!}{1!\, 4!\, 4!\, 2!} = \frac{39\,916\,800}{1 \cdot 24 \cdot 24 \cdot 2}
  = 34\,650.
\]
En effet : M apparaît 1~fois, I~4~fois, S~4~fois, P~2~fois.
\end{exemple}

%=====================================================
\section{Méthode des boules et barres}\label{sec:boules-barres}
%=====================================================

\begin{theoreme}[Principe d'inclusion-exclusion]\label{thm:inclusion-exclusion}
\index{inclusion-exclusion!principe}
Pour des ensembles finis $A_1, \ldots, A_n$ :
\[
  \card{A_1 \cup \cdots \cup A_n}
  = \sum_{i}\card{A_i}
  - \sum_{i<j}\card{A_i \cap A_j}
  + \sum_{i<j<k}\card{A_i \cap A_j \cap A_k}
  - \cdots
  + (-1)^{n+1}\card{A_1 \cap \cdots \cap A_n}.
\]
\end{theoreme}

\begin{proof}
Montrons que chaque élément $x \in A_1 \cup \cdots \cup A_n$ est compté
exactement une fois dans le membre de droite. Supposons que $x$ appartient à
exactement $m$ des ensembles ($m \geq 1$). Sa contribution au membre de
droite est
\[
  \binom{m}{1} - \binom{m}{2} + \binom{m}{3} - \cdots + (-1)^{m+1}\binom{m}{m}
  = -\sum_{k=0}^{m}(-1)^k\binom{m}{k} + 1 = -(1-1)^m + 1 = 1,
\]
par le théorème binomial. Donc chaque élément est compté exactement une fois.
\end{proof}

\begin{exemple}[Inclusion-exclusion pour deux ensembles]\label{ex:incl-excl-2}
Pour deux ensembles finis : $\card{A \cup B} = \card{A} + \card{B} - \card{A \cap B}$.
\end{exemple}

% --- Diagramme de Venn pour inclusion-exclusion ---
\begin{figure}[ht]
\centering
\begin{tikzpicture}[scale=0.8]
  \draw[fill=red!15, draw=red!60, thick] (210:1.2) circle (1.8);
  \draw[fill=green!15, draw=green!60, thick] (330:1.2) circle (1.8);
  \draw[fill=blue!15, draw=blue!60, thick] (90:1.2) circle (1.8);
  \node at (210:2.5) {$A_1$};
  \node at (330:2.5) {$A_2$};
  \node at (90:2.5) {$A_3$};
  \node[font=\footnotesize] at (270:0.4) {$A_1 \!\cap\! A_2$};
  \node[font=\footnotesize] at (30:0.6) {$A_2 \!\cap\! A_3$};
  \node[font=\footnotesize] at (150:0.6) {$A_1 \!\cap\! A_3$};
  \node[font=\footnotesize] at (0:0) {$\cap$};
\end{tikzpicture}
\caption{Diagramme de Venn pour trois ensembles (inclusion-exclusion).}
\label{fig:venn-incl-excl}
\end{figure}

\begin{theoreme}[Boules et barres (stars and bars)]\label{thm:boules-barres}
\index{boules et barres}\index{stars and bars}
Le nombre de solutions entières de
\[
  x_1 + x_2 + \cdots + x_k = n, \quad x_i \geq 0,
\]
est $\displaystyle\binom{n+k-1}{k-1}$.
\end{theoreme}

\begin{proof}
Représentons chaque solution par une suite de $n$ étoiles (\textasteriskcentered)
et $k-1$ barres ($|$). Les étoiles avant la première barre représentent
$x_1$, celles entre les barres $i-1$ et $i$ représentent $x_i$, etc.

Par exemple, pour $k=3$, $n=5$, la configuration
$\textasteriskcentered\textasteriskcentered\,|\,\textasteriskcentered\,|\,\textasteriskcentered\textasteriskcentered$
correspond à $(x_1, x_2, x_3) = (2, 1, 2)$.

Une telle configuration est une suite de $n + k - 1$ symboles dont $k-1$
sont des barres. Le nombre de configurations est donc
$\binom{n+k-1}{k-1}$.
\end{proof}

\begin{remarque}\label{rem:boules-barres-strict}
Si l'on impose $x_i \geq 1$, le changement de variable $y_i = x_i - 1$
ramène au cas $y_1 + \cdots + y_k = n - k$ avec $y_i \geq 0$, donnant
$\binom{n-1}{k-1}$ solutions.
\end{remarque}

\begin{exemple}\label{ex:boules-barres}
Combien y a-t-il de façons de répartir 10~bonbons identiques entre 4~enfants ?
Réponse : $\binom{10+4-1}{4-1} = \binom{13}{3} = 286$.
\end{exemple}

% --- Diagramme boules et barres ---
\begin{figure}[ht]
\centering
\begin{tikzpicture}[scale=0.6]
  \foreach \i in {1,...,5} {
    \draw[fill=blue!30] (\i*0.8, 0) circle (0.3);
  }
  \draw[very thick, red] (5*0.8+0.6, -0.5) -- (5*0.8+0.6, 0.5);
  \foreach \i in {1,...,3} {
    \pgfmathsetmacro{\xpos}{5*0.8+0.6+0.2+\i*0.8}
    \draw[fill=blue!30] (\xpos, 0) circle (0.3);
  }
  \draw[very thick, red] (5*0.8+0.6+0.2+3*0.8+0.4, -0.5) -- (5*0.8+0.6+0.2+3*0.8+0.4, 0.5);
  \foreach \i in {1,...,2} {
    \pgfmathsetmacro{\xpos}{5*0.8+0.6+0.2+3*0.8+0.4+0.2+\i*0.8}
    \draw[fill=blue!30] (\xpos, 0) circle (0.3);
  }
  \node at (6, -1.3) {$\textasteriskcentered\textasteriskcentered\textasteriskcentered\textasteriskcentered\textasteriskcentered
    \;\mid\; \textasteriskcentered\textasteriskcentered\textasteriskcentered
    \;\mid\; \textasteriskcentered\textasteriskcentered
    \;\longleftrightarrow\; (5,3,2)$};
\end{tikzpicture}
\caption{Illustration de la méthode des boules et barres pour $n=10$, $k=3$.}
\label{fig:boules-barres}
\end{figure}

%=====================================================
\section{Dérangements}\label{sec:derangements}
%=====================================================

\begin{definition}[Dérangement]\label{def:derangement}
\index{dérangement}\index{sous-factorielle}
Un \textbf{dérangement} d'un ensemble $\{1, 2, \ldots, n\}$ est une
permutation $\sigma$ sans point fixe, c'est-à-dire telle que
$\sigma(i) \neq i$ pour tout $i$. Le nombre de dérangements de $n$ éléments
est noté $D_n$ (ou $!n$, la \textbf{sous-factorielle}).
\end{definition}

\begin{theoreme}[Formule des dérangements]\label{thm:derangements}
\index{dérangement!formule}
Pour tout $n \geq 0$ :
\[
  D_n = n! \sum_{k=0}^{n} \frac{(-1)^k}{k!}.
\]
En particulier, $D_0 = 1$, $D_1 = 0$, $D_2 = 1$, $D_3 = 2$, $D_4 = 9$,
$D_5 = 44$.
\end{theoreme}

\begin{proof}
Soit $A_i$ l'ensemble des permutations de $\{1,\ldots,n\}$ qui fixent $i$.
Par le principe d'inclusion-exclusion\index{inclusion-exclusion} :
\begin{align*}
  D_n &= n! - \card{A_1 \cup A_2 \cup \cdots \cup A_n} \\
  &= n! - \sum_{i}\card{A_i} + \sum_{i<j}\card{A_i \cap A_j}
    - \cdots + (-1)^n \card{A_1 \cap \cdots \cap A_n}.
\end{align*}
Or $\card{A_{i_1} \cap \cdots \cap A_{i_k}} = (n-k)!$ (on fixe les positions
$i_1, \ldots, i_k$ et on permute les $n-k$ restantes). Il y a $\binom{n}{k}$
façons de choisir les $k$ indices. Donc :
\[
  D_n = \sum_{k=0}^{n} (-1)^k \binom{n}{k}(n-k)!
  = \sum_{k=0}^{n} (-1)^k \frac{n!}{k!}
  = n! \sum_{k=0}^{n} \frac{(-1)^k}{k!}. \qedhere
\]
\end{proof}

\begin{remarque}\label{rem:derangement-approx}
Comme $\displaystyle\sum_{k=0}^{\infty}\frac{(-1)^k}{k!} = e^{-1}$, on a
l'approximation $D_n \approx n!/e$, et plus précisément
$D_n$ est l'entier le plus proche de $n!/e$.
\end{remarque}

\begin{proposition}[Récurrence pour $D_n$]\label{prop:derangement-rec}
\index{dérangement!récurrence}
Pour $n \geq 2$ :
\[
  D_n = (n-1)(D_{n-1} + D_{n-2}).
\]
\end{proposition}

\begin{proof}
Considérons un dérangement $\sigma$ de $\{1,\ldots,n\}$. Posons
$\sigma(1) = j$ avec $j \neq 1$ (il y a $n-1$ choix pour $j$). Deux cas :
\begin{itemize}
  \item Si $\sigma(j) = 1$ : les éléments $1$ et $j$ s'échangent, et les
    $n-2$ restants forment un dérangement. Il y a $D_{n-2}$ possibilités.
  \item Si $\sigma(j) \neq 1$ : on peut interpréter la restriction aux
    éléments $\{2,\ldots,n\}$ comme un dérangement de $n-1$ éléments (en
    remplaçant la contrainte $\sigma(j)\neq 1$ par $\sigma(j)\neq j$). Il y
    a $D_{n-1}$ possibilités.
\end{itemize}
Au total, $D_n = (n-1)(D_{n-1} + D_{n-2})$.
\end{proof}

%=====================================================
\section{Exemples résolus}\label{sec:exemples-denombrement}
%=====================================================

\begin{exemple}[Sous-ensembles d'un ensemble]\label{ex:sous-ensembles-comptage}
Combien de sous-ensembles de $\{1,2,\ldots,10\}$ contiennent exactement
3~éléments pairs et 2~éléments impairs ?

\textbf{Solution.} Il y a 5~nombres pairs ($2,4,6,8,10$) et 5~nombres impairs
($1,3,5,7,9$). Le nombre de sous-ensembles cherchés est
$\binom{5}{3}\binom{5}{2} = 10 \cdot 10 = 100$.
\end{exemple}

\begin{exemple}[Diagonales d'un polygone]\label{ex:diagonales}
\index{diagonales d'un polygone}
Un polygone convexe à $n$ sommets possède $\binom{n}{2} - n = \frac{n(n-3)}{2}$
diagonales. En effet, $\binom{n}{2}$ est le nombre total de segments reliant
deux sommets, et $n$ d'entre eux sont des côtés.
\end{exemple}

% --- Diagramme arborescent de comptage ---
\begin{figure}[ht]
\centering
\begin{tikzpicture}[
    level distance=1.5cm,
    level 1/.style={sibling distance=4cm},
    level 2/.style={sibling distance=1.3cm},
    every node/.style={draw, rounded corners, minimum width=1cm, minimum height=0.6cm, font=\small},
    edge from parent/.style={draw, ->, >=stealth}
  ]
  \node {Entrée}
    child { node {V}
      child { node {VV} }
      child { node {VC} }
      child { node {VCh} }
    }
    child { node {C}
      child { node {CV} }
      child { node {CC} }
      child { node {CCh} }
    }
    child { node {Ch}
      child { node {ChV} }
      child { node {ChC} }
      child { node {ChCh} }
    };
\end{tikzpicture}
\caption{Arbre de dénombrement pour un menu à 2~plats choisis parmi viande (V),
poisson (C) et champignons (Ch) --- principe de multiplication : $3 \times 3 = 9$.}
\label{fig:arbre-comptage}
\end{figure}

\begin{exemple}[Comité avec conditions]\label{ex:comite}
Dans un groupe de 8~hommes et 6~femmes, on veut former un comité de
5~personnes contenant au moins 2~femmes. Combien de comités possibles ?

\textbf{Solution.} On compte selon le nombre de femmes $f \in \{2,3,4,5\}$ :
\[
  \sum_{f=2}^{5} \binom{6}{f}\binom{8}{5-f}
  = \binom{6}{2}\binom{8}{3} + \binom{6}{3}\binom{8}{2}
  + \binom{6}{4}\binom{8}{1} + \binom{6}{5}\binom{8}{0}.
\]
Calcul :
$15 \cdot 56 + 20 \cdot 28 + 15 \cdot 8 + 6 \cdot 1
= 840 + 560 + 120 + 6 = 1526$.
\end{exemple}

\begin{exemple}[Chemins sur une grille]\label{ex:chemins-grille}
\index{chemins sur une grille}
Combien de chemins mènent du point $(0,0)$ au point $(m,n)$ sur une grille,
si l'on ne peut se déplacer que vers la droite (D) ou vers le haut (H) ?

\textbf{Solution.} Chaque chemin est une suite de $m$ pas D et $n$ pas H,
soit une suite de $m+n$ pas dont $m$ sont D. Le nombre de chemins est
$\binom{m+n}{m}$.

Par exemple, de $(0,0)$ à $(4,3)$ : $\binom{7}{4} = 35$ chemins.
\end{exemple}

\begin{exemple}[Distribution dans des boîtes distinctes]\label{ex:distribution}
De combien de façons peut-on répartir 12~livres \emph{identiques} en
4~étagères (pouvant être vides) ?

\textbf{Solution.} C'est le problème $x_1 + x_2 + x_3 + x_4 = 12$ avec
$x_i \geq 0$. Par la méthode des boules et barres :
$\binom{12+4-1}{4-1} = \binom{15}{3} = 455$.
\end{exemple}

\begin{exemple}[Mots de passe]\label{ex:mdp}
Un mot de passe contient exactement 8~caractères choisis parmi 26~lettres
minuscules et 10~chiffres. Il doit contenir au moins un chiffre. Combien de
mots de passe valides ?

\textbf{Solution.} Total de mots de passe de longueur 8 : $36^8$. Ceux sans
chiffre : $26^8$. Donc le nombre de mots de passe valides est
$36^8 - 26^8 = 2\,821\,109\,907\,456 - 208\,827\,064\,576 = 2\,612\,282\,842\,880$.
\end{exemple}

\begin{exemple}[Le problème des lettres et enveloppes]\label{ex:lettres-enveloppes}
\index{problème des lettres et enveloppes}
On place au hasard 5~lettres dans 5~enveloppes adressées. Quelle est la
probabilité qu'aucune lettre ne soit dans la bonne enveloppe ?

\textbf{Solution.} C'est le problème des dérangements. Le nombre de
dérangements de 5~éléments est
$D_5 = 5!\left(1 - 1 + \frac{1}{2} - \frac{1}{6} + \frac{1}{24} - \frac{1}{120}\right) = 120 \cdot \frac{11}{30} = 44$.
La probabilité est $D_5/5! = 44/120 = 11/30 \approx 0{,}367$.
\end{exemple}

\begin{exemple}[Fonction surjective]\label{ex:surjection-comptage}
\index{surjection!dénombrement}
Le nombre de fonctions surjectives de $\{1,\ldots,n\}$ vers $\{1,\ldots,k\}$
($n \geq k$) est donné par la formule d'inclusion-exclusion :
\[
  S(n,k) = \sum_{j=0}^{k}(-1)^j \binom{k}{j}(k-j)^n.
\]
Par exemple, le nombre de surjections de $\{1,\ldots,4\}$ vers $\{1,2,3\}$ est
\[
  \binom{3}{0}3^4 - \binom{3}{1}2^4 + \binom{3}{2}1^4 - \binom{3}{3}0^4
  = 81 - 48 + 3 - 0 = 36.
\]
\end{exemple}

%=====================================================
\section{Exercices}\label{sec:exo-ch4}
%=====================================================

\begin{exercice}[Principes de base]\label{exo:principes}
\begin{enumerate}
  \item Combien de chaînes binaires de longueur 10 commencent par 11 ou
    finissent par 00 ?
  \item Combien de nombres entre 1 et 1000 sont divisibles par 3 ou par 5 ?
\end{enumerate}
\end{exercice}

\begin{exercice}[Permutations]\label{exo:permutations}
De combien de façons peut-on asseoir 8~personnes autour d'une table ronde
(deux dispositions sont identiques si l'on peut passer de l'une à l'autre
par rotation) ?
\end{exercice}

\begin{exercice}[Identités binomiales]\label{exo:identites-binom}
Démontrer les identités suivantes :
\begin{enumerate}
  \item $\displaystyle\sum_{k=0}^{n} k\binom{n}{k} = n\,2^{n-1}$.
  \item $\displaystyle\sum_{k=0}^{n} \binom{n}{k}^2 = \binom{2n}{n}$.
  \item $\displaystyle\binom{n}{k}\binom{k}{j} = \binom{n}{j}\binom{n-j}{k-j}$
    pour $0 \leq j \leq k \leq n$.
\end{enumerate}
\end{exercice}

\begin{exercice}[Anagrammes avec contraintes]\label{exo:anagrammes}
\begin{enumerate}
  \item Combien d'anagrammes du mot \textsc{ARRANGEMENT} ?
  \item Combien d'anagrammes du mot \textsc{COMBINATOIRE} ne commencent
    pas par une voyelle ?
\end{enumerate}
\end{exercice}

\begin{exercice}[Boules et barres]\label{exo:boules-barres}
\begin{enumerate}
  \item Combien de solutions entières a l'équation $x_1+x_2+x_3+x_4 = 20$
    avec $x_i \geq 0$ ?
  \item Même question avec $x_i \geq 2$ pour tout $i$.
  \item Combien de solutions avec $0 \leq x_i \leq 5$ ?
    \textit{(Utiliser l'inclusion-exclusion.)}
\end{enumerate}
\end{exercice}

\begin{exercice}[Dérangements]\label{exo:derangements}
\begin{enumerate}
  \item Calculer $D_6$ et $D_7$ en utilisant la formule de récurrence.
  \item Montrer que $D_n = nD_{n-1} + (-1)^n$ pour $n \geq 1$.
  \item Soit $p_n = D_n/n!$ la probabilité qu'une permutation aléatoire soit
    un dérangement. Montrer que $\lim_{n\to\infty} p_n = 1/e$.
\end{enumerate}
\end{exercice}

\begin{exercice}[Chemins avec obstacles]\label{exo:chemins-obstacles}
\index{chemins!avec obstacles}
Combien de chemins sur grille de $(0,0)$ à $(6,6)$ (pas D et H) passent par
le point $(3,3)$ ?
\end{exercice}

\begin{exercice}[Dénombrement avancé]\label{exo:denombrement-avance}
On dispose de $n$ billes \emph{distinctes} et de $k$ boîtes \emph{distinctes}.
\begin{enumerate}
  \item Combien de répartitions si les boîtes peuvent être vides ?
  \item Combien si chaque boîte contient au plus une bille ($n \leq k$) ?
  \item Combien si chaque boîte contient au moins une bille ($n \geq k$) ?
\end{enumerate}
\end{exercice}

\begin{exercice}[Identité de Vandermonde généralisée]\label{exo:vandermonde-gen}
\index{Vandermonde!identité généralisée}
En utilisant le théorème binomial appliqué à $(1+x)^m(1+x)^n = (1+x)^{m+n}$,
retrouver l'identité de Vandermonde
$\displaystyle\binom{m+n}{r} = \sum_{k=0}^{r}\binom{m}{k}\binom{n}{r-k}$.
\end{exercice}

\begin{exercice}[Surjections]\label{exo:surjections}
\begin{enumerate}
  \item Combien de surjections y a-t-il de $\{1,\ldots,5\}$ vers $\{a,b,c\}$ ?
  \item Montrer que le nombre de surjections de $\{1,\ldots,n\}$ vers
    $\{1,\ldots,n\}$ est $n!$.
\end{enumerate}
\end{exercice}

\begin{exercice}[Nombres à chiffres distincts]\label{exo:chiffres-distincts}
Combien de nombres de 4~chiffres (entre 1000 et 9999) ont tous leurs chiffres
distincts ?
\end{exercice}

\begin{exercice}[Problème de Lucas]\label{exo:lucas}
\index{problème des ménages}
Combien y a-t-il de façons de choisir $k$ éléments non consécutifs dans
$\{1, 2, \ldots, n\}$ (disposés en ligne) ? Montrer que la réponse est
$\binom{n-k+1}{k}$.
\end{exercice}

\begin{exercice}[Problème du chapelier]\label{exo:chapelier}
\index{problème du chapelier}
$n$ personnes déposent leur chapeau au vestiaire. Les chapeaux sont rendus
au hasard. Calculer la probabilité qu'au moins une personne reçoive son propre
chapeau. \textit{(Utiliser l'inclusion-exclusion et les dérangements.)}
\end{exercice}

%=====================================================
\section*{Résumé du chapitre~\thechapter}\label{sec:resume-ch4}
%=====================================================
\addcontentsline{toc}{section}{Résumé du chapitre~\thechapter}

\begin{itemize}
  \item Les \textbf{principes d'addition et de multiplication} sont les outils
    fondamentaux du dénombrement.
  \item \textbf{Permutations} : $n!$ ordres de $n$ éléments ;
    \textbf{arrangements} : $A(n,k) = n!/(n-k)!$.
  \item \textbf{Combinaisons} : $\binom{n}{k} = \frac{n!}{k!(n-k)!}$ ; la
    formule de Pascal donne le triangle du même nom.
  \item Le \textbf{théorème binomial} :
    $(a+b)^n = \sum_{k}\binom{n}{k}a^{n-k}b^k$.
  \item L'\textbf{identité de Vandermonde} :
    $\binom{m+n}{r} = \sum_k \binom{m}{k}\binom{n}{r-k}$.
  \item \textbf{Combinaisons avec répétition} :
    $\binom{n+k-1}{k}$ choix de $k$ éléments parmi $n$ types.
  \item \textbf{Permutations avec répétition} (multinomiales) :
    $\frac{n!}{n_1!\cdots n_r!}$.
  \item La méthode des \textbf{boules et barres} résout les équations
    $x_1 + \cdots + x_k = n$ en entiers non négatifs :
    $\binom{n+k-1}{k-1}$ solutions.
  \item Les \textbf{dérangements} $D_n = n!\sum_{k=0}^{n}(-1)^k/k!$
    comptent les permutations sans point fixe ; $D_n \approx n!/e$.
\end{itemize}

% === FIN DU CHUNK ch3-4 ===
% === DÉBUT DU CHUNK ch5-6 ===

%%%%%%%%%%%%%%%%%%%%%%%%%%%%%%%%%%%%%%%%%%%%%%%%%%%%%%%%%%%%%%%%%%%%
%  CHAPITRE 5 — Principe d'Inclusion-Exclusion et Dénombrement Avancé
%%%%%%%%%%%%%%%%%%%%%%%%%%%%%%%%%%%%%%%%%%%%%%%%%%%%%%%%%%%%%%%%%%%%
\chapter{Principe d'Inclusion-Exclusion et Dénombrement Avancé}
\label{ch:inclusion-exclusion}
\index{inclusion-exclusion}

\section{Le principe d'inclusion-exclusion}
\label{sec:ie-principe}

Le principe d'inclusion-exclusion\index{inclusion-exclusion!principe} est l'un des outils
les plus puissants du dénombrement. Il permet de calculer le cardinal de l'union
de plusieurs ensembles à partir des cardinaux de leurs intersections.

\subsection{Cas de deux ensembles}

\begin{theoreme}[Inclusion-exclusion pour deux ensembles]
\label{thm:ie-deux}
Soient $A$ et $B$ deux ensembles finis. Alors
\[
  \card{A \cup B} = \card{A} + \card{B} - \card{A \cap B}.
\]
\end{theoreme}

\begin{proof}
Chaque élément $x$ de $A \cup B$ appartient à exactement l'une des trois
régions disjointes : $A \setminus B$, $B \setminus A$ ou $A \cap B$.
On a donc
\[
  \card{A \cup B} = \card{A \setminus B} + \card{B \setminus A} + \card{A \cap B}.
\]
Or $\card{A} = \card{A \setminus B} + \card{A \cap B}$ et
$\card{B} = \card{B \setminus A} + \card{A \cap B}$, d'où
\[
  \card{A} + \card{B} = \card{A \setminus B} + \card{B \setminus A} + 2\card{A \cap B}.
\]
En soustrayant $\card{A \cap B}$, on obtient le résultat.
\end{proof}

% --- Diagramme de Venn pour 2 ensembles ---
\begin{figure}[ht]
\centering
\begin{tikzpicture}[scale=0.9]
  \begin{scope}[fill opacity=0.3]
    \fill[blue!50] (-0.8,0) circle (1.8);
    \fill[red!50] (0.8,0) circle (1.8);
  \end{scope}
  \draw (-0.8,0) circle (1.8) node[left=1.2cm] {$A$};
  \draw (0.8,0) circle (1.8) node[right=1.2cm] {$B$};
  \node at (-1.5,0) {$a$};
  \node at (0,0) {$c$};
  \node at (1.5,0) {$b$};
  \node[below] at (0,-2.2) {$\card{A \cup B} = a + b + c = \card{A}+\card{B}-\card{A\cap B}$};
\end{tikzpicture}
\caption{Diagramme de Venn pour deux ensembles : $a = \card{A \setminus B}$,
$b = \card{B \setminus A}$, $c = \card{A \cap B}$.}
\label{fig:venn-deux}
\end{figure}

\begin{exemple}[Divisibilité]
\label{ex:ie-div}
Combien d'entiers de $\set{1,2,\dots,100}$ sont divisibles par $2$ ou par $3$ ?
Posons $A = \set{n \in \set{1,\dots,100} : 2 \mid n}$ et
$B = \set{n \in \set{1,\dots,100} : 3 \mid n}$.
Alors $\card{A} = 50$, $\card{B} = 33$ et $\card{A \cap B} = \lfloor 100/6 \rfloor = 16$.
Par inclusion-exclusion,
\[
  \card{A \cup B} = 50 + 33 - 16 = 67.
\]
\end{exemple}

\subsection{Cas de trois ensembles}

\begin{theoreme}[Inclusion-exclusion pour trois ensembles]
\label{thm:ie-trois}
Soient $A$, $B$ et $C$ trois ensembles finis. Alors
\[
  \card{A \cup B \cup C}
    = \card{A} + \card{B} + \card{C}
    - \card{A \cap B} - \card{A \cap C} - \card{B \cap C}
    + \card{A \cap B \cap C}.
\]
\end{theoreme}

\begin{proof}
Appliquons le théorème~\ref{thm:ie-deux} à $A \cup B$ et $C$ :
\[
  \card{(A \cup B) \cup C}
    = \card{A \cup B} + \card{C} - \card{(A \cup B) \cap C}.
\]
Or $(A \cup B) \cap C = (A \cap C) \cup (B \cap C)$, donc
\[
  \card{(A \cup B) \cap C}
    = \card{A \cap C} + \card{B \cap C} - \card{A \cap B \cap C}.
\]
En combinant avec $\card{A \cup B} = \card{A} + \card{B} - \card{A \cap B}$,
on obtient la formule annoncée.
\end{proof}

% --- Diagramme de Venn pour 3 ensembles ---
\begin{figure}[ht]
\centering
\begin{tikzpicture}[scale=0.85]
  \begin{scope}[fill opacity=0.25]
    \fill[blue!60] (90:1.2) circle (2);
    \fill[red!60] (210:1.2) circle (2);
    \fill[green!60] (330:1.2) circle (2);
  \end{scope}
  \draw (90:1.2) circle (2) node[above=1.4cm] {$A$};
  \draw (210:1.2) circle (2) node[below left=1cm] {$B$};
  \draw (330:1.2) circle (2) node[below right=1cm] {$C$};
  % Labels des régions
  \node at (90:2.0) {\small $a$};
  \node at (210:2.0) {\small $b$};
  \node at (330:2.0) {\small $c$};
  \node at (150:0.7) {\small $d$};
  \node at (30:0.7) {\small $e$};
  \node at (270:0.7) {\small $f$};
  \node at (0,0.15) {\small $g$};
\end{tikzpicture}
\caption{Diagramme de Venn pour trois ensembles. Les sept régions correspondent
aux différentes intersections.}
\label{fig:venn-trois}
\end{figure}

\begin{exemple}[Langues parlées]
\label{ex:ie-langues}
Dans une classe de $100$ étudiants, $60$ parlent anglais, $45$ parlent espagnol
et $35$ parlent allemand. De plus, $20$ parlent anglais et espagnol, $15$ parlent
anglais et allemand, $12$ parlent espagnol et allemand, et $5$ parlent les trois langues.
Combien d'étudiants parlent au moins une de ces langues ?

Par inclusion-exclusion :
\[
  \card{A \cup E \cup G} = 60 + 45 + 35 - 20 - 15 - 12 + 5 = 98.
\]
Ainsi, $2$ étudiants ne parlent aucune de ces trois langues.
\end{exemple}

\subsection{Le cas général}

\begin{theoreme}[Principe d'inclusion-exclusion général]
\label{thm:ie-general}
\index{inclusion-exclusion!formule générale}
Soient $A_1, A_2, \dots, A_n$ des ensembles finis. Alors
\[
  \card*{\bigcup_{i=1}^{n} A_i}
    = \sum_{k=1}^{n} (-1)^{k+1}
      \sum_{1 \leqslant i_1 < i_2 < \cdots < i_k \leqslant n}
      \card{A_{i_1} \cap A_{i_2} \cap \cdots \cap A_{i_k}}.
\]
\end{theoreme}

\begin{proof}
Montrons que chaque élément $x \in \bigcup_{i=1}^n A_i$ est compté exactement
une fois dans le membre de droite. Supposons que $x$ appartient à exactement
$m$ ensembles parmi $A_1, \dots, A_n$, avec $m \geqslant 1$.
L'élément $x$ contribue au terme d'indice $k$ exactement $\binom{m}{k}$ fois
(il faut choisir $k$ ensembles parmi les $m$ qui le contiennent). Sa contribution
totale est
\[
  \sum_{k=1}^{m} (-1)^{k+1} \binom{m}{k}
    = -\sum_{k=1}^{m} (-1)^{k} \binom{m}{k}
    = -\left[\sum_{k=0}^{m} (-1)^{k} \binom{m}{k} - 1\right]
    = -(0 - 1) = 1,
\]
car $\sum_{k=0}^{m} (-1)^k \binom{m}{k} = (1-1)^m = 0$ pour $m \geqslant 1$
d'après le binôme de Newton.
\end{proof}

\begin{remarque}
\label{rem:ie-complementaire}
Le principe s'écrit souvent sous la forme complémentaire. Si tous les $A_i$
sont des sous-ensembles d'un univers fini $U$, alors le nombre d'éléments
n'appartenant à \emph{aucun} des $A_i$ est
\[
  \card*{U \setminus \bigcup_{i=1}^{n} A_i}
    = \card{U} - \sum_{k=1}^{n} (-1)^{k+1}
      \sum_{1 \leqslant i_1 < \cdots < i_k \leqslant n}
      \card{A_{i_1} \cap \cdots \cap A_{i_k}}.
\]
C'est cette forme que l'on utilise le plus souvent en pratique.
\end{remarque}

\section{Applications du principe d'inclusion-exclusion}
\label{sec:ie-applications}

\subsection{Dénombrement des surjections}
\index{surjection!dénombrement}

\begin{theoreme}[Nombre de surjections]
\label{thm:surjections}
Le nombre de surjections d'un ensemble à $n$ éléments vers un ensemble à $k$
éléments (avec $n \geqslant k$) est
\[
  S(n,k) = \sum_{j=0}^{k} (-1)^{j} \binom{k}{j} (k-j)^n.
\]
\end{theoreme}

\begin{proof}
Soit $B = \set{b_1, \dots, b_k}$ l'ensemble d'arrivée. Pour chaque $i$,
posons $A_i = \set{f : \set{1,\dots,n} \to B \mid b_i \notin \text{Im}(f)}$,
l'ensemble des fonctions qui \og{}évitent\fg{} $b_i$.
Une fonction $f$ est surjective si et seulement si elle n'appartient à aucun $A_i$.
Par inclusion-exclusion (forme complémentaire),
\begin{align*}
  S(n,k) &= k^n - \card*{\bigcup_{i=1}^{k} A_i} \\
  &= \sum_{j=0}^{k} (-1)^{j} \binom{k}{j} (k-j)^n,
\end{align*}
car l'intersection $A_{i_1} \cap \cdots \cap A_{i_j}$ consiste en les fonctions
évitant $j$ éléments fixés, donc à valeurs dans un ensemble de $k-j$ éléments,
et $\card{A_{i_1} \cap \cdots \cap A_{i_j}} = (k-j)^n$.
\end{proof}

\begin{exemple}
\label{ex:surjections-calcul}
Le nombre de surjections de $\set{1,2,3,4}$ vers $\set{a,b,c}$ est
\[
  S(4,3) = \binom{3}{0} 3^4 - \binom{3}{1} 2^4 + \binom{3}{2} 1^4 - \binom{3}{3} 0^4
  = 81 - 48 + 3 - 0 = 36.
\]
\end{exemple}

\subsection{L'indicatrice d'Euler via inclusion-exclusion}
\index{Euler!indicatrice}
\index{fonction d'Euler|see{Euler, indicatrice}}

\begin{definition}[Indicatrice d'Euler]
\label{def:euler-totient}
Pour $n \geqslant 1$, l'indicatrice d'Euler\index{indicatrice d'Euler} $\varphi(n)$ est
le nombre d'entiers $k$ tels que $1 \leqslant k \leqslant n$ et $\pgcd(k,n) = 1$.
\end{definition}

\begin{theoreme}[Formule d'Euler via inclusion-exclusion]
\label{thm:euler-ie}
Si $n = p_1^{a_1} p_2^{a_2} \cdots p_r^{a_r}$ est la factorisation en nombres
premiers de $n$, alors
\[
  \varphi(n) = n \prod_{i=1}^{r} \left(1 - \frac{1}{p_i}\right).
\]
\end{theoreme}

\begin{proof}
Soit $U = \set{1, 2, \dots, n}$ et pour chaque $i$, posons
$A_i = \set{k \in U : p_i \mid k}$. Un entier $k \in U$ satisfait
$\pgcd(k,n) > 1$ si et seulement si $k \in A_i$ pour au moins un $i$.
Donc $\varphi(n) = \card{U \setminus \bigcup_{i=1}^r A_i}$.

On a $\card{A_i} = n/p_i$, et plus généralement, pour
$1 \leqslant i_1 < \cdots < i_s \leqslant r$ :
\[
  \card{A_{i_1} \cap \cdots \cap A_{i_s}} = \frac{n}{p_{i_1} \cdots p_{i_s}}.
\]
Par inclusion-exclusion,
\begin{align*}
  \varphi(n) &= n - \sum_i \frac{n}{p_i} + \sum_{i<j} \frac{n}{p_i p_j}
    - \cdots + (-1)^r \frac{n}{p_1 \cdots p_r} \\
  &= n \prod_{i=1}^{r} \left(1 - \frac{1}{p_i}\right).
  \qedhere
\end{align*}
\end{proof}

\begin{exemple}
\label{ex:euler-360}
Calculons $\varphi(360)$. Puisque $360 = 2^3 \cdot 3^2 \cdot 5$,
\[
  \varphi(360) = 360 \cdot \left(1 - \frac{1}{2}\right)
    \left(1 - \frac{1}{3}\right)\left(1 - \frac{1}{5}\right)
  = 360 \cdot \frac{1}{2} \cdot \frac{2}{3} \cdot \frac{4}{5} = 96.
\]
\end{exemple}

\subsection{Les dérangements}
\label{sec:derangements}
\index{dérangement}

\begin{definition}[Dérangement]
\label{def:derangement}
Un \emph{dérangement} de $\set{1, 2, \dots, n}$ est une permutation $\sigma$
telle que $\sigma(i) \neq i$ pour tout $i$. On note $D_n$\index{$D_n$}
le nombre de dérangements de $\set{1, \dots, n}$.
\end{definition}

\begin{theoreme}[Formule des dérangements]
\label{thm:derangements}
\index{dérangement!formule}
Pour tout $n \geqslant 0$,
\[
  D_n = n! \sum_{k=0}^{n} \frac{(-1)^k}{k!}.
\]
En particulier, $D_n$ est l'entier le plus proche de $n!/e$.
\end{theoreme}

\begin{proof}
Soit $A_i = \set{\sigma \in \mathfrak{S}_n : \sigma(i) = i}$ l'ensemble
des permutations fixant $i$. Un dérangement est une permutation
n'appartenant à aucun $A_i$, donc
\[
  D_n = n! - \card*{\bigcup_{i=1}^n A_i}.
\]
Pour tout sous-ensemble $\set{i_1, \dots, i_k}$,
$\card{A_{i_1} \cap \cdots \cap A_{i_k}} = (n-k)!$
(les $k$ éléments sont fixés, les $n-k$ restants permutent librement).
Il y a $\binom{n}{k}$ tels sous-ensembles, d'où par inclusion-exclusion :
\begin{align*}
  D_n &= \sum_{k=0}^{n} (-1)^k \binom{n}{k}(n-k)!
       = \sum_{k=0}^{n} (-1)^k \frac{n!}{k!}
       = n! \sum_{k=0}^{n} \frac{(-1)^k}{k!}. \qedhere
\end{align*}
\end{proof}

\begin{exemple}[Premières valeurs]
\label{ex:derangements-valeurs}
\begin{center}
\begin{tabular}{c|cccccccc}
\toprule
$n$ & $0$ & $1$ & $2$ & $3$ & $4$ & $5$ & $6$ & $7$ \\
\midrule
$D_n$ & $1$ & $0$ & $1$ & $2$ & $9$ & $44$ & $265$ & $1\,854$ \\
$D_n/n!$ & $1$ & $0$ & $0{,}5$ & $0{,}333$ & $0{,}375$ & $0{,}367$ & $0{,}368$ & $0{,}368$ \\
\bottomrule
\end{tabular}
\end{center}
On observe que $D_n / n! \to 1/e \approx 0{,}3679$.
\end{exemple}

\begin{corollaire}[Récurrence des dérangements]
\label{cor:derangement-rec}
Pour $n \geqslant 2$, on a $D_n = (n-1)(D_{n-1} + D_{n-2})$.
\end{corollaire}

\begin{proof}
Soit $\sigma$ un dérangement de $\set{1, \dots, n}$. Posons $\sigma(1) = j$
avec $j \neq 1$, ce qui donne $n-1$ choix pour $j$.
\begin{itemize}[nosep]
  \item Si $\sigma(j) = 1$ : les éléments restants forment un dérangement
    de $\set{2, \dots, n} \setminus \set{j}$, d'où $D_{n-2}$ possibilités.
  \item Si $\sigma(j) \neq 1$ : on effectue la substitution $\sigma'(i) = \sigma(i)$
    pour $i \neq 1$ et $\sigma'(j) = \sigma(j)$ sauf que l'interdiction
    $\sigma(j) \neq 1$ remplace $\sigma(j) \neq j$, donnant $D_{n-1}$ possibilités.
\end{itemize}
D'où $D_n = (n-1)(D_{n-1} + D_{n-2})$.
\end{proof}

\subsection{Le problème des ménages (aperçu)}
\label{sec:menage}
\index{problème des ménages}

\begin{remarque}[Problème des ménages]
\label{rem:menage}
Le \emph{problème des ménages} demande de compter le nombre de façons
d'asseoir $n$ couples autour d'une table ronde de sorte qu'aucun époux
ne soit à côté de son conjoint et que les places alternent hommes et femmes.
La réponse, notée $M_n$, est donnée par la formule de Touchard\index{Touchard, formule de} :
\[
  M_n = n! \sum_{k=0}^{n} (-1)^k \frac{2n}{2n-k} \binom{2n-k}{k} (n-k)!
\]
pour $n \geqslant 3$. La démonstration, qui repose sur une application
raffinée de l'inclusion-exclusion, dépasse le cadre de ce cours.
\end{remarque}

\section{Nombres de Stirling de seconde espèce}
\label{sec:stirling}
\index{Stirling!nombres de seconde espèce}

\begin{definition}[Nombre de Stirling de seconde espèce]
\label{def:stirling-second}
Le \emph{nombre de Stirling de seconde espèce}\index{$S(n,k)$|see{Stirling}} $\genfrac{\{}{\}}{0pt}{}{n}{k}$
est le nombre de partitions d'un ensemble à $n$ éléments en exactement $k$
blocs non vides. On le note aussi $S(n,k)$.
\end{definition}

\begin{theoreme}[Récurrence de Stirling]
\label{thm:stirling-rec}
\index{Stirling!récurrence}
Pour $n \geqslant 1$ et $1 \leqslant k \leqslant n$,
\[
  \genfrac{\{}{\}}{0pt}{}{n}{k}
    = k \genfrac{\{}{\}}{0pt}{}{n-1}{k}
    + \genfrac{\{}{\}}{0pt}{}{n-1}{k-1},
\]
avec les conditions initiales $\genfrac{\{}{\}}{0pt}{}{0}{0} = 1$ et
$\genfrac{\{}{\}}{0pt}{}{n}{0} = 0$ pour $n \geqslant 1$.
\end{theoreme}

\begin{proof}
Considérons l'élément $n$ dans une partition de $\set{1,\dots,n}$ en $k$ blocs.
\begin{itemize}[nosep]
  \item Si $\set{n}$ est un bloc à lui seul : les $n-1$ éléments restants forment
    une partition en $k-1$ blocs, d'où $\genfrac{\{}{\}}{0pt}{}{n-1}{k-1}$ possibilités.
  \item Si $n$ est dans un bloc contenant d'autres éléments : on partitionne
    d'abord $\set{1,\dots,n-1}$ en $k$ blocs ($\genfrac{\{}{\}}{0pt}{}{n-1}{k}$ façons),
    puis on insère $n$ dans l'un des $k$ blocs. \qedhere
\end{itemize}
\end{proof}

% --- Triangle de Stirling ---
\begin{figure}[ht]
\centering
\begin{tikzpicture}[scale=0.75, every node/.style={minimum size=8mm}]
  % En-tête
  \node at (-1.5,0) {\footnotesize $n \backslash k$};
  \foreach \k in {0,...,6} {
    \node at (1.3*\k+0.5, 0) {\footnotesize $\k$};
  }
  % Lignes
  \def\data{%
    {0/{1,0,0,0,0,0,0}},
    {1/{0,1,0,0,0,0,0}},
    {2/{0,1,1,0,0,0,0}},
    {3/{0,1,3,1,0,0,0}},
    {4/{0,1,7,6,1,0,0}},
    {5/{0,1,15,25,10,1,0}},
    {6/{0,1,31,90,65,15,1}}%
  }
  \foreach \n/\row [count=\i from 0] in \data {
    \node at (-1.5, -0.8*\i-0.8) {\footnotesize $\n$};
    \foreach \val [count=\j from 0] in \row {
      \ifnum\val>0
        \node at (1.3*\j+0.5, -0.8*\i-0.8) {\footnotesize $\val$};
      \else
        \node[gray!50] at (1.3*\j+0.5, -0.8*\i-0.8) {\footnotesize $0$};
      \fi
    }
  }
  % Cadre
  \draw[gray,thin] (-2.2,0.4) -- (9.0,0.4);
  \draw[gray,thin] (-0.6,0.4) -- (-0.6,-6.4);
\end{tikzpicture}
\caption{Triangle des nombres de Stirling de seconde espèce $\genfrac{\{}{\}}{0pt}{}{n}{k}$ pour $0 \leqslant n \leqslant 6$.}
\label{fig:triangle-stirling}
\end{figure}

\begin{theoreme}[Formule explicite de Stirling]
\label{thm:stirling-explicite}
Pour $n \geqslant k \geqslant 1$,
\[
  \genfrac{\{}{\}}{0pt}{}{n}{k}
    = \frac{1}{k!} \sum_{j=0}^{k} (-1)^{j} \binom{k}{j} (k-j)^n.
\]
\end{theoreme}

\begin{proof}
Le nombre de surjections de $\set{1,\dots,n}$ vers $\set{1,\dots,k}$ est,
d'une part, donné par le théorème~\ref{thm:surjections} :
$\sum_{j=0}^{k} (-1)^j \binom{k}{j}(k-j)^n$.
D'autre part, chaque partition en $k$ blocs donne lieu à $k!$ surjections
(en étiquetant les blocs). Donc le nombre de surjections est
$k! \genfrac{\{}{\}}{0pt}{}{n}{k}$, ce qui donne la formule.
\end{proof}

\begin{exemple}
\label{ex:stirling-calcul}
Calculons $\genfrac{\{}{\}}{0pt}{}{4}{2}$ :
\[
  \genfrac{\{}{\}}{0pt}{}{4}{2} = \frac{1}{2!}\bigl[\binom{2}{0}2^4 - \binom{2}{1}1^4 + \binom{2}{2}0^4\bigr]
  = \frac{1}{2}(16 - 2 + 0) = 7.
\]
Les $7$ partitions de $\set{1,2,3,4}$ en $2$ blocs sont :
$\set{\set{1},\set{2,3,4}}$, $\set{\set{2},\set{1,3,4}}$,
$\set{\set{3},\set{1,2,4}}$, $\set{\set{4},\set{1,2,3}}$,
$\set{\set{1,2},\set{3,4}}$, $\set{\set{1,3},\set{2,4}}$,
$\set{\set{1,4},\set{2,3}}$.
\end{exemple}

\section{Nombres de Bell}
\label{sec:bell}
\index{Bell, nombre de}

\begin{definition}[Nombre de Bell]
\label{def:bell}
Le \emph{nombre de Bell}\index{$B_n$} $B_n$ est le nombre total de partitions
d'un ensemble à $n$ éléments :
\[
  B_n = \sum_{k=0}^{n} \genfrac{\{}{\}}{0pt}{}{n}{k}.
\]
\end{definition}

\begin{proposition}[Récurrence de Bell]
\label{prop:bell-rec}
Pour $n \geqslant 1$,
\[
  B_{n} = \sum_{k=0}^{n-1} \binom{n-1}{k} B_k.
\]
\end{proposition}

\begin{proof}
Le bloc contenant l'élément $n$ contient $n - 1 - k$ autres éléments
(choisis parmi les $n-1$ restants de $\binom{n-1}{n-1-k} = \binom{n-1}{k}$ façons),
et les $k$ éléments restants forment une partition arbitraire ($B_k$ possibilités).
En sommant sur $k = 0, \dots, n-1$, on obtient la récurrence.
\end{proof}

\begin{exemple}[Premières valeurs]
\label{ex:bell-valeurs}
\[
  B_0 = 1,\quad B_1 = 1,\quad B_2 = 2,\quad B_3 = 5,\quad
  B_4 = 15,\quad B_5 = 52,\quad B_6 = 203.
\]
\end{exemple}

\begin{remarque}[Formule de Dobinski]
\label{rem:dobinski}
\index{Dobinski, formule de}
Les nombres de Bell admettent la représentation en série :
\[
  B_n = \frac{1}{e} \sum_{k=0}^{\infty} \frac{k^n}{k!}.
\]
Cette formule, due à Dobinski (1877), sera retrouvée via les séries génératrices
exponentielles au chapitre~\ref{ch:generatrices}.
\end{remarque}

\section{Partitions d'entiers (introduction)}
\label{sec:partitions-entiers}
\index{partition!d'un entier}

\begin{definition}[Partition d'un entier]
\label{def:partition-entier}
Une \emph{partition} d'un entier positif $n$ est une écriture
$n = \lambda_1 + \lambda_2 + \cdots + \lambda_r$ avec
$\lambda_1 \geqslant \lambda_2 \geqslant \cdots \geqslant \lambda_r \geqslant 1$.
On note $p(n)$\index{$p(n)$} le nombre de partitions de $n$.
\end{definition}

\begin{exemple}
\label{ex:partitions-5}
Les partitions de $5$ sont : $5$, $4+1$, $3+2$, $3+1+1$, $2+2+1$,
$2+1+1+1$, $1+1+1+1+1$. Donc $p(5) = 7$.
\end{exemple}

\begin{theoreme}[Fonction génératrice des partitions]
\label{thm:partition-gen}
La série génératrice des $p(n)$ est
\[
  \sum_{n=0}^{\infty} p(n) x^n = \prod_{k=1}^{\infty} \frac{1}{1 - x^k}.
\]
\end{theoreme}

\begin{proof}[Idée de preuve]
Chaque facteur $\frac{1}{1-x^k} = 1 + x^k + x^{2k} + \cdots$ encode
le nombre de fois que la part $k$ apparaît dans la partition. En multipliant
sur tous les $k \geqslant 1$, on obtient exactement la somme sur toutes les partitions.
\end{proof}

\begin{theoreme}[Théorème d'Euler sur les partitions]
\label{thm:euler-partitions}
\index{Euler!théorème sur les partitions}
Le nombre de partitions de $n$ en parts impaires est égal au nombre de partitions
de $n$ en parts distinctes.
\end{theoreme}

\begin{proof}
Les partitions en parts distinctes ont pour génératrice
$\prod_{k=1}^\infty (1 + x^k)$. On a
\[
  \prod_{k=1}^{\infty}(1+x^k)
  = \prod_{k=1}^{\infty} \frac{1 - x^{2k}}{1-x^k}
  = \frac{\prod_{k=1}^{\infty}(1-x^{2k})}{\prod_{k=1}^{\infty}(1-x^k)}
  = \prod_{\substack{k=1 \\ k \text{ impair}}}^{\infty} \frac{1}{1-x^k},
\]
qui est la génératrice des partitions en parts impaires.
\end{proof}

\section{Exercices}
\label{sec:ie-exercices}

\begin{exercice}[Crible d'Ératosthène]
\label{exo:crible}
\index{crible d'Ératosthène}
Utilisez le principe d'inclusion-exclusion pour déterminer le nombre
de nombres premiers dans $\set{1, 2, \dots, 120}$.
\emph{Indication} : il suffit de cribler par $2$, $3$, $5$ et $7$.
\end{exercice}

\begin{exercice}[Fonctions à image prescrite]
\label{exo:fonctions-image}
Combien de fonctions $f : \set{1,\dots,6} \to \set{a,b,c,d}$ ont pour image
exactement $\set{a,b,c}$ ?
\end{exercice}

\begin{exercice}[Dérangements et probabilités]
\label{exo:derangement-proba}
Un professeur rend au hasard $n$ copies à $n$ étudiants. Quelle est la probabilité
qu'aucun étudiant ne reçoive sa propre copie ? Calculer la limite quand $n \to \infty$.
\end{exercice}

\begin{exercice}[Stirling et identités]
\label{exo:stirling-id}
Montrez que pour tout $n \geqslant 1$ :
\[
  \sum_{k=0}^{n} \genfrac{\{}{\}}{0pt}{}{n}{k} x^{\underline{k}} = x^n,
\]
où $x^{\underline{k}} = x(x-1)\cdots(x-k+1)$ désigne la puissance descendante.
\end{exercice}

\begin{exercice}[Surjections et nombres de Stirling]
\label{exo:surj-stirling}
En utilisant la relation entre surjections et nombres de Stirling, montrez que
$\genfrac{\{}{\}}{0pt}{}{n}{n-1} = \binom{n}{2}$ pour tout $n \geqslant 2$.
Interprétez combinatoirement ce résultat.
\end{exercice}

\begin{exercice}[Triangle de Bell]
\label{exo:triangle-bell}
Le \emph{triangle de Bell}\index{Bell!triangle de} est construit comme suit :
la première ligne contient $B_0 = 1$, et chaque ligne suivante commence par
le dernier élément de la ligne précédente, les éléments suivants étant la somme
de l'élément à gauche et de l'élément au-dessus. Construisez les $6$ premières
lignes et vérifiez que les premiers éléments de chaque ligne sont les nombres de Bell.
\end{exercice}

\begin{exercice}[Indicatrice d'Euler]
\label{exo:euler-totient}
\begin{enumerate}[label=\alph*)]
  \item Calculez $\varphi(1000)$ et $\varphi(2025)$.
  \item Montrez que $\sum_{d \mid n} \varphi(d) = n$ en utilisant
    l'inclusion-exclusion.
\end{enumerate}
\end{exercice}

\begin{exercice}[Partitions en au plus $k$ parts]
\label{exo:partitions-k}
Montrez, par un argument de génératrices, que le nombre de partitions de $n$
en au plus $k$ parts est égal au nombre de partitions de $n$ dont chaque
part est $\leqslant k$.
\end{exercice}

\section*{Résumé du chapitre}
\addcontentsline{toc}{section}{Résumé du chapitre}

\begin{itemize}[nosep]
  \item Le \textbf{principe d'inclusion-exclusion} donne
    $\card{\bigcup A_i}$ en termes des cardinaux des intersections, avec
    des signes alternés.
  \item Applications principales : \textbf{surjections}
    ($\sum (-1)^j \binom{k}{j}(k-j)^n$), \textbf{indicatrice d'Euler}
    ($\varphi(n) = n\prod(1-1/p_i)$), \textbf{dérangements}
    ($D_n = n!\sum (-1)^k/k!$).
  \item Les \textbf{nombres de Stirling de seconde espèce}
    $\genfrac{\{}{\}}{0pt}{}{n}{k}$ comptent les partitions d'ensembles
    en $k$ blocs et sont liés aux surjections par
    $\genfrac{\{}{\}}{0pt}{}{n}{k} = S(n,k)/k!$.
  \item Les \textbf{nombres de Bell} $B_n = \sum_k \genfrac{\{}{\}}{0pt}{}{n}{k}$
    comptent toutes les partitions d'un ensemble.
  \item Les \textbf{partitions d'entiers} $p(n)$ ont pour génératrice
    $\prod_k (1-x^k)^{-1}$.
\end{itemize}


%%%%%%%%%%%%%%%%%%%%%%%%%%%%%%%%%%%%%%%%%%%%%%%%%%%%%%%%%%%%%%%%%%%%
%  CHAPITRE 6 — Génératrices Ordinaires et Exponentielles
%%%%%%%%%%%%%%%%%%%%%%%%%%%%%%%%%%%%%%%%%%%%%%%%%%%%%%%%%%%%%%%%%%%%
\chapter{Séries Génératrices Ordinaires et Exponentielles}
\label{ch:generatrices}
\index{série génératrice}

Les séries génératrices constituent l'un des outils les plus élégants et
les plus puissants de la combinatoire. Elles transforment des problèmes
de dénombrement en problèmes d'algèbre et d'analyse, permettant de résoudre
des récurrences, d'obtenir des formules explicites et de démontrer des identités.

\section{Séries génératrices ordinaires}
\label{sec:ogf}
\index{série génératrice!ordinaire}

\begin{definition}[Série génératrice ordinaire]
\label{def:ogf}
Soit $(a_n)_{n \geqslant 0}$ une suite de nombres. La \emph{série génératrice
ordinaire} (SGO) de cette suite est la série formelle
\[
  A(x) = \sum_{n=0}^{\infty} a_n x^n.
\]
\end{definition}

\begin{notation}
\label{not:coefficient}
On note $[x^n] A(x) = a_n$ le coefficient de $x^n$ dans $A(x)$.
\end{notation}

\begin{exemple}[Exemples fondamentaux]
\label{ex:ogf-fondamentaux}
\begin{enumerate}[label=\roman*)]
  \item La suite constante $a_n = 1$ a pour SGO
    $\displaystyle\frac{1}{1-x} = \sum_{n=0}^{\infty} x^n$.
  \item La suite $a_n = r^n$ a pour SGO
    $\displaystyle\frac{1}{1-rx} = \sum_{n=0}^{\infty} r^n x^n$.
  \item La suite $a_n = \binom{n+k-1}{k-1}$ a pour SGO
    $\displaystyle\frac{1}{(1-x)^k}$.
  \item La suite $a_n = \binom{m}{n}$ (avec $m$ fixé) a pour SGO $(1+x)^m$.
\end{enumerate}
\end{exemple}

\subsection{Opérations sur les SGO}
\label{sec:ogf-operations}

\begin{proposition}[Addition et multiplication scalaire]
\label{prop:ogf-lineaire}
Si $A(x) = \sum a_n x^n$ et $B(x) = \sum b_n x^n$, alors
\[
  \alpha A(x) + \beta B(x) = \sum_{n=0}^{\infty} (\alpha a_n + \beta b_n) x^n.
\]
\end{proposition}

\begin{theoreme}[Produit de Cauchy (convolution)]
\label{thm:ogf-produit}
\index{convolution!produit de Cauchy}
Si $A(x) = \sum a_n x^n$ et $B(x) = \sum b_n x^n$, alors
\[
  A(x) \cdot B(x) = \sum_{n=0}^{\infty} c_n x^n,
  \quad \text{où} \quad
  c_n = \sum_{k=0}^{n} a_k b_{n-k}.
\]
\end{theoreme}

\begin{proof}
C'est la formule standard du produit de séries formelles, qui résulte de la
distributivité : le terme $a_k x^k \cdot b_{n-k} x^{n-k} = a_k b_{n-k} x^n$
contribue au coefficient de $x^n$.
\end{proof}

\begin{proposition}[Décalage]
\label{prop:ogf-decalage}
\index{série génératrice!décalage}
Si $A(x) = \sum_{n=0}^\infty a_n x^n$, alors :
\begin{enumerate}[label=\roman*)]
  \item $\displaystyle\frac{A(x) - a_0}{x} = \sum_{n=0}^{\infty} a_{n+1} x^n$
    \quad (décalage à gauche).
  \item $x A(x) = \sum_{n=1}^{\infty} a_{n-1} x^n$
    \quad (décalage à droite).
\end{enumerate}
\end{proposition}

\begin{proposition}[Dérivation]
\label{prop:ogf-derivation}
\index{série génératrice!dérivation}
Si $A(x) = \sum_{n=0}^\infty a_n x^n$, alors
\[
  A'(x) = \sum_{n=0}^{\infty} (n+1) a_{n+1} x^n,
  \qquad
  x A'(x) = \sum_{n=0}^{\infty} n a_n x^n.
\]
Plus généralement, l'opérateur $x\frac{d}{dx}$ revient à multiplier
$a_n$ par $n$.
\end{proposition}

\begin{exemple}[Application de la dérivation]
\label{ex:ogf-derivation}
Calculons $\sum_{n=0}^\infty n x^n$. On part de $\frac{1}{1-x} = \sum x^n$
et on applique $x\frac{d}{dx}$ :
\[
  x \cdot \frac{d}{dx}\frac{1}{1-x} = \frac{x}{(1-x)^2} = \sum_{n=0}^{\infty} n x^n.
\]
En posant $x = 1/2$, on retrouve $\sum_{n=0}^\infty n/2^n = 2$.
\end{exemple}

\subsection{Résolution de récurrences par SGO}
\label{sec:ogf-recurrences}

\begin{exemple}[Suite de Fibonacci via SGO]
\label{ex:fibonacci-ogf}
\index{Fibonacci!série génératrice}
Résolvons la récurrence $F_0 = 0$, $F_1 = 1$, $F_{n} = F_{n-1} + F_{n-2}$
pour $n \geqslant 2$. Posons $F(x) = \sum_{n=0}^\infty F_n x^n$.
La récurrence donne
\[
  F(x) - F_0 - F_1 x = x(F(x) - F_0) + x^2 F(x),
\]
soit $F(x) - x = xF(x) + x^2 F(x)$, d'où
\[
  F(x) = \frac{x}{1 - x - x^2}.
\]

En décomposant en éléments simples avec les racines
$\alpha = \frac{1+\sqrt{5}}{2}$ et $\beta = \frac{1-\sqrt{5}}{2}$
de $1 - x - x^2 = -(x - 1/\alpha)(x - 1/\beta)$, on obtient
\[
  F(x) = \frac{1}{\sqrt{5}}\left(\frac{1}{1-\alpha x} - \frac{1}{1-\beta x}\right),
\]
ce qui donne la formule de Binet\index{Binet, formule de} :
\[
  F_n = \frac{\alpha^n - \beta^n}{\sqrt{5}}
      = \frac{1}{\sqrt{5}}\left[\left(\frac{1+\sqrt{5}}{2}\right)^{\!n}
        - \left(\frac{1-\sqrt{5}}{2}\right)^{\!n}\right].
\]
\end{exemple}

\begin{exemple}[Récurrence linéaire d'ordre 2]
\label{ex:ogf-rec-ordre2}
Résolvons $a_0 = 1$, $a_1 = 3$, $a_n = 4a_{n-1} - 4a_{n-2}$ pour $n \geqslant 2$.
Posons $A(x) = \sum a_n x^n$. La récurrence donne
\[
  A(x) - 1 - 3x = 4x(A(x) - 1) - 4x^2 A(x),
\]
d'où $A(x)(1 - 4x + 4x^2) = 1 - x$, soit
\[
  A(x) = \frac{1-x}{(1-2x)^2}.
\]
En décomposant : $A(x) = \frac{1}{1-2x} + \frac{1}{(1-2x)^2}$ (après calcul).
Or $[x^n]\frac{1}{(1-2x)^2} = (n+1)2^n$, donc $a_n = 2^n + (n+1)2^n = (n+2)2^n$.
\end{exemple}

\subsection{Les nombres de Catalan}
\label{sec:catalan}
\index{Catalan, nombre de}

\begin{definition}[Nombre de Catalan]
\label{def:catalan}
Le $n$-ième \emph{nombre de Catalan}\index{$C_n$} est
\[
  C_n = \frac{1}{n+1}\binom{2n}{n}, \qquad n \geqslant 0.
\]
\end{definition}

\begin{theoreme}[SGO des nombres de Catalan]
\label{thm:catalan-ogf}
\index{Catalan!série génératrice}
La série génératrice ordinaire $C(x) = \sum_{n=0}^\infty C_n x^n$ satisfait
l'équation fonctionnelle
\[
  C(x) = 1 + x C(x)^2,
\]
et on a
\[
  C(x) = \frac{1 - \sqrt{1 - 4x}}{2x}.
\]
\end{theoreme}

\begin{proof}
Les nombres de Catalan satisfont la récurrence
$C_0 = 1$ et $C_{n+1} = \sum_{k=0}^{n} C_k C_{n-k}$ pour $n \geqslant 0$
(décomposition récursive des chemins de Dyck, par exemple). En termes de
séries génératrices, cela se traduit par
\[
  \frac{C(x) - 1}{x} = C(x)^2,
  \quad\text{soit}\quad
  C(x) = 1 + x C(x)^2.
\]
C'est une équation du second degré en $C(x)$ :
$x C(x)^2 - C(x) + 1 = 0$, de discriminant $\Delta = 1 - 4x$, d'où
\[
  C(x) = \frac{1 \pm \sqrt{1-4x}}{2x}.
\]
Comme $C(0) = 1$, on choisit le signe $-$ (le signe $+$ donnerait $C(x) \to \infty$
quand $x \to 0$).

Vérifions la formule du coefficient. En utilisant le binôme généralisé\index{binôme généralisé},
\[
  \sqrt{1-4x} = \sum_{n=0}^{\infty} \binom{1/2}{n}(-4x)^n,
\]
on obtient, après simplification,
\[
  [x^n] C(x) = \frac{1}{n+1}\binom{2n}{n} = C_n. \qedhere
\]
\end{proof}

\begin{exemple}[Premières valeurs de Catalan]
\label{ex:catalan-valeurs}
\[
  C_0 = 1,\; C_1 = 1,\; C_2 = 2,\; C_3 = 5,\; C_4 = 14,\;
  C_5 = 42,\; C_6 = 132,\; C_7 = 429.
\]
\end{exemple}

\begin{proposition}[Interprétations combinatoires de $C_n$]
\label{prop:catalan-interpretations}
\index{Catalan!interprétations}
Le nombre $C_n$ compte, entre autres :
\begin{enumerate}[label=\roman*)]
  \item le nombre de chemins de Dyck\index{chemin de Dyck} de longueur $2n$
    (chemins de $(0,0)$ à $(2n,0)$ avec des pas $+1$ et $-1$, ne passant
    jamais sous l'axe des abscisses) ;
  \item le nombre de parenthésages corrects de $n$ paires de parenthèses ;
  \item le nombre d'arbres binaires complets à $n$ n\oe uds internes ;
  \item le nombre de triangulations d'un polygone convexe à $n+2$ côtés.
\end{enumerate}
\end{proposition}

% --- Chemins de Dyck pour C_3 = 5 ---
\begin{figure}[ht]
\centering
\begin{tikzpicture}[scale=0.5]
  % Chemin 1 : UUUDDD
  \begin{scope}[shift={(0,0)}]
    \draw[gray!30] (0,0) grid (6,3);
    \draw[thick,blue!70!black] (0,0)--(1,1)--(2,2)--(3,3)--(4,2)--(5,1)--(6,0);
    \node[below] at (3,-0.5) {\tiny UUUDDD};
  \end{scope}
  % Chemin 2 : UUDUUD -> non, UUDUDD
  \begin{scope}[shift={(7.5,0)}]
    \draw[gray!30] (0,0) grid (6,3);
    \draw[thick,blue!70!black] (0,0)--(1,1)--(2,2)--(3,1)--(4,2)--(5,1)--(6,0);
    \node[below] at (3,-0.5) {\tiny UUDUDD};
  \end{scope}
  % Chemin 3 : UUDDUD
  \begin{scope}[shift={(15,0)}]
    \draw[gray!30] (0,0) grid (6,3);
    \draw[thick,blue!70!black] (0,0)--(1,1)--(2,2)--(3,1)--(4,0)--(5,1)--(6,0);
    \node[below] at (3,-0.5) {\tiny UUDDUD};
  \end{scope}
  % Chemin 4 : UDUDUD
  \begin{scope}[shift={(3.5,-5.5)}]
    \draw[gray!30] (0,0) grid (6,3);
    \draw[thick,blue!70!black] (0,0)--(1,1)--(2,0)--(3,1)--(4,0)--(5,1)--(6,0);
    \node[below] at (3,-0.5) {\tiny UDUDUD};
  \end{scope}
  % Chemin 5 : UDUUDD
  \begin{scope}[shift={(11.5,-5.5)}]
    \draw[gray!30] (0,0) grid (6,3);
    \draw[thick,blue!70!black] (0,0)--(1,1)--(2,0)--(3,1)--(4,2)--(5,1)--(6,0);
    \node[below] at (3,-0.5) {\tiny UDUUDD};
  \end{scope}
\end{tikzpicture}
\caption{Les $C_3 = 5$ chemins de Dyck de longueur $6$. U = montée, D = descente.}
\label{fig:dyck-paths}
\end{figure}

% --- Parenthésages pour C_3 = 5 ---
\begin{figure}[ht]
\centering
\begin{tikzpicture}
  \node at (0,0) {\texttt{((()))}};
  \node at (3,0) {\texttt{(()())}};
  \node at (6,0) {\texttt{(())()}};
  \node at (9,0) {\texttt{()(())}};
  \node at (12,0) {\texttt{()()()}};
\end{tikzpicture}
\caption{Les $C_3 = 5$ parenthésages corrects de $3$ paires de parenthèses.}
\label{fig:parenthesages}
\end{figure}

\section{Séries génératrices exponentielles}
\label{sec:egf}
\index{série génératrice!exponentielle}

\subsection{Définition et motivation}

\begin{definition}[Série génératrice exponentielle]
\label{def:egf}
Soit $(a_n)_{n \geqslant 0}$ une suite. La \emph{série génératrice
exponentielle} (SGE) de cette suite est
\[
  \hat{A}(x) = \sum_{n=0}^{\infty} a_n \frac{x^n}{n!}.
\]
\end{definition}

\begin{remarque}
\label{rem:egf-motivation}
La SGE est particulièrement adaptée au dénombrement de structures
\emph{étiquetées} (permutations, graphes étiquetés, etc.), car le facteur
$1/n!$ dans la série compense le nombre de réétiquetages.
\end{remarque}

\begin{exemple}[SGE fondamentales]
\label{ex:egf-fondamentales}
\begin{enumerate}[label=\roman*)]
  \item La suite $a_n = 1$ a pour SGE $e^x = \sum \frac{x^n}{n!}$.
  \item La suite $a_n = n!$ (nombre de permutations) a pour SGE
    $\frac{1}{1-x} = \sum x^n$.
  \item La suite $a_n = r^n$ a pour SGE $e^{rx}$.
\end{enumerate}
\end{exemple}

\begin{theoreme}[Produit de SGE et structures étiquetées]
\label{thm:egf-produit}
\index{série génératrice!exponentielle!produit}
Si $\hat{A}(x) = \sum a_n \frac{x^n}{n!}$ et $\hat{B}(x) = \sum b_n \frac{x^n}{n!}$,
alors
\[
  \hat{A}(x) \cdot \hat{B}(x) = \sum_{n=0}^{\infty} c_n \frac{x^n}{n!},
  \quad \text{où} \quad
  c_n = \sum_{k=0}^{n} \binom{n}{k} a_k b_{n-k}.
\]
\end{theoreme}

\begin{proof}
On a
\[
  \hat{A}(x) \cdot \hat{B}(x)
    = \sum_{n=0}^{\infty} \left(\sum_{k=0}^{n} \frac{a_k}{k!} \cdot \frac{b_{n-k}}{(n-k)!}\right) x^n
    = \sum_{n=0}^{\infty} \frac{1}{n!} \left(\sum_{k=0}^{n} \binom{n}{k} a_k b_{n-k}\right) \frac{x^n}{1}.
\]
En identifiant, on obtient $c_n = \sum_k \binom{n}{k} a_k b_{n-k}$, qui est
la \emph{convolution binomiale}\index{convolution!binomiale}.
\end{proof}

\begin{remarque}[Interprétation combinatoire]
\label{rem:egf-produit-interpretation}
Si $a_n$ compte les structures de type $\mathcal{A}$ sur $n$ éléments et $b_n$
celles de type $\mathcal{B}$, alors $c_n$ compte les structures obtenues en
partitionnant $\set{1,\dots,n}$ en deux parts, plaçant une $\mathcal{A}$-structure
sur la première et une $\mathcal{B}$-structure sur la seconde. C'est le
\emph{produit étiqueté} de la combinatoire analytique.
\end{remarque}

\subsection{SGE des permutations}
\label{sec:egf-permutations}
\index{permutation!série génératrice}

\begin{exemple}[Permutations et SGE]
\label{ex:egf-permutations}
Puisque le nombre de permutations de $\set{1,\dots,n}$ est $n!$,
la SGE des permutations est
\[
  \sum_{n=0}^{\infty} n! \cdot \frac{x^n}{n!} = \sum_{n=0}^{\infty} x^n = \frac{1}{1-x}.
\]
\end{exemple}

\begin{exemple}[Involutions]
\label{ex:egf-involutions}
\index{involution}
Une \emph{involution} est une permutation $\sigma$ telle que $\sigma^2 = \mathrm{id}$.
Si $I_n$ désigne le nombre d'involutions de $\set{1,\dots,n}$, on a la récurrence
$I_n = I_{n-1} + (n-1)I_{n-2}$ (l'élément $n$ est soit un point fixe, soit
échangé avec l'un des $n-1$ autres éléments). La SGE satisfait
\[
  \hat{I}'(x) = (1+x)\hat{I}(x), \qquad \hat{I}(0) = 1,
\]
d'où $\hat{I}(x) = e^{x + x^2/2}$.
\end{exemple}

\subsection{SGE des dérangements}
\label{sec:egf-derangements}
\index{dérangement!série génératrice}

\begin{theoreme}[SGE des dérangements]
\label{thm:egf-derangements}
La SGE des dérangements est
\[
  \hat{D}(x) = \sum_{n=0}^{\infty} D_n \frac{x^n}{n!} = \frac{e^{-x}}{1-x}.
\]
\end{theoreme}

\begin{proof}
Toute permutation se décompose de manière unique en un ensemble de points fixes
et un dérangement du reste. Si $\hat{D}(x)$ est la SGE des dérangements,
la SGE des permutations (qui est $1/(1-x)$) s'écrit comme le produit
de la SGE de l'identité sur les points fixes ($e^x$, car un ensemble de $k$
points fixes se choisit de $\binom{n}{k}$ manières) et de $\hat{D}(x)$ :
\[
  \frac{1}{1-x} = e^x \cdot \hat{D}(x),
\]
d'où $\hat{D}(x) = \frac{e^{-x}}{1-x}$.

Vérifions : en développant,
\[
  \hat{D}(x) = \left(\sum_{k=0}^{\infty} \frac{(-x)^k}{k!}\right)
    \left(\sum_{m=0}^{\infty} x^m\right)
  = \sum_{n=0}^{\infty} \left(\sum_{k=0}^{n} \frac{(-1)^k}{k!}\right) x^n.
\]
Le coefficient de $x^n$ est $\sum_{k=0}^n \frac{(-1)^k}{k!} = D_n/n!$,
ce qui correspond bien à la formule du théorème~\ref{thm:derangements}.
\end{proof}

\subsection{SGE des nombres de Bell}
\label{sec:egf-bell}
\index{Bell, nombre de!série génératrice}

\begin{theoreme}[SGE des nombres de Bell]
\label{thm:egf-bell}
La SGE des nombres de Bell est
\[
  \sum_{n=0}^{\infty} B_n \frac{x^n}{n!} = e^{e^x - 1}.
\]
\end{theoreme}

\begin{proof}
Une partition d'un ensemble étiqueté peut être vue comme un ensemble
(non ordonné) de blocs non vides. En combinatoire analytique, la SGE
d'un ensemble de structures ayant pour SGE $\hat{B}(x)$ est $e^{\hat{B}(x)}$.
La SGE d'un bloc non vide (ensemble non vide) est $e^x - 1$.
Donc la SGE des partitions est $e^{e^x - 1}$.

Plus directement, on vérifie que la récurrence $B_{n+1} = \sum_{k=0}^n \binom{n}{k} B_k$
(proposition~\ref{prop:bell-rec}) est compatible avec cette SGE en dérivant :
\[
  \frac{d}{dx} e^{e^x - 1} = e^x \cdot e^{e^x - 1},
\]
ce qui signifie $B_{n+1}/n! = \sum_{k=0}^n B_k/(k! \cdot (n-k)!)$, soit
$B_{n+1} = \sum_{k=0}^n \binom{n}{k} B_k$.
\end{proof}

\begin{corollaire}[Formule de Dobinski]
\label{cor:dobinski}
\index{Dobinski, formule de}
Pour tout $n \geqslant 0$,
\[
  B_n = \frac{1}{e} \sum_{k=0}^{\infty} \frac{k^n}{k!}.
\]
\end{corollaire}

\begin{proof}
On développe $e^{e^x - 1} = e^{-1} e^{e^x} = e^{-1} \sum_{k=0}^\infty \frac{e^{kx}}{k!}
= e^{-1} \sum_{k=0}^\infty \frac{1}{k!} \sum_{n=0}^\infty \frac{(kx)^n}{n!}$.
En échangeant les sommations et en identifiant le coefficient de $x^n/n!$ :
$B_n = e^{-1} \sum_{k=0}^\infty k^n / k!$.
\end{proof}

\section{Séries formelles : structure d'anneau}
\label{sec:series-formelles}
\index{série formelle!anneau}

\begin{definition}[Anneau des séries formelles]
\label{def:anneau-series}
L'ensemble $\R[[x]]$\index{$\R[[x]]$} des séries formelles à coefficients réels,
muni de l'addition et de la multiplication (produit de Cauchy), forme un
anneau commutatif unitaire. Plus précisément, si
$A(x) = \sum a_n x^n$ et $B(x) = \sum b_n x^n$, alors :
\begin{itemize}[nosep]
  \item $(A+B)(x) = \sum (a_n + b_n) x^n$,
  \item $(A \cdot B)(x) = \sum_{n=0}^\infty \left(\sum_{k=0}^n a_k b_{n-k}\right) x^n$.
\end{itemize}
L'élément neutre pour l'addition est $0$ et pour la multiplication est $1$.
\end{definition}

\begin{proposition}[Inversibilité]
\label{prop:serie-inverse}
\index{série formelle!inverse}
Une série formelle $A(x) = \sum_{n=0}^\infty a_n x^n$ est inversible dans
$\R[[x]]$ si et seulement si $a_0 \neq 0$.
\end{proposition}

\begin{proof}
Si $A(x) B(x) = 1$, en identifiant le terme constant on obtient $a_0 b_0 = 1$,
donc $a_0 \neq 0$ est nécessaire. Réciproquement, si $a_0 \neq 0$, on peut
construire les coefficients $b_n$ de $B(x)$ par récurrence :
$b_0 = 1/a_0$ et pour $n \geqslant 1$,
\[
  b_n = -\frac{1}{a_0}\sum_{k=1}^{n} a_k b_{n-k}. \qedhere
\]
\end{proof}

\begin{proposition}[Composition]
\label{prop:serie-composition}
\index{série formelle!composition}
Si $A(x) = \sum a_n x^n$ et $B(x) = \sum_{n=1}^\infty b_n x^n$ (avec $b_0 = 0$),
alors la série composée $A(B(x))$ est bien définie dans $\R[[x]]$.
\end{proposition}

\begin{remarque}
\label{rem:anneau-integre}
L'anneau $\R[[x]]$ est intègre (produit de séries non nulles est non nul)
mais pas un corps. Cependant, le corps des fractions de $\R[[x]]$ est l'ensemble
$\R((x))$ des \emph{séries de Laurent formelles}\index{série de Laurent}.
\end{remarque}

\section{Exemples développés}
\label{sec:gen-exemples}

\subsection{Fibonacci via SGO : méthode complète}
\label{sec:fibonacci-complet}

\begin{exemple}[Fibonacci : résolution détaillée]
\label{ex:fibonacci-detail}
\index{Fibonacci!résolution complète par SGO}
Reprenons $F(x) = x/(1-x-x^2)$. Les racines de $1-x-x^2 = 0$ sont
$x = (-1 \pm \sqrt{5})/2$, donc $1/\alpha$ et $1/\beta$ où $\alpha = (1+\sqrt{5})/2$
(nombre d'or\index{nombre d'or}) et $\beta = (1-\sqrt{5})/2$.
Ainsi $1 - x - x^2 = -(x-1/\alpha)(x-1/\beta) = (1-\alpha x)(1-\beta x)/(\alpha\beta)$.
Comme $\alpha\beta = -1$, on a $1-x-x^2 = -(1-\alpha x)(1-\beta x) \cdot (-1) = (1-\alpha x)(1-\beta x)$.

Vérifions : $(1-\alpha x)(1-\beta x) = 1 - (\alpha+\beta)x + \alpha\beta x^2 = 1 - x - x^2$. \checkmark

Décomposition en éléments simples :
\[
  F(x) = \frac{x}{(1-\alpha x)(1-\beta x)}
    = \frac{A}{1-\alpha x} + \frac{B}{1-\beta x}.
\]
En posant $x = 1/\alpha$ : $A = \frac{1/\alpha}{1-\beta/\alpha} = \frac{1}{\alpha-\beta} = \frac{1}{\sqrt{5}}$.
De même, $B = -1/\sqrt{5}$. Donc
\[
  F_n = \frac{\alpha^n - \beta^n}{\sqrt{5}}.
\]
Puisque $\abs{\beta} < 1$, on a $F_n = \lfloor \alpha^n/\sqrt{5} + 1/2 \rfloor$
pour tout $n \geqslant 0$.
\end{exemple}

\subsection{Dérangements via SGE : vérification}
\label{sec:derangements-egf-verif}

\begin{exemple}[Dérangements : extraction de coefficients]
\label{ex:derangements-egf}
\index{dérangement!extraction de coefficients}
À partir de $\hat{D}(x) = e^{-x}/(1-x)$, extrayons $D_n$ directement.
On a
\[
  \hat{D}(x) = \left(\sum_{j=0}^{\infty} \frac{(-1)^j x^j}{j!}\right)
    \left(\sum_{m=0}^{\infty} x^m\right).
\]
Le coefficient de $x^n$ est
\[
  \sum_{j=0}^{n} \frac{(-1)^j}{j!} \cdot 1 = \sum_{j=0}^{n} \frac{(-1)^j}{j!} = \frac{D_n}{n!},
\]
d'où $D_n = n! \sum_{j=0}^n (-1)^j / j!$, conformément au théorème~\ref{thm:derangements}.

\medskip
\noindent\textbf{Vérification numérique.}
Pour $n = 4$ : $D_4 = 24(1 - 1 + 1/2 - 1/6 + 1/24) = 24 \cdot 3/8 = 9$.
Effectivement, les $9$ dérangements de $\set{1,2,3,4}$ sont :
$(2,1,4,3)$, $(2,3,4,1)$, $(2,4,1,3)$, $(3,1,4,2)$, $(3,4,1,2)$,
$(3,4,2,1)$, $(4,1,2,3)$, $(4,3,1,2)$, $(4,3,2,1)$.
\end{exemple}

\subsection{Nombres de Catalan : récurrence et formule close}

\begin{exemple}[Catalan via l'équation fonctionnelle]
\label{ex:catalan-rec}
\index{Catalan!récurrence}
De l'équation $C(x) = 1 + xC(x)^2$, on peut extraire la récurrence
directement. Le coefficient de $x^{n+1}$ dans $xC(x)^2$ est
$[x^n]C(x)^2 = \sum_{k=0}^n C_k C_{n-k}$, d'où
\[
  C_{n+1} = \sum_{k=0}^{n} C_k C_{n-k}.
\]
Calculons : $C_0 = 1$, $C_1 = C_0^2 = 1$, $C_2 = C_0 C_1 + C_1 C_0 = 2$,
$C_3 = C_0 C_2 + C_1^2 + C_2 C_0 = 2 + 1 + 2 = 5$, $C_4 = 2 \cdot 5 + 2 \cdot 2 + 5 \cdot 1 = 14$.
\end{exemple}

\subsection{Application : chemins dans un réseau}

\begin{exemple}[Chemins de réseau]
\label{ex:chemins-reseau}
\index{chemin de réseau}
Le nombre de chemins du point $(0,0)$ au point $(m,n)$ dans $\Z^2$,
n'utilisant que des pas vers la droite $(1,0)$ ou vers le haut $(0,1)$,
est $\binom{m+n}{m}$. La SGO en deux variables est
\[
  \sum_{m,n \geqslant 0} \binom{m+n}{m} x^m y^n = \frac{1}{1-x-y}.
\]
C'est un cas particulier du produit de SGO : le chemin est une suite de
$m+n$ pas, dont $m$ sont horizontaux.
\end{exemple}

\section{Techniques avancées}
\label{sec:gen-avancees}

\begin{proposition}[Extraction via le binôme généralisé]
\label{prop:binome-generalise}
\index{binôme généralisé}
Pour tout $\alpha \in \R$ et $\abs{x} < 1$,
\[
  (1+x)^{\alpha} = \sum_{n=0}^{\infty} \binom{\alpha}{n} x^n,
  \quad \text{où} \quad
  \binom{\alpha}{n} = \frac{\alpha(\alpha-1)\cdots(\alpha-n+1)}{n!}.
\]
\end{proposition}

\begin{exemple}[Coefficients centraux]
\label{ex:coefficients-centraux}
De $(1-4x)^{-1/2} = \sum_{n=0}^\infty \binom{2n}{n} x^n$, on déduit que
\[
  \binom{-1/2}{n}(-4)^n = \frac{(-1/2)(-3/2)\cdots(-(2n-1)/2)}{n!}(-4)^n = \binom{2n}{n}.
\]
\end{exemple}

\begin{proposition}[Méthode du serpent (snake oil)]
\label{prop:snake-oil}
\index{méthode du serpent}
Pour évaluer une somme $S = \sum_k f(n,k)$ faisant intervenir des coefficients
binomiaux, on forme la SGO $F(x) = \sum_n S_n x^n$, on intervertit les sommations,
et on simplifie le résultat à l'aide de séries connues.
\end{proposition}

\begin{exemple}[Identité de Vandermonde via SGO]
\label{ex:vandermonde-ogf}
\index{Vandermonde, identité de}
Montrons que $\sum_{k=0}^{r} \binom{m}{k}\binom{n}{r-k} = \binom{m+n}{r}$.
On compare les coefficients de $x^r$ dans $(1+x)^m (1+x)^n = (1+x)^{m+n}$.
Le membre gauche donne $\sum_k \binom{m}{k}\binom{n}{r-k}$ (produit de Cauchy)
et le membre droit donne $\binom{m+n}{r}$.
\end{exemple}

\section{Exercices}
\label{sec:gen-exercices}

\begin{exercice}[SGO et coefficients binomiaux]
\label{exo:ogf-binomial}
Montrez que la SGO de la suite $a_n = \binom{2n}{n}$ est
$\frac{1}{\sqrt{1-4x}}$.
\end{exercice}

\begin{exercice}[Récurrence de Catalan]
\label{exo:catalan-rec}
Démontrez que $C_n = \frac{4n-2}{n+1} C_{n-1}$ pour $n \geqslant 1$ à partir
de la formule $C_n = \binom{2n}{n}/(n+1)$.
\end{exercice}

\begin{exercice}[Suite de tribonacci]
\label{exo:tribonacci}
\index{tribonacci}
La suite de tribonacci est définie par $T_0 = 0$, $T_1 = 0$, $T_2 = 1$ et
$T_n = T_{n-1} + T_{n-2} + T_{n-3}$ pour $n \geqslant 3$. Déterminez sa SGO.
\end{exercice}

\begin{exercice}[Nombre de façons de rendre la monnaie]
\label{exo:monnaie}
\index{problème de la monnaie}
Avec des pièces de $1$, $2$ et $5$ centimes, de combien de façons peut-on
constituer une somme de $n$ centimes ? Exprimez la réponse comme coefficient
d'une SGO et calculez le résultat pour $n = 10$.
\end{exercice}

\begin{exercice}[SGE et permutations sans point fixe]
\label{exo:egf-sans-point-fixe}
En utilisant la SGE des dérangements, montrez que la probabilité qu'une
permutation aléatoire de $\set{1,\dots,n}$ soit un dérangement tend
vers $1/e$ quand $n \to \infty$.
\end{exercice}

\begin{exercice}[Nombres de Bell et congruence]
\label{exo:bell-mod}
\index{Bell, nombre de!congruence}
Montrez que $B_{p+1} \equiv B_1 + B_0 \pmod{p}$ pour tout nombre premier $p$,
c'est-à-dire $B_{p+1} \equiv 2 \pmod{p}$ (théorème de Touchard).
\emph{Indication} : utiliser la formule de Dobinski modulo $p$.
\end{exercice}

\begin{exercice}[Produit de SGE]
\label{exo:egf-produit}
Soient $a_n = 2^n$ et $b_n = n!$. Calculez la suite $c_n = \sum_{k=0}^n \binom{n}{k} a_k b_{n-k}$
à l'aide du produit des SGE et donnez les cinq premières valeurs.
\end{exercice}

\begin{exercice}[Arbres binaires et Catalan]
\label{exo:arbres-catalan}
\index{arbre binaire}
Soit $T_n$ le nombre d'arbres binaires à $n$ n\oe uds internes.
\begin{enumerate}[label=\alph*)]
  \item Montrez que $T_n = C_n$ en établissant que la SGO $T(x)$ satisfait
    $T(x) = 1 + xT(x)^2$.
  \item Déduisez-en que le nombre d'arbres binaires à $10$ n\oe uds internes
    est $C_{10} = 16\,796$.
\end{enumerate}
\end{exercice}

\begin{exercice}[Identité de Cauchy]
\label{exo:cauchy}
En utilisant les SGE, montrez l'identité
\[
  \sum_{k=0}^{n} \binom{n}{k}^2 = \binom{2n}{n}.
\]
\emph{Indication} : considérer $[x^n](1+x)^n (1+x)^n$.
\end{exercice}

\begin{exercice}[Série génératrice des partitions en parts distinctes]
\label{exo:parts-distinctes}
\index{partition!en parts distinctes}
Montrez que la SGO des partitions de $n$ en parts distinctes est
$\prod_{k=1}^{\infty}(1 + x^k)$ et calculez le nombre de telles partitions
pour $n = 8$.
\end{exercice}

\section*{Résumé du chapitre}
\addcontentsline{toc}{section}{Résumé du chapitre}

\begin{itemize}[nosep]
  \item La \textbf{série génératrice ordinaire} (SGO) $A(x) = \sum a_n x^n$
    encode une suite $(a_n)$. Le produit de SGO correspond à la
    \textbf{convolution} $c_n = \sum_k a_k b_{n-k}$.
  \item Les SGO permettent de résoudre des \textbf{récurrences linéaires} en
    les transformant en équations algébriques, puis en extrayant les coefficients
    par décomposition en éléments simples.
  \item Les \textbf{nombres de Catalan} $C_n = \binom{2n}{n}/(n+1)$ ont pour SGO
    $C(x) = (1-\sqrt{1-4x})/(2x)$, satisfaisant $C = 1 + xC^2$.
  \item La \textbf{série génératrice exponentielle} (SGE)
    $\hat{A}(x) = \sum a_n x^n/n!$ est adaptée aux structures étiquetées.
    Le produit de SGE correspond à la \textbf{convolution binomiale}.
  \item SGE remarquables : permutations $\to$ $1/(1-x)$, dérangements $\to$
    $e^{-x}/(1-x)$, nombres de Bell $\to$ $e^{e^x-1}$.
  \item L'anneau $\R[[x]]$ des séries formelles est commutatif, unitaire,
    intègre, et une série y est inversible ssi son terme constant est non nul.
\end{itemize}

% === FIN DU CHUNK ch5-6 ===
% === DÉBUT DU CHUNK ch7-8 ===

%%%%%%%%%%%%%%%%%%%%%%%%%%%%%%%%%%%%%%%%%%%%%%%%%%%%%%%%%%%%%%%%%%%%
%  CHAPITRE 7 — Récurrences et Équations de Récurrence
%%%%%%%%%%%%%%%%%%%%%%%%%%%%%%%%%%%%%%%%%%%%%%%%%%%%%%%%%%%%%%%%%%%%
\chapter{Récurrences et Équations de Récurrence}
\label{ch:recurrences}

\minitoc

\section*{Introduction}
\addcontentsline{toc}{section}{Introduction}

Les suites définies par récurrence sont omniprésentes en mathématiques discrètes
et en informatique. Qu'il s'agisse de compter des configurations combinatoires,
d'analyser la complexité d'un algorithme ou de modéliser un phénomène discret,
on est souvent amené à exprimer le $n$-ième terme d'une suite en fonction des
termes précédents. L'objectif de ce chapitre est de développer les techniques
permettant de résoudre de telles \emph{équations de récurrence}\index{equation de recurrence@équation de récurrence},
c'est-à-dire de trouver une formule explicite (dite \emph{forme close}\index{forme close})
pour le terme général.

%% ──────────────────────────────────────────────────────────────────
\section{Récurrences linéaires à coefficients constants}
\label{sec:rec-lineaires}
%% ──────────────────────────────────────────────────────────────────

\subsection{Définitions et cadre général}

\begin{definition}[Récurrence linéaire d'ordre $k$]
\label{def:rec-lineaire}
\index{recurrence lineaire@récurrence linéaire}
Une \textbf{récurrence linéaire d'ordre~$k$ à coefficients constants}
est une relation de la forme
\[
  a_n = c_1\, a_{n-1} + c_2\, a_{n-2} + \cdots + c_k\, a_{n-k} + f(n),
  \qquad n \geq k,
\]
où $c_1, c_2, \ldots, c_k \in \R$ avec $c_k \neq 0$, et $f\colon \N \to \R$.
La récurrence est dite \textbf{homogène}\index{recurrence homogene@récurrence homogène}
si $f(n) = 0$ pour tout~$n$, et \textbf{non homogène} sinon.
\end{definition}

\begin{remarque}
Pour qu'une telle récurrence détermine une suite unique, il faut spécifier
$k$~conditions initiales : $a_0, a_1, \ldots, a_{k-1}$.
\end{remarque}

\subsection{Récurrences homogènes : l'équation caractéristique}
\label{subsec:rec-homogene}

Considérons la récurrence homogène d'ordre~$k$ :
\begin{equation}
\label{eq:rec-hom}
  a_n - c_1\, a_{n-1} - c_2\, a_{n-2} - \cdots - c_k\, a_{n-k} = 0.
\end{equation}

\begin{definition}[Équation caractéristique]
\label{def:eq-carac}
\index{equation caracteristique@équation caractéristique}
L'\textbf{équation caractéristique} associée à~\eqref{eq:rec-hom} est le polynôme
\begin{equation}
\label{eq:poly-carac}
  r^k - c_1\, r^{k-1} - c_2\, r^{k-2} - \cdots - c_k = 0.
\end{equation}
\end{definition}

L'idée fondatrice est de chercher des solutions de la forme $a_n = r^n$ pour une
constante $r \neq 0$. En substituant dans~\eqref{eq:rec-hom}, on obtient
\[
  r^n - c_1\, r^{n-1} - c_2\, r^{n-2} - \cdots - c_k\, r^{n-k} = 0.
\]
En divisant par $r^{n-k} \neq 0$, on retrouve exactement~\eqref{eq:poly-carac}.

\begin{theoreme}[Cas des racines distinctes]
\label{thm:racines-distinctes}
\index{racines distinctes (récurrence)}
Si l'équation caractéristique~\eqref{eq:poly-carac} admet $k$~racines
distinctes $r_1, r_2, \ldots, r_k$ (dans $\C$), alors la solution générale
de~\eqref{eq:rec-hom} est
\[
  a_n = \alpha_1\, r_1^n + \alpha_2\, r_2^n + \cdots + \alpha_k\, r_k^n,
\]
où les constantes $\alpha_1, \ldots, \alpha_k$ sont déterminées par les
conditions initiales.
\end{theoreme}

\begin{proof}
Montrons d'abord que l'ensemble~$\mathcal{S}$ des suites satisfaisant~\eqref{eq:rec-hom}
est un espace vectoriel (sur $\R$ ou $\C$) de dimension~$k$.

\emph{Sous-espace vectoriel.} Si $(a_n)$ et $(b_n)$ satisfont~\eqref{eq:rec-hom},
alors pour tous $\lambda, \mu \in \R$, la suite $(\lambda a_n + \mu b_n)$
satisfait aussi~\eqref{eq:rec-hom} par linéarité.

\emph{Dimension~$k$.} L'application $\varphi \colon \mathcal{S} \to \R^k$ définie
par $\varphi\bigl((a_n)\bigr) = (a_0, a_1, \ldots, a_{k-1})$ est un isomorphisme
d'espaces vectoriels : elle est linéaire, injective (les conditions initiales
déterminent la suite de manière unique) et surjective (pour tout choix de
conditions initiales, il existe une unique suite solution).

Or les $k$~suites $(r_i^n)_{n \geq 0}$ pour $i = 1, \ldots, k$ sont solutions
(par construction de l'équation caractéristique). Elles sont linéairement
indépendantes car la matrice de Vandermonde
\[
  V = \begin{pmatrix}
    1 & 1 & \cdots & 1 \\
    r_1 & r_2 & \cdots & r_k \\
    r_1^2 & r_2^2 & \cdots & r_k^2 \\
    \vdots & \vdots & \ddots & \vdots \\
    r_1^{k-1} & r_2^{k-1} & \cdots & r_k^{k-1}
  \end{pmatrix}
\]
a pour déterminant $\det(V) = \prod_{1 \leq i < j \leq k}(r_j - r_i) \neq 0$
puisque les racines sont distinctes. Ainsi les $(r_i^n)$ forment une base
de~$\mathcal{S}$.
\end{proof}

\begin{theoreme}[Cas des racines répétées]
\label{thm:racines-repetees}
\index{racines repetees@racines répétées (récurrence)}
Si l'équation caractéristique~\eqref{eq:poly-carac} admet les racines
$r_1, r_2, \ldots, r_s$ de multiplicités respectives $m_1, m_2, \ldots, m_s$
(avec $m_1 + m_2 + \cdots + m_s = k$), alors la solution générale est
\[
  a_n = \sum_{i=1}^{s} \left( \alpha_{i,0} + \alpha_{i,1}\, n
    + \alpha_{i,2}\, n^2 + \cdots + \alpha_{i,m_i - 1}\, n^{m_i - 1} \right) r_i^n.
\]
\end{theoreme}

\begin{proof}
Il suffit de montrer que si $r$ est une racine de multiplicité~$m$ du polynôme
caractéristique $p(r)$, alors les suites $n^j r^n$ pour $j = 0, 1, \ldots, m-1$
sont solutions de la récurrence.

Posons $a_n = n^j r^n$. L'opérateur de décalage $E$ agit par $(Ea)_n = a_{n+1}$.
La récurrence~\eqref{eq:rec-hom} se réécrit $p(E)\, a = 0$ où
$p(x) = x^k - c_1 x^{k-1} - \cdots - c_k$. Puisque $r$ est racine de
multiplicité~$m$, on a $p(x) = (x - r)^m\, q(x)$. Or l'opérateur $(E - r)$
appliqué à la suite $n^j r^n$ donne :
\[
  (E - r)(n^j r^n) = (n+1)^j r^{n+1} - r \cdot n^j r^n
  = r^{n+1}\bigl((n+1)^j - n^j\bigr),
\]
qui est de la forme $r^{n+1} \cdot (\text{polynôme de degré } j-1 \text{ en } n)$.
Par récurrence, $(E - r)^m (n^j r^n) = 0$ dès que $j < m$. Donc
$p(E)(n^j r^n) = (E-r)^m q(E)(n^j r^n)$ et il suffit de vérifier que
$(E-r)^m$ annule le résultat, ce qui découle de l'observation précédente
appliquée au polynôme en~$n$ de degré au plus~$j < m$.

L'indépendance linéaire des $k$~solutions ainsi obtenues se vérifie par un
argument de type Vandermonde généralisé (matrice confluente de Vandermonde).
\end{proof}

\begin{exemple}[Récurrence d'ordre~2 à racines distinctes]
\label{ex:rec-ordre2-distinctes}
Résolvons $a_n = 5a_{n-1} - 6a_{n-2}$ avec $a_0 = 1$, $a_1 = 4$.

L'équation caractéristique est $r^2 - 5r + 6 = 0$, soit $(r-2)(r-3) = 0$.
Les racines sont $r_1 = 2$ et $r_2 = 3$. La solution générale est
$a_n = \alpha\, 2^n + \beta\, 3^n$.

Conditions initiales :
\[
  \begin{cases} \alpha + \beta = 1, \\ 2\alpha + 3\beta = 4. \end{cases}
  \implies \alpha = -1,\quad \beta = 2.
\]
Donc $a_n = -2^n + 2 \cdot 3^n = 2 \cdot 3^n - 2^n$.
\end{exemple}

\begin{exemple}[Récurrence à racine double]
\label{ex:rec-racine-double}
Résolvons $a_n = 4a_{n-1} - 4a_{n-2}$ avec $a_0 = 1$, $a_1 = 6$.

L'équation caractéristique est $r^2 - 4r + 4 = 0$, soit $(r - 2)^2 = 0$.
La racine double est $r = 2$. La solution générale est
$a_n = (\alpha + \beta\, n)\, 2^n$.

Conditions initiales :
\[
  \begin{cases} \alpha = 1, \\ (\alpha + \beta)\cdot 2 = 6. \end{cases}
  \implies \alpha = 1,\quad \beta = 2.
\]
Donc $a_n = (1 + 2n)\, 2^n$.
\end{exemple}

\begin{exemple}[Récurrence d'ordre~3]
\label{ex:rec-ordre3}
Résolvons $a_n = 6a_{n-1} - 11a_{n-2} + 6a_{n-3}$ avec $a_0 = 2$,
$a_1 = 5$, $a_2 = 15$.

L'équation caractéristique est $r^3 - 6r^2 + 11r - 6 = 0$, qui se factorise
en $(r-1)(r-2)(r-3) = 0$. La solution générale est
$a_n = \alpha + \beta\, 2^n + \gamma\, 3^n$.

Le système des conditions initiales donne $\alpha = 1$, $\beta = -1$,
$\gamma = 2$, d'où $a_n = 1 - 2^n + 2 \cdot 3^n$.
\end{exemple}

\subsection{Récurrences non homogènes}
\label{subsec:rec-non-hom}

\begin{theoreme}[Structure des solutions d'une récurrence non homogène]
\label{thm:structure-non-hom}
\index{recurrence non homogene@récurrence non homogène}
La solution générale de la récurrence non homogène
\[
  a_n - c_1 a_{n-1} - \cdots - c_k a_{n-k} = f(n)
\]
est de la forme $a_n = a_n^{(h)} + a_n^{(p)}$, où $a_n^{(h)}$ est la solution
générale de la récurrence homogène associée et $a_n^{(p)}$ est une solution
particulière quelconque de la récurrence non homogène.
\end{theoreme}

\begin{proof}
Si $(a_n)$ et $(b_n)$ sont deux solutions de la récurrence non homogène, alors
leur différence $(a_n - b_n)$ satisfait la récurrence homogène, car
\[
  (a_n - b_n) - c_1(a_{n-1} - b_{n-1}) - \cdots - c_k(a_{n-k} - b_{n-k})
  = f(n) - f(n) = 0. \qedhere
\]
\end{proof}

\subsubsection{Méthode des coefficients indéterminés}
\label{subsubsec:coeff-indet}
\index{coefficients indetermines@coefficients indéterminés (méthode)}

La forme de la solution particulière dépend de la forme du second membre~$f(n)$.

\begin{center}
\renewcommand{\arraystretch}{1.4}
\begin{tabular}{@{}ll@{}}
\toprule
\textbf{Forme de $f(n)$} & \textbf{Essai pour $a_n^{(p)}$} \\
\midrule
$d$ (constante) & $A$ \\
$dn + e$ & $An + B$ \\
Polynôme de degré $s$ & Polynôme de degré $s$ \\
$d\, \beta^n$ & $A\, \beta^n$ \\
$n^s \, \beta^n$ & $(A_0 + A_1 n + \cdots + A_s n^s)\, \beta^n$ \\
\bottomrule
\end{tabular}
\end{center}

\begin{remarque}
Si l'essai pour $a_n^{(p)}$ est lui-même solution de la récurrence homogène
(c'est-à-dire si $\beta$ est racine de l'équation caractéristique de
multiplicité~$m$), on multiplie l'essai par $n^m$.
\end{remarque}

\begin{exemple}[Récurrence non homogène avec second membre constant]
\label{ex:non-hom-constante}
Résolvons $a_n = 3a_{n-1} + 4$ avec $a_0 = 1$.

La récurrence homogène $a_n = 3a_{n-1}$ a pour solution $a_n^{(h)} = \alpha\, 3^n$.
On cherche une solution particulière constante $a_n^{(p)} = A$.
En substituant : $A = 3A + 4$, d'où $A = -2$.

Solution générale : $a_n = \alpha\, 3^n - 2$. Avec $a_0 = 1$ :
$\alpha - 2 = 1$, soit $\alpha = 3$.
Conclusion : $a_n = 3^{n+1} - 2$.
\end{exemple}

\begin{exemple}[Récurrence non homogène avec second membre polynomial]
\label{ex:non-hom-poly}
Résolvons $a_n = 2a_{n-1} + n$ avec $a_0 = 0$.

Solution homogène : $a_n^{(h)} = \alpha\, 2^n$. On essaie
$a_n^{(p)} = An + B$. En substituant :
\[
  An + B = 2(A(n-1) + B) + n = (2A + 1)n + (-2A + 2B).
\]
Par identification : $A = 2A + 1$ et $B = -2A + 2B$, d'où $A = -1$, $B = -2$.

Solution générale : $a_n = \alpha\, 2^n - n - 2$. Avec $a_0 = 0$ :
$\alpha = 2$, d'où $a_n = 2^{n+1} - n - 2$.
\end{exemple}

\begin{exemple}[Second membre exponentiel coïncidant avec la racine]
\label{ex:non-hom-expo-racine}
Résolvons $a_n = 2a_{n-1} + 2^n$ avec $a_0 = 1$.

La racine de l'équation caractéristique est $r = 2$, qui coïncide avec la base
de l'exponentielle du second membre. On essaie donc $a_n^{(p)} = An \cdot 2^n$.
En substituant :
\[
  An\, 2^n = 2 \cdot A(n-1)\, 2^{n-1} + 2^n = A(n-1)\, 2^n + 2^n.
\]
Simplifiant par $2^n$ : $An = An - A + 1$, soit $A = 1$.

Solution générale : $a_n = \alpha\, 2^n + n\, 2^n = (\alpha + n)\, 2^n$.
Avec $a_0 = 1$ : $\alpha = 1$, d'où $a_n = (n+1)\, 2^n$.
\end{exemple}

%% ──────────────────────────────────────────────────────────────────
\section{La suite de Fibonacci et la formule de Binet}
\label{sec:fibonacci}
\index{Fibonacci (suite de)}
%% ──────────────────────────────────────────────────────────────────

\begin{definition}[Suite de Fibonacci]
\label{def:fibonacci}
La \textbf{suite de Fibonacci} $(F_n)_{n \geq 0}$ est définie par
\[
  F_0 = 0, \qquad F_1 = 1, \qquad F_n = F_{n-1} + F_{n-2} \quad (n \geq 2).
\]
\end{definition}

Les premiers termes sont : $0, 1, 1, 2, 3, 5, 8, 13, 21, 34, 55, 89, \ldots$

\begin{theoreme}[Formule de Binet]
\label{thm:binet}
\index{Binet (formule de)}
Pour tout $n \geq 0$,
\[
  F_n = \frac{\varphi^n - \psi^n}{\sqrt{5}},
\]
où $\varphi = \dfrac{1+\sqrt{5}}{2}$ (le \emph{nombre d'or}\index{nombre d'or})
et $\psi = \dfrac{1-\sqrt{5}}{2}$.
\end{theoreme}

\begin{proof}
L'équation caractéristique de $F_n = F_{n-1} + F_{n-2}$ est
$r^2 - r - 1 = 0$, de discriminant $\Delta = 5$. Les racines sont
\[
  \varphi = \frac{1+\sqrt{5}}{2} \approx 1{,}618 \quad \text{et} \quad
  \psi = \frac{1-\sqrt{5}}{2} \approx -0{,}618.
\]
Elles sont distinctes, donc $F_n = \alpha\, \varphi^n + \beta\, \psi^n$.
Les conditions initiales donnent le système :
\[
  \begin{cases}
    \alpha + \beta = F_0 = 0, \\
    \alpha\, \varphi + \beta\, \psi = F_1 = 1.
  \end{cases}
\]
De la première équation, $\beta = -\alpha$. En substituant dans la seconde :
$\alpha(\varphi - \psi) = 1$, soit $\alpha\, \sqrt{5} = 1$, d'où
$\alpha = 1/\sqrt{5}$ et $\beta = -1/\sqrt{5}$.
\end{proof}

\begin{corollaire}
\label{cor:fibonacci-arrondi}
Pour tout $n \geq 0$, $F_n$ est l'entier le plus proche de
$\varphi^n / \sqrt{5}$. Plus précisément, $\abs{F_n - \varphi^n/\sqrt{5}} < 1/2$.
\end{corollaire}

\begin{proof}
On a $\abs{\psi} = \abs{(1-\sqrt{5})/2} \approx 0{,}618 < 1$, donc
$\abs{\psi^n/\sqrt{5}} < 1/\sqrt{5} < 1/2$ pour tout $n \geq 0$.
\end{proof}

\begin{proposition}[Identité de Cassini]
\label{prop:cassini}
\index{Cassini (identite de)@Cassini (identité de)}
Pour tout $n \geq 1$,
\[
  F_{n-1}\, F_{n+1} - F_n^2 = (-1)^n.
\]
\end{proposition}

\begin{proof}
Par récurrence sur~$n$. Pour $n = 1$ : $F_0 \cdot F_2 - F_1^2 = 0 \cdot 1 - 1 = -1 = (-1)^1$.

Supposons le résultat vrai au rang~$n$. Alors
\begin{align*}
  F_n F_{n+2} - F_{n+1}^2
  &= F_n(F_{n+1} + F_n) - F_{n+1}^2 \\
  &= F_n F_{n+1} + F_n^2 - F_{n+1}^2 \\
  &= F_n F_{n+1} + F_n^2 - F_{n+1}(F_n + F_{n-1}) \\
  &= F_n^2 - F_{n+1} F_{n-1} \\
  &= -(F_{n-1} F_{n+1} - F_n^2) = -(-1)^n = (-1)^{n+1}. \qedhere
\end{align*}
\end{proof}

%% ──────────────────────────────────────────────────────────────────
\section{Exemples classiques de récurrences}
\label{sec:rec-classiques}
%% ────────────────────────────────────────────────────────────────

\subsection{Les tours de Hanoï}
\label{subsec:hanoi}
\index{tours de Hanoi@tours de Hanoï}

\begin{definition}[Tours de Hanoï]
\label{def:hanoi}
Soit $T_n$ le nombre minimal de déplacements nécessaires pour transférer une tour
de $n$~disques d'un piquet à un autre, selon les règles classiques (un seul disque
déplacé à la fois, jamais un grand disque sur un petit).
\end{definition}

\begin{proposition}
\label{prop:hanoi}
La suite $(T_n)$ satisfait la récurrence $T_n = 2T_{n-1} + 1$ avec $T_1 = 1$,
et a pour forme close $T_n = 2^n - 1$.
\end{proposition}

\begin{proof}
\emph{Récurrence.} Pour déplacer $n$~disques du piquet~$A$ au piquet~$C$ :
\begin{enumerate}[label=(\roman*)]
  \item déplacer les $n-1$ disques supérieurs de $A$ vers $B$ ($T_{n-1}$ coups) ;
  \item déplacer le disque le plus grand de $A$ vers $C$ ($1$ coup) ;
  \item déplacer les $n-1$ disques de $B$ vers $C$ ($T_{n-1}$ coups).
\end{enumerate}
Donc $T_n = 2T_{n-1} + 1$.

\emph{Forme close.} Posons $S_n = T_n + 1$. Alors $S_n = 2S_{n-1}$ avec
$S_1 = 2$, d'où $S_n = 2^n$ et $T_n = 2^n - 1$.
\end{proof}

\subsection{Les nombres de Catalan}
\label{subsec:catalan}
\index{Catalan (nombres de)}

\begin{definition}[Nombres de Catalan]
\label{def:catalan}
Les \textbf{nombres de Catalan} $(C_n)_{n \geq 0}$ sont définis par
$C_0 = 1$ et la récurrence
\begin{equation}
\label{eq:catalan-rec}
  C_{n+1} = \sum_{i=0}^{n} C_i\, C_{n-i}, \qquad n \geq 0.
\end{equation}
\end{definition}

Les premiers termes sont : $1, 1, 2, 5, 14, 42, 132, 429, \ldots$

\begin{theoreme}[Forme close des nombres de Catalan]
\label{thm:catalan-close}
Pour tout $n \geq 0$,
\[
  C_n = \frac{1}{n+1}\binom{2n}{n}.
\]
\end{theoreme}

\begin{proof}
Soit $C(x) = \sum_{n \geq 0} C_n x^n$ la série génératrice des nombres de Catalan.
La récurrence~\eqref{eq:catalan-rec} signifie que le coefficient de $x^{n+1}$
dans $C(x)$ vaut la convolution des coefficients, c'est-à-dire
\[
  C(x) = 1 + x\, C(x)^2.
\]
En résolvant cette équation quadratique en $C(x)$ :
\[
  C(x) = \frac{1 - \sqrt{1 - 4x}}{2x}
\]
(on choisit la racine donnant $C(0) = 1$). En développant $\sqrt{1-4x}$ par la
formule du binôme généralisé
$\sqrt{1 - 4x} = \sum_{n=0}^{\infty} \binom{1/2}{n}(-4x)^n$
et en identifiant les coefficients, on obtient $C_n = \frac{1}{n+1}\binom{2n}{n}$.
\end{proof}

\begin{exemple}[Interprétations combinatoires de $C_n$]
\label{ex:catalan-interp}
Le nombre de Catalan $C_n$ compte, entre autres :
\begin{enumerate}[label=(\roman*)]
  \item le nombre de parenthésages corrects de $n$~paires de parenthèses ;
  \item le nombre d'arbres binaires complets à $n$~nœuds internes ;
  \item le nombre de chemins de Dyck de longueur~$2n$ (chemins monotones de
    $(0,0)$ à $(2n,0)$ par pas $(1,1)$ et $(1,-1)$ ne passant jamais sous
    l'axe des abscisses) ;
  \item le nombre de triangulations d'un polygone convexe à $(n+2)$~sommets.
\end{enumerate}
\end{exemple}

%% ──────────────────────────────────────────────────────────────────
\section{Récurrences de type diviser pour régner}
\label{sec:diviser-regner}
\index{diviser pour regner@diviser pour régner}
%% ──────────────────────────────────────────────────────────────────

De nombreux algorithmes suivent le paradigme \emph{diviser pour régner} : on divise
un problème de taille~$n$ en $a$~sous-problèmes de taille~$n/b$, puis on combine
les résultats en temps $\Theta(n^d)$.

\begin{definition}[Récurrence de diviser pour régner]
\label{def:rec-diviser}
Une récurrence de la forme
\begin{equation}
\label{eq:rec-diviser}
  T(n) = a\, T(n/b) + f(n),
\end{equation}
avec $a \geq 1$, $b > 1$ et $f$ une fonction positive, est appelée
\textbf{récurrence de diviser pour régner}\index{recurrence diviser pour regner@récurrence diviser pour régner}.
\end{definition}

\begin{theoreme}[Théorème maître]
\label{thm:maitre}
\index{theoreme maitre@théorème maître}
Soit $T(n) = a\, T(n/b) + \Theta(n^d)$ avec $a \geq 1$, $b > 1$, $d \geq 0$.
Posons $c = \log_b a$. Alors :
\[
  T(n) = \begin{cases}
    \Theta(n^d) & \text{si } d > c, \\[4pt]
    \Theta(n^d \log n) & \text{si } d = c, \\[4pt]
    \Theta(n^c) & \text{si } d < c.
  \end{cases}
\]
\end{theoreme}

\begin{remarque}
Intuitivement, on compare le coût de la combinaison ($n^d$) au nombre de
feuilles de l'arbre de récursion ($a^{\log_b n} = n^{\log_b a} = n^c$).
Si le coût de combinaison domine ($d > c$), $T(n) = \Theta(n^d)$.
Si les feuilles dominent ($d < c$), $T(n) = \Theta(n^c)$.
Si les deux sont équilibrés ($d = c$), on accumule un facteur logarithmique.
\end{remarque}

\begin{exemple}[Tri par fusion]
\label{ex:tri-fusion}
\index{tri par fusion}
Le tri par fusion satisfait $T(n) = 2T(n/2) + \Theta(n)$, soit $a = 2$,
$b = 2$, $d = 1$. On a $c = \log_2 2 = 1 = d$, donc
$T(n) = \Theta(n \log n)$.
\end{exemple}

\begin{exemple}[Recherche binaire]
\label{ex:rech-binaire}
\index{recherche binaire}
La recherche binaire satisfait $T(n) = T(n/2) + \Theta(1)$, soit $a = 1$,
$b = 2$, $d = 0$. On a $c = \log_2 1 = 0 = d$, donc
$T(n) = \Theta(\log n)$.
\end{exemple}

\begin{exemple}[Multiplication de Karatsuba]
\label{ex:karatsuba}
\index{Karatsuba}
L'algorithme de Karatsuba satisfait $T(n) = 3T(n/2) + \Theta(n)$, soit
$a = 3$, $b = 2$, $d = 1$. On a $c = \log_2 3 \approx 1{,}585 > 1 = d$,
donc $T(n) = \Theta(n^{\log_2 3})$.
\end{exemple}

\begin{exemple}[Multiplication matricielle de Strassen]
\label{ex:strassen}
\index{Strassen}
L'algorithme de Strassen satisfait $T(n) = 7T(n/2) + \Theta(n^2)$, soit
$a = 7$, $b = 2$, $d = 2$. On a $c = \log_2 7 \approx 2{,}807 > 2 = d$,
donc $T(n) = \Theta(n^{\log_2 7})$.
\end{exemple}

%% ──────────────────────────────────────────────────────────────────
\section{Résolution par fonctions génératrices}
\label{sec:rec-gen}
\index{fonctions generatrices@fonctions génératrices!résolution de récurrences}
%% ──────────────────────────────────────────────────────────────────

Les fonctions génératrices (voir chapitre~6) fournissent un outil puissant et
systématique pour résoudre les récurrences. La méthode se décompose en trois
étapes.

\begin{enumerate}[label=\textbf{Étape \arabic*.}]
  \item \textbf{Traduire} la récurrence en une équation pour la série
    génératrice $A(x) = \sum_{n \geq 0} a_n x^n$.
  \item \textbf{Résoudre} l'équation algébrique (ou fonctionnelle) en $A(x)$.
  \item \textbf{Extraire} les coefficients de $A(x)$ par décomposition en
    fractions partielles ou développement en série.
\end{enumerate}

\begin{exemple}[Fibonacci par fonctions génératrices]
\label{ex:fib-gen}
Soit $F(x) = \sum_{n \geq 0} F_n x^n$. La récurrence $F_n = F_{n-1} + F_{n-2}$
(pour $n \geq 2$) avec $F_0 = 0$, $F_1 = 1$ se traduit par :
\[
  F(x) - F_0 - F_1 x = x\bigl(F(x) - F_0\bigr) + x^2 F(x),
\]
soit $F(x) - x = xF(x) + x^2 F(x)$. D'où
\[
  F(x) = \frac{x}{1 - x - x^2}.
\]
En décomposant en fractions partielles avec $1 - x - x^2 = -(x - \varphi^{-1})(x - \psi^{-1})$ (en notant que $-1/\varphi = \psi$ et $-1/\psi = \varphi$ ne sont pas les racines directes du dénominateur, mais $1/\varphi$ et $1/\psi$ le sont en inversant), on retrouve la formule de Binet.
\end{exemple}

\begin{exemple}[Récurrence $a_n = 3a_{n-1} + 2$ par fonctions génératrices]
\label{ex:gen-non-hom}
Soit $a_0 = 1$ et $a_n = 3a_{n-1} + 2$ pour $n \geq 1$.
Posons $A(x) = \sum_{n \geq 0} a_n x^n$. Alors
\[
  A(x) - 1 = 3x\, A(x) + \frac{2x}{1-x}.
\]
D'où
\[
  A(x) = \frac{1 - x + 2x}{(1-x)(1-3x)} = \frac{1+x}{(1-x)(1-3x)}.
\]
Par fractions partielles : $A(x) = \frac{-1}{1-x} + \frac{2}{1-3x}$, d'où
$a_n = -1 + 2 \cdot 3^n = 2 \cdot 3^n - 1$.

Vérification : $a_0 = 2 - 1 = 1$, $a_1 = 6 - 1 = 5 = 3 \cdot 1 + 2$. \checkmark
\end{exemple}

%% ──────────────────────────────────────────────────────────────────
\section{Exercices}
\label{sec:rec-exercices}
%% ──────────────────────────────────────────────────────────────────

\begin{exercice}[Récurrence d'ordre~2]
\label{exo:rec-ordre2}
Résoudre les récurrences suivantes avec les conditions initiales données :
\begin{enumerate}[label=(\alph*)]
  \item $a_n = 7a_{n-1} - 10a_{n-2}$, \quad $a_0 = 2$, $a_1 = 1$.
  \item $a_n = 6a_{n-1} - 9a_{n-2}$, \quad $a_0 = 1$, $a_1 = 6$.
  \item $a_n = -2a_{n-1} - a_{n-2}$, \quad $a_0 = 0$, $a_1 = 1$.
\end{enumerate}
\end{exercice}

\begin{exercice}[Récurrences non homogènes]
\label{exo:rec-non-hom}
Résoudre les récurrences suivantes :
\begin{enumerate}[label=(\alph*)]
  \item $a_n = 2a_{n-1} + 3$, \quad $a_0 = 1$.
  \item $a_n = 3a_{n-1} + n$, \quad $a_0 = 0$.
  \item $a_n = 2a_{n-1} + 3^n$, \quad $a_0 = 1$.
  \item $a_n = 4a_{n-1} - 4a_{n-2} + 2^n$, \quad $a_0 = 0$, $a_1 = 1$.
\end{enumerate}
\end{exercice}

\begin{exercice}[Identités de Fibonacci]
\label{exo:fib-identites}
Montrer les identités suivantes pour la suite de Fibonacci :
\begin{enumerate}[label=(\alph*)]
  \item $\sum_{k=0}^{n} F_k = F_{n+2} - 1$.
  \item $\sum_{k=0}^{n} F_k^2 = F_n\, F_{n+1}$.
  \item $F_{m+n} = F_{m-1}\, F_n + F_m\, F_{n+1}$ pour $m \geq 1$.
  \item $\pgcd(F_m, F_n) = F_{\pgcd(m,n)}$.
\end{enumerate}
\end{exercice}

\begin{exercice}[Théorème maître]
\label{exo:maitre}
Utiliser le théorème maître pour déterminer l'ordre de grandeur asymptotique
de $T(n)$ dans chacun des cas suivants :
\begin{enumerate}[label=(\alph*)]
  \item $T(n) = 4T(n/2) + n$.
  \item $T(n) = 4T(n/2) + n^2$.
  \item $T(n) = 4T(n/2) + n^3$.
  \item $T(n) = 8T(n/2) + n^2$.
  \item $T(n) = 2T(n/4) + \sqrt{n}$.
\end{enumerate}
\end{exercice}

\begin{exercice}[Suite de Lucas]
\label{exo:lucas}
\index{Lucas (suite de)}
La \textbf{suite de Lucas} est définie par $L_0 = 2$, $L_1 = 1$ et
$L_n = L_{n-1} + L_{n-2}$.
\begin{enumerate}[label=(\alph*)]
  \item Trouver la forme close de $L_n$.
  \item Montrer que $L_n = F_{n-1} + F_{n+1}$ pour $n \geq 1$.
  \item Montrer que $F_{2n} = F_n L_n$.
\end{enumerate}
\end{exercice}

\begin{exercice}[Problème des régions du plan]
\label{exo:regions-plan}
Soit $R_n$ le nombre maximal de régions créées dans le plan par $n$~droites
en position générale (pas deux parallèles, pas trois concourantes).
\begin{enumerate}[label=(\alph*)]
  \item Montrer que $R_n = R_{n-1} + n$ avec $R_0 = 1$.
  \item En déduire que $R_n = \binom{n}{2} + n + 1 = \frac{n^2 + n + 2}{2}$.
\end{enumerate}
\end{exercice}

\begin{exercice}[Fonctions génératrices et récurrences]
\label{exo:gen-rec}
Résoudre par la méthode des fonctions génératrices :
\begin{enumerate}[label=(\alph*)]
  \item $a_n = 5a_{n-1} - 6a_{n-2}$ avec $a_0 = 2$, $a_1 = 5$.
  \item $a_n = a_{n-1} + 2a_{n-2}$ avec $a_0 = 0$, $a_1 = 1$.
\end{enumerate}
\end{exercice}

\begin{exercice}[Nombres de Catalan]
\label{exo:catalan}
\begin{enumerate}[label=(\alph*)]
  \item Vérifier que $C_4 = 14$ en énumérant les parenthésages de $4$~paires
    de parenthèses.
  \item Montrer que $C_n = \frac{4n-2}{n+1}\, C_{n-1}$ pour $n \geq 1$.
  \item Montrer que $C_n \sim \frac{4^n}{n^{3/2}\sqrt{\pi}}$ quand $n \to \infty$.
\end{enumerate}
\end{exercice}

%% ──────────────────────────────────────────────────────────────────
\section*{Résumé du chapitre}
\addcontentsline{toc}{section}{Résumé du chapitre}
%% ──────────────────────────────────────────────────────────────────

\begin{itemize}
  \item Une \textbf{récurrence linéaire homogène} d'ordre~$k$ se résout via son
    \textbf{équation caractéristique} : racines distinctes $\to$ combinaison de
    puissances ; racines répétées $\to$ on multiplie par des polynômes en~$n$.
  \item Pour une récurrence \textbf{non homogène}, la solution générale est
    la somme de la solution homogène et d'une solution particulière, trouvée par
    la \textbf{méthode des coefficients indéterminés}.
  \item La \textbf{formule de Binet} donne la forme close de la suite de
    Fibonacci : $F_n = (\varphi^n - \psi^n)/\sqrt{5}$.
  \item Les \textbf{tours de Hanoï} nécessitent $2^n - 1$ déplacements ; les
    \textbf{nombres de Catalan} valent $C_n = \frac{1}{n+1}\binom{2n}{n}$.
  \item Le \textbf{théorème maître} résout les récurrences de diviser pour régner
    $T(n) = aT(n/b) + \Theta(n^d)$ en comparant $d$ à $\log_b a$.
  \item Les \textbf{fonctions génératrices} offrent une méthode systématique :
    traduire la récurrence en équation algébrique, résoudre, puis extraire les
    coefficients.
\end{itemize}

\newpage

%%%%%%%%%%%%%%%%%%%%%%%%%%%%%%%%%%%%%%%%%%%%%%%%%%%%%%%%%%%%%%%%%%%%
%  CHAPITRE 8 — Théorie des Graphes : Notions de Base
%%%%%%%%%%%%%%%%%%%%%%%%%%%%%%%%%%%%%%%%%%%%%%%%%%%%%%%%%%%%%%%%%%%%
\chapter{Théorie des Graphes --- Notions de Base}
\label{ch:graphes}

\minitoc

\section*{Introduction}
\addcontentsline{toc}{section}{Introduction}

La théorie des graphes\index{graphe} est l'une des branches les plus vibrantes
et les plus applicables des mathématiques discrètes. Née d'un problème
récréatif posé par Euler au XVIII\textsuperscript{e}~siècle, elle est devenue
un outil fondamental en informatique, en recherche opérationnelle, en biologie
computationnelle et dans bien d'autres domaines.

%% ──────────────────────────────────────────────────────────────────
\section{Perspective historique : les ponts de Königsberg}
\label{sec:konigsberg}
\index{Konigsberg (ponts de)@Königsberg (ponts de)}
%% ──────────────────────────────────────────────────────────────────

En 1736, Leonhard Euler\index{Euler} résolut le célèbre \emph{problème des ponts
de Königsberg} : est-il possible de traverser chacun des sept ponts de la ville
exactement une fois et de revenir au point de départ ? Euler montra que c'est
impossible et, ce faisant, posa les fondements de la théorie des graphes.

\begin{center}
\begin{tikzpicture}[
  scale=1.2,
  every node/.style={draw, circle, fill=blue!20, minimum size=8mm, inner sep=0pt,
    font=\small\bfseries},
  every edge/.style={draw, thick}
]
  % Sommets représentant les régions de Königsberg
  \node (A) at (0, 2) {$A$};
  \node (B) at (0, 0) {$B$};
  \node (C) at (-2.5, 1) {$C$};
  \node (D) at (2.5, 1) {$D$};

  % Arêtes représentant les ponts
  \draw[thick] (A) to[bend right=20] (C);
  \draw[thick] (A) to[bend left=20] (C);
  \draw[thick] (B) to[bend right=20] (C);
  \draw[thick] (B) to[bend left=20] (C);
  \draw[thick] (A) to[bend left=20] (D);
  \draw[thick] (B) to[bend right=20] (D);
  \draw[thick] (A) -- (B);
\end{tikzpicture}
\captionof{figure}{Le multigraphe des ponts de Königsberg : chaque sommet
  représente une masse terrestre, chaque arête un pont.}
\label{fig:konigsberg}
\end{center}

L'argument d'Euler repose sur une observation simple : pour qu'un tel parcours
(appelé aujourd'hui \emph{circuit eulérien}\index{circuit eulerien@circuit eulérien})
existe, chaque sommet doit avoir un degré pair, ce qui n'est le cas d'aucun
sommet dans le graphe de Königsberg.

%% ──────────────────────────────────────────────────────────────────
\section{Définitions fondamentales}
\label{sec:graph-def}
%% ──────────────────────────────────────────────────────────────────

\begin{definition}[Graphe simple]
\label{def:graphe-simple}
\index{graphe!simple}
Un \textbf{graphe simple} (ou simplement un \textbf{graphe}) est un couple
$G = (V, E)$ où :
\begin{itemize}
  \item $V$ est un ensemble fini non vide dont les éléments sont appelés
    \textbf{sommets}\index{sommet} (ou \emph{nœuds}\index{noeud@nœud}) ;
  \item $E \subseteq \binom{V}{2} = \set{\{u,v\} : u, v \in V,\; u \neq v}$
    est un ensemble de paires non ordonnées de sommets distincts, appelées
    \textbf{arêtes}\index{arete@arête}.
\end{itemize}
\end{definition}

\begin{notation}
\label{not:graph}
On note $\abs{V} = n$ le \textbf{nombre de sommets} (ou \emph{ordre}\index{ordre (graphe)})
et $\abs{E} = m$ le \textbf{nombre d'arêtes} (ou \emph{taille}\index{taille (graphe)})
du graphe.
\end{notation}

\begin{definition}[Adjacence et incidence]
\label{def:adjacence}
\index{adjacence}\index{incidence}
Soit $G = (V, E)$ un graphe et $e = \{u, v\} \in E$ une arête.
\begin{itemize}
  \item Les sommets $u$ et $v$ sont dits \textbf{adjacents} (ou \emph{voisins}).
  \item L'arête $e$ est \textbf{incidente} aux sommets $u$ et $v$.
  \item On note $N(v) = \set{u \in V : \{u,v\} \in E}$ le \textbf{voisinage}
    de~$v$\index{voisinage}.
\end{itemize}
\end{definition}

\subsection{Variantes de graphes}

\begin{definition}[Multigraphe]
\label{def:multigraphe}
\index{multigraphe}
Un \textbf{multigraphe} est un graphe dans lequel on autorise :
\begin{itemize}
  \item les \textbf{arêtes multiples}\index{arete multiple@arête multiple}
    (plusieurs arêtes entre la même paire de sommets) ;
  \item les \textbf{boucles}\index{boucle} (arêtes reliant un sommet à lui-même).
\end{itemize}
Formellement, $E$ est un multi-ensemble de paires (non nécessairement distinctes)
d'éléments de~$V$.
\end{definition}

\begin{definition}[Graphe orienté]
\label{def:digraphe}
\index{graphe!oriente@orienté}\index{digraphe}
Un \textbf{graphe orienté} (ou \textbf{digraphe}) est un couple $G = (V, A)$ où
$A \subseteq V \times V$ est un ensemble de couples ordonnés, appelés
\textbf{arcs}\index{arc}. Un arc $(u, v)$ va de~$u$ (origine) à~$v$ (extrémité).
\end{definition}

\begin{definition}[Graphe pondéré]
\label{def:graphe-pondere}
\index{graphe!pondere@pondéré}
Un \textbf{graphe pondéré} est un graphe $G = (V, E)$ muni d'une fonction de
poids $w \colon E \to \R$ qui attribue un poids (ou coût) à chaque arête.
\end{definition}

%% ──────────────────────────────────────────────────────────────────
\section{Représentations d'un graphe}
\label{sec:representations}
\index{representation d'un graphe@représentation d'un graphe}
%% ──────────────────────────────────────────────────────────────────

\subsection{Matrice d'adjacence}
\label{subsec:matrice-adj}
\index{matrice d'adjacence}

\begin{definition}[Matrice d'adjacence]
\label{def:matrice-adj}
Soit $G = (V, E)$ un graphe simple avec $V = \{v_1, \ldots, v_n\}$.
La \textbf{matrice d'adjacence} de~$G$ est la matrice $A \in \{0,1\}^{n \times n}$
définie par
\[
  A_{ij} = \begin{cases} 1 & \text{si } \{v_i, v_j\} \in E, \\ 0 & \text{sinon.} \end{cases}
\]
\end{definition}

\begin{remarque}
Pour un graphe simple, $A$ est symétrique ($A_{ij} = A_{ji}$) et à diagonale
nulle ($A_{ii} = 0$). L'espace mémoire est $\Theta(n^2)$. La matrice d'adjacence
est efficace pour tester l'existence d'une arête en $O(1)$, mais coûteuse en
mémoire pour les graphes creux.
\end{remarque}

\begin{exemple}
\label{ex:matrice-adj-K4}
La matrice d'adjacence du graphe complet $K_4$ est
\[
  A = \begin{pmatrix}
    0 & 1 & 1 & 1 \\
    1 & 0 & 1 & 1 \\
    1 & 1 & 0 & 1 \\
    1 & 1 & 1 & 0
  \end{pmatrix}.
\]
\end{exemple}

\subsection{Listes d'adjacence}
\label{subsec:listes-adj}
\index{liste d'adjacence}

\begin{definition}[Liste d'adjacence]
\label{def:liste-adj}
La \textbf{représentation par listes d'adjacence} d'un graphe $G = (V, E)$
consiste à associer à chaque sommet $v \in V$ la liste de ses voisins $N(v)$.
\end{definition}

\begin{remarque}
L'espace mémoire est $\Theta(n + m)$. Cette représentation est préférable pour
les graphes creux ($m \ll n^2$), qui sont les plus fréquents en pratique.
\end{remarque}

\subsection{Matrice d'incidence}
\label{subsec:matrice-inc}
\index{matrice d'incidence}

\begin{definition}[Matrice d'incidence]
\label{def:matrice-inc}
Soit $G = (V, E)$ un graphe avec $V = \{v_1, \ldots, v_n\}$ et
$E = \{e_1, \ldots, e_m\}$. La \textbf{matrice d'incidence} $B \in \{0,1\}^{n \times m}$
est définie par
\[
  B_{ij} = \begin{cases} 1 & \text{si le sommet $v_i$ est incident à l'arête $e_j$,} \\
    0 & \text{sinon.}
  \end{cases}
\]
\end{definition}

\begin{remarque}
Chaque colonne de $B$ contient exactement deux~$1$ (les deux extrémités de
l'arête correspondante). La somme de la ligne~$i$ donne le degré de~$v_i$.
\end{remarque}

%% ──────────────────────────────────────────────────────────────────
\section{Degré d'un sommet}
\label{sec:degre}
%% ──────────────────────────────────────────────────────────────────

\begin{definition}[Degré]
\label{def:degre}
\index{degre@degré}
Le \textbf{degré} d'un sommet~$v$ dans un graphe $G = (V, E)$, noté
$\deg(v)$\index{deg@$\deg$}, est le nombre d'arêtes incidentes à~$v$ :
\[
  \deg(v) = \abs{N(v)} = \card\set{e \in E : v \in e}.
\]
Un sommet de degré~$0$ est dit \textbf{isolé}\index{sommet!isole@isolé},
un sommet de degré~$1$ est une \textbf{feuille}\index{feuille} (ou sommet pendant).
\end{definition}

\begin{definition}[Degré dans un graphe orienté]
\label{def:degre-oriente}
\index{degre entrant@degré entrant}\index{degre sortant@degré sortant}
Dans un digraphe $G = (V, A)$, on distingue :
\begin{itemize}
  \item le \textbf{degré entrant} de~$v$ :
    $\deg^-(v) = \card\set{u \in V : (u,v) \in A}$ ;
  \item le \textbf{degré sortant} de~$v$ :
    $\deg^+(v) = \card\set{u \in V : (v,u) \in A}$.
\end{itemize}
\end{definition}

\begin{theoreme}[Lemme des poignées de main]
\label{thm:handshaking}
\index{lemme des poignees de main@lemme des poignées de main}
\index{handshaking lemma}
Soit $G = (V, E)$ un graphe. Alors
\[
  \sum_{v \in V} \deg(v) = 2\, \abs{E}.
\]
\end{theoreme}

\begin{proof}
Chaque arête $e = \{u, v\}$ contribue exactement~$1$ au degré de~$u$ et~$1$ au
degré de~$v$. Donc, en sommant les degrés de tous les sommets, chaque arête est
comptée exactement deux fois.

Formellement, on peut écrire
\[
  \sum_{v \in V} \deg(v) = \sum_{v \in V} \sum_{e \in E} \mathbf{1}_{v \in e}
  = \sum_{e \in E} \sum_{v \in V} \mathbf{1}_{v \in e}
  = \sum_{e \in E} 2 = 2\, \abs{E}. \qedhere
\]
\end{proof}

\begin{corollaire}
\label{cor:degre-pair}
Dans tout graphe, le nombre de sommets de degré impair est pair.
\end{corollaire}

\begin{proof}
Séparons les sommets en deux ensembles : $V_{\text{pair}}$ (sommets de degré
pair) et $V_{\text{impair}}$ (sommets de degré impair). Alors
\[
  2\abs{E} = \sum_{v \in V_{\text{pair}}} \deg(v)
  + \sum_{v \in V_{\text{impair}}} \deg(v).
\]
Le membre de gauche est pair, et la première somme du membre de droite est pair
(somme de nombres pairs). Donc la seconde somme est pair. Comme chaque terme de
cette somme est impair, le nombre de termes doit être pair.
\end{proof}

\begin{exemple}
\label{ex:degre}
Dans le graphe de Königsberg (figure~\ref{fig:konigsberg}), les degrés sont :
$\deg(A) = 3$, $\deg(B) = 3$, $\deg(C) = 4$, $\deg(D) = 2$ (dans le cas d'un
multigraphe à 7~arêtes : $3+3+4+2 \neq 14$). En fait, avec les 7~ponts :
$\deg(A) = 3$, $\deg(B) = 3$, $\deg(C) = 3$, $\deg(D) = 5$ dans la version
historique originale. La somme des degrés est $3+3+3+5 = 14 = 2 \times 7$.

\emph{Remarque :} comme les quatre sommets ont tous un degré impair, aucun
circuit eulérien n'existe.
\end{exemple}

%% ──────────────────────────────────────────────────────────────────
\section{Chemins, cycles et connexité}
\label{sec:chemins}
%% ──────────────────────────────────────────────────────────────────

\begin{definition}[Marche, chemin, cycle]
\label{def:marche-chemin-cycle}
\index{marche}\index{chemin}\index{cycle}
Soit $G = (V, E)$ un graphe.
\begin{itemize}
  \item Une \textbf{marche} de longueur~$k$ est une suite
    $v_0, e_1, v_1, e_2, \ldots, e_k, v_k$ alternant sommets et arêtes, où
    $e_i = \{v_{i-1}, v_i\} \in E$ pour tout~$i$. On la note souvent
    $v_0 v_1 \cdots v_k$.
  \item Un \textbf{chemin} (ou \emph{chaîne élémentaire}) est une marche dont
    tous les sommets sont distincts.
  \item Un \textbf{cycle} (ou \emph{circuit élémentaire}) est une marche fermée
    ($v_0 = v_k$) de longueur $k \geq 3$ dont tous les sommets intermédiaires
    sont distincts.
\end{itemize}
\end{definition}

\begin{definition}[Connexité]
\label{def:connexite}
\index{connexite@connexité}\index{graphe!connexe}
Un graphe $G = (V, E)$ est \textbf{connexe} si, pour toute paire de sommets
$u, v \in V$, il existe un chemin de~$u$ à~$v$ dans~$G$.
\end{definition}

\begin{definition}[Composantes connexes]
\label{def:comp-connexe}
\index{composante connexe}
Les \textbf{composantes connexes} d'un graphe~$G$ sont les sous-graphes connexes
maximaux de~$G$. Formellement, ce sont les classes d'équivalence de la relation
\og être relié par un chemin\fg{}, qui est une relation d'équivalence sur~$V$.
\end{definition}

\begin{proposition}
\label{prop:comp-connexe-equiv}
La relation $u \sim v$ définie par \og il existe un chemin de~$u$ à~$v$\fg{}
est une relation d'équivalence sur $V$.
\end{proposition}

\begin{proof}
\emph{Réflexivité} : le chemin de longueur~$0$ relie $v$ à lui-même.
\emph{Symétrie} : si $v_0 v_1 \cdots v_k$ est un chemin de $u$ à $v$, alors
$v_k v_{k-1} \cdots v_0$ est un chemin de $v$ à $u$.
\emph{Transitivité} : si un chemin relie $u$ à $v$ et un autre $v$ à $w$,
leur concaténation contient une marche de $u$ à $w$, et toute marche contient
un chemin (en supprimant les boucles).
\end{proof}

\begin{proposition}
\label{prop:arete-connexe}
Un graphe connexe d'ordre~$n$ possède au moins $n - 1$ arêtes.
\end{proposition}

\begin{proof}
Par récurrence sur~$n$. Pour $n = 1$, $m \geq 0 = n - 1$. Supposons $n \geq 2$
et soit $v$ un sommet de $G$. Le graphe $G - v$ a au plus $\deg(v)$~composantes
connexes (chacune devait être reliée à~$v$). Pour que $G$ soit connexe, il faut
au moins une arête entre $v$ et chaque composante. Par hypothèse de récurrence,
les composantes de $G - v$ (ayant au total $n-1$ sommets) nécessitent au moins
$(n-1) - (\text{nombre de composantes})$ arêtes en tout. En ajoutant les arêtes
incidentes à~$v$, on obtient $m \geq (n-1-c) + c = n - 1$ où $c$ est le nombre
de composantes de $G - v$.
\end{proof}

\begin{definition}[Distance dans un graphe]
\label{def:distance-graphe}
\index{distance (graphe)}
La \textbf{distance} entre deux sommets $u$ et $v$ dans un graphe~$G$, notée
$d(u,v)$, est la longueur du plus court chemin de~$u$ à~$v$. Si aucun chemin
n'existe, on pose $d(u,v) = +\infty$.
\end{definition}

%% ──────────────────────────────────────────────────────────────────
\section{Isomorphisme de graphes}
\label{sec:isomorphisme}
\index{isomorphisme de graphes}
%% ──────────────────────────────────────────────────────────────────

\begin{definition}[Isomorphisme]
\label{def:isomorphisme-graphe}
Deux graphes $G_1 = (V_1, E_1)$ et $G_2 = (V_2, E_2)$ sont \textbf{isomorphes},
noté $G_1 \simeq G_2$, s'il existe une bijection $\varphi \colon V_1 \to V_2$
telle que
\[
  \{u, v\} \in E_1 \iff \{\varphi(u), \varphi(v)\} \in E_2.
\]
Une telle bijection est appelée un \textbf{isomorphisme}\index{isomorphisme}.
\end{definition}

\begin{remarque}
Deux graphes isomorphes ont nécessairement :
\begin{itemize}
  \item le même nombre de sommets et le même nombre d'arêtes ;
  \item la même suite des degrés (triée) ;
  \item le même nombre de composantes connexes ;
  \item les mêmes longueurs de cycles.
\end{itemize}
Ces conditions sont nécessaires mais non suffisantes.
\end{remarque}

%% ──────────────────────────────────────────────────────────────────
\section{Familles classiques de graphes}
\label{sec:familles}
%% ──────────────────────────────────────────────────────────────────

\subsection{Graphes complets}
\label{subsec:graphe-complet}
\index{graphe!complet}\index{Kn@$K_n$}

\begin{definition}[Graphe complet $K_n$]
\label{def:graphe-complet}
Le \textbf{graphe complet} $K_n$ est le graphe simple sur $n$~sommets dans lequel
chaque paire de sommets distincts est reliée par une arête. Il possède
$\binom{n}{2} = \frac{n(n-1)}{2}$ arêtes.
\end{definition}

\begin{center}
\begin{tikzpicture}[
  every node/.style={draw, circle, fill=blue!20, minimum size=6mm,
    inner sep=0pt, font=\footnotesize},
  every edge/.style={draw, thick}
]
  % K_4
  \begin{scope}[shift={(-4,0)}]
    \node[label=below:{$K_4$}] at (0,-1.6) [draw=none, fill=none] {};
    \node (a) at (90:1.2)  {1};
    \node (b) at (0:1.2)   {2};
    \node (c) at (270:1.2) {3};
    \node (d) at (180:1.2) {4};
    \draw (a) -- (b) -- (c) -- (d) -- (a) -- (c);
    \draw (b) -- (d);
  \end{scope}

  % K_5
  \begin{scope}[shift={(0,0)}]
    \node[label=below:{$K_5$}] at (0,-1.6) [draw=none, fill=none] {};
    \foreach \i in {1,...,5} {
      \node (v\i) at ({90 + (\i-1)*72}:1.2) {\i};
    }
    \foreach \i in {1,...,5} {
      \foreach \j in {\i,...,5} {
        \draw (v\i) -- (v\j);
      }
    }
  \end{scope}

  % K_3
  \begin{scope}[shift={(4,0)}]
    \node[label=below:{$K_3$}] at (0,-1.6) [draw=none, fill=none] {};
    \node (w1) at (90:1.2)  {1};
    \node (w2) at (210:1.2) {2};
    \node (w3) at (330:1.2) {3};
    \draw (w1) -- (w2) -- (w3) -- (w1);
  \end{scope}
\end{tikzpicture}
\captionof{figure}{Les graphes complets $K_3$, $K_4$ et $K_5$.}
\label{fig:graphes-complets}
\end{center}

\subsection{Graphes bipartis}
\label{subsec:graphe-biparti}
\index{graphe!biparti}\index{Kmn@$K_{m,n}$}

\begin{definition}[Graphe biparti]
\label{def:graphe-biparti}
Un graphe $G = (V, E)$ est \textbf{biparti} s'il existe une partition
$V = X \sqcup Y$ telle que toute arête de~$E$ relie un sommet de~$X$ à un
sommet de~$Y$.
\end{definition}

\begin{definition}[Graphe biparti complet $K_{m,n}$]
\label{def:biparti-complet}
Le \textbf{graphe biparti complet} $K_{m,n}$ est le graphe biparti avec
$\abs{X} = m$, $\abs{Y} = n$, et toute paire $(x, y) \in X \times Y$ est
reliée. Il possède $mn$~arêtes.
\end{definition}

\begin{center}
\begin{tikzpicture}[
  every node/.style={draw, circle, fill=red!20, minimum size=6mm,
    inner sep=0pt, font=\footnotesize},
  every edge/.style={draw, thick}
]
  % K_{3,3}
  \node[label=below:{$K_{3,3}$}] at (1,-2.2) [draw=none, fill=none] {};
  \foreach \i in {1,2,3} {
    \node[fill=blue!20] (x\i) at ({(\i-1)*1.2}, 1.2) {$x_\i$};
    \node[fill=red!20]  (y\i) at ({(\i-1)*1.2}, -1.2) {$y_\i$};
  }
  \foreach \i in {1,2,3} {
    \foreach \j in {1,2,3} {
      \draw (x\i) -- (y\j);
    }
  }

  % K_{2,3}
  \begin{scope}[shift={(5.5,0)}]
    \node[label=below:{$K_{2,3}$}] at (1,-2.2) [draw=none, fill=none] {};
    \foreach \i in {1,2} {
      \node[fill=blue!20] (a\i) at ({(\i-1)*1.5 + 0.3}, 1.2) {$x_\i$};
    }
    \foreach \j in {1,2,3} {
      \node[fill=red!20] (b\j) at ({(\j-1)*1.2}, -1.2) {$y_\j$};
    }
    \foreach \i in {1,2} {
      \foreach \j in {1,2,3} {
        \draw (a\i) -- (b\j);
      }
    }
  \end{scope}
\end{tikzpicture}
\captionof{figure}{Les graphes bipartis complets $K_{3,3}$ et $K_{2,3}$.}
\label{fig:bipartis}
\end{center}

\begin{proposition}[Caractérisation des graphes bipartis]
\label{prop:biparti-carac}
\index{graphe!biparti!caractérisation}
Un graphe est biparti si et seulement s'il ne contient aucun cycle de
longueur impaire.
\end{proposition}

\begin{proof}
$(\Rightarrow)$ Dans un graphe biparti avec partition $V = X \sqcup Y$, tout
cycle alterne entre $X$ et $Y$, donc sa longueur est paire.

$(\Leftarrow)$ Supposons que $G$ est connexe et ne contient aucun cycle impair.
Fixons un sommet~$s$ et posons $X = \set{v : d(s,v) \text{ pair}}$ et
$Y = \set{v : d(s,v) \text{ impair}}$. Si une arête $\{u,v\}$ reliait deux
sommets de~$X$ (resp.\ de~$Y$), un plus court chemin de~$s$ à~$u$, l'arête
$\{u,v\}$, et un plus court chemin de~$v$ à~$s$ formeraient un cycle de longueur
$d(s,u) + 1 + d(s,v)$, qui serait impaire puisque $d(s,u)$ et $d(s,v)$ ont même
parité. Contradiction.
\end{proof}

\subsection{Graphes cycles et chemins}
\label{subsec:cycles-chemins}
\index{Cn@$C_n$}\index{Pn@$P_n$}

\begin{definition}[Graphe cycle $C_n$ et graphe chemin $P_n$]
\label{def:cycle-chemin}
\begin{itemize}
  \item Le \textbf{graphe cycle} $C_n$ ($n \geq 3$) a pour sommets
    $\{v_1, \ldots, v_n\}$ et pour arêtes
    $\{v_1 v_2, v_2 v_3, \ldots, v_{n-1} v_n, v_n v_1\}$.
    Chaque sommet est de degré~$2$.
  \item Le \textbf{graphe chemin} $P_n$ ($n \geq 1$) a pour sommets
    $\{v_1, \ldots, v_n\}$ et pour arêtes
    $\{v_1 v_2, v_2 v_3, \ldots, v_{n-1} v_n\}$.
    Les sommets extrêmes sont de degré~$1$, les autres de degré~$2$.
\end{itemize}
\end{definition}

\begin{center}
\begin{tikzpicture}[
  every node/.style={draw, circle, fill=green!20, minimum size=6mm,
    inner sep=0pt, font=\footnotesize},
  every edge/.style={draw, thick}
]
  % C_6
  \begin{scope}[shift={(-3,0)}]
    \node[label=below:{$C_6$}] at (0,-1.8) [draw=none, fill=none] {};
    \foreach \i in {1,...,6} {
      \node (c\i) at ({90 + (\i-1)*60}:1.3) {\i};
    }
    \draw (c1)--(c2)--(c3)--(c4)--(c5)--(c6)--(c1);
  \end{scope}

  % P_5
  \begin{scope}[shift={(3,0)}]
    \node[label=below:{$P_5$}] at (1.6,-1.8) [draw=none, fill=none] {};
    \foreach \i in {1,...,5} {
      \node (p\i) at ({(\i-1)*0.8}, 0) {\i};
    }
    \draw (p1)--(p2)--(p3)--(p4)--(p5);
  \end{scope}
\end{tikzpicture}
\captionof{figure}{Le graphe cycle $C_6$ et le graphe chemin $P_5$.}
\label{fig:cycles-chemins}
\end{center}

%% ──────────────────────────────────────────────────────────────────
\section{Sous-graphes}
\label{sec:sous-graphes}
\index{sous-graphe}
%% ──────────────────────────────────────────────────────────────────

\begin{definition}[Sous-graphe]
\label{def:sous-graphe}
Un graphe $H = (V', E')$ est un \textbf{sous-graphe} de $G = (V, E)$ si
$V' \subseteq V$ et $E' \subseteq E$.
\end{definition}

\begin{definition}[Sous-graphe induit]
\label{def:sous-graphe-induit}
\index{sous-graphe!induit}
Le \textbf{sous-graphe induit} par un sous-ensemble $S \subseteq V$, noté
$G[S]$, est le graphe $(S, E_S)$ où
$E_S = \set{\{u,v\} \in E : u \in S \text{ et } v \in S}$.
Autrement dit, on garde toutes les arêtes de~$G$ dont les deux extrémités
sont dans~$S$.
\end{definition}

\begin{definition}[Sous-graphe couvrant]
\label{def:sous-graphe-couvrant}
\index{sous-graphe!couvrant}
Un sous-graphe $H = (V, E')$ de $G = (V, E)$ est dit \textbf{couvrant}
(ou \emph{engendrant}) si $V' = V$, c'est-à-dire s'il contient tous les
sommets de~$G$.
\end{definition}

%% ──────────────────────────────────────────────────────────────────
\section{Le graphe de Petersen}
\label{sec:petersen}
\index{Petersen (graphe de)}
%% ──────────────────────────────────────────────────────────────────

Le \textbf{graphe de Petersen} est l'un des graphes les plus importants en
théorie des graphes. Il sert fréquemment de contre-exemple et possède de
nombreuses propriétés remarquables.

\begin{definition}[Graphe de Petersen]
\label{def:petersen}
Le graphe de Petersen a $10$~sommets et $15$~arêtes. Ses sommets sont les
paires $\{i,j\}$ avec $1 \leq i < j \leq 5$, et deux sommets sont adjacents
si et seulement si les paires correspondantes sont disjointes.
\end{definition}

\begin{center}
\begin{tikzpicture}[
  every node/.style={draw, circle, fill=orange!20, minimum size=6mm,
    inner sep=0pt, font=\scriptsize},
  every edge/.style={draw, thick}
]
  % Pentagone extérieur
  \foreach \i in {1,...,5} {
    \node (o\i) at ({90 + (\i-1)*72}:2.5) {\i};
  }
  \draw (o1)--(o2)--(o3)--(o4)--(o5)--(o1);

  % Pentagramme intérieur
  \foreach \i in {1,...,5} {
    \node (i\i) at ({90 + (\i-1)*72}:1.2) {$\i'$};
  }
  \draw (i1)--(i3)--(i5)--(i2)--(i4)--(i1);

  % Rayons
  \foreach \i in {1,...,5} {
    \draw (o\i) -- (i\i);
  }
\end{tikzpicture}
\captionof{figure}{Le graphe de Petersen : pentagone extérieur, pentagramme
  intérieur, et cinq rayons. C'est un graphe $3$-régulier avec $10$~sommets
  et $15$~arêtes.}
\label{fig:petersen}
\end{center}

\begin{proposition}[Propriétés du graphe de Petersen]
\label{prop:petersen-props}
Le graphe de Petersen satisfait les propriétés suivantes :
\begin{enumerate}[label=(\roman*)]
  \item il est $3$-régulier (chaque sommet a degré~$3$) ;
  \item il est connexe ;
  \item son diamètre est~$2$ (la distance maximale entre deux sommets) ;
  \item sa maille est~$5$ (longueur du plus court cycle) ;
  \item il n'est pas planaire ;
  \item il n'est pas hamiltonien (il ne possède pas de cycle passant par
    tous les sommets exactement une fois).
\end{enumerate}
\end{proposition}

%% ──────────────────────────────────────────────────────────────────
\section{Quelques illustrations supplémentaires}
\label{sec:illustrations}
%% ──────────────────────────────────────────────────────────────────

\begin{exemple}[Un graphe et ses représentations]
\label{ex:graphe-representations}
Considérons le graphe $G$ suivant :

\begin{center}
\begin{tikzpicture}[
  every node/.style={draw, circle, fill=cyan!20, minimum size=7mm,
    inner sep=0pt, font=\small\bfseries},
  every edge/.style={draw, thick}
]
  \node (a) at (0, 2)   {$a$};
  \node (b) at (2, 2)   {$b$};
  \node (c) at (3, 0.5) {$c$};
  \node (d) at (1, -0.5){$d$};
  \node (e) at (-1, 0.5){$e$};

  \draw (a) -- (b);
  \draw (a) -- (e);
  \draw (b) -- (c);
  \draw (b) -- (d);
  \draw (c) -- (d);
  \draw (d) -- (e);
\end{tikzpicture}
\end{center}

\textbf{Matrice d'adjacence} (avec l'ordre $a, b, c, d, e$) :
\[
  A = \begin{pmatrix}
    0 & 1 & 0 & 0 & 1 \\
    1 & 0 & 1 & 1 & 0 \\
    0 & 1 & 0 & 1 & 0 \\
    0 & 1 & 1 & 0 & 1 \\
    1 & 0 & 0 & 1 & 0
  \end{pmatrix}.
\]

\textbf{Listes d'adjacence} :
\begin{center}
\begin{tabular}{@{}cl@{}}
\toprule
\textbf{Sommet} & \textbf{Voisins} \\
\midrule
$a$ & $b, e$ \\
$b$ & $a, c, d$ \\
$c$ & $b, d$ \\
$d$ & $b, c, e$ \\
$e$ & $a, d$ \\
\bottomrule
\end{tabular}
\end{center}

\textbf{Degrés} : $\deg(a) = 2$, $\deg(b) = 3$, $\deg(c) = 2$, $\deg(d) = 3$,
$\deg(e) = 2$. Somme : $2+3+2+3+2 = 12 = 2 \times 6 = 2m$. \checkmark
\end{exemple}

\begin{exemple}[Graphe orienté]
\label{ex:digraphe}
\begin{center}
\begin{tikzpicture}[
  every node/.style={draw, circle, fill=violet!15, minimum size=7mm,
    inner sep=0pt, font=\small\bfseries},
  >={Stealth[length=3mm]},
  every edge/.style={draw, thick, ->}
]
  \node (1) at (0, 1.5)  {$1$};
  \node (2) at (2, 1.5)  {$2$};
  \node (3) at (3, 0)    {$3$};
  \node (4) at (1, -0.5) {$4$};

  \draw[->] (1) -- (2);
  \draw[->] (2) -- (3);
  \draw[->] (3) -- (4);
  \draw[->] (4) -- (1);
  \draw[->] (1) -- (3);
\end{tikzpicture}
\end{center}

Ce digraphe a $4$~sommets et $5$~arcs. Les degrés sont :
\begin{center}
\begin{tabular}{@{}cccc@{}}
\toprule
\textbf{Sommet} & $\deg^+$ & $\deg^-$ & $\deg^+ + \deg^-$ \\
\midrule
$1$ & $2$ & $1$ & $3$ \\
$2$ & $1$ & $1$ & $2$ \\
$3$ & $1$ & $2$ & $3$ \\
$4$ & $1$ & $1$ & $2$ \\
\bottomrule
\end{tabular}
\end{center}
On vérifie : $\sum \deg^+ = \sum \deg^- = 5 = \abs{A}$.
\end{exemple}

\begin{exemple}[Graphe pondéré]
\label{ex:graphe-pondere}
\begin{center}
\begin{tikzpicture}[
  every node/.style={draw, circle, fill=yellow!20, minimum size=7mm,
    inner sep=0pt, font=\small\bfseries},
  every edge/.style={draw, thick}
]
  \node (a) at (0, 1.5)  {$a$};
  \node (b) at (2.5, 1.5){$b$};
  \node (c) at (3.5, 0)  {$c$};
  \node (d) at (1.5, -1) {$d$};
  \node (e) at (-0.5, 0) {$e$};

  \draw (a) -- node[draw=none, fill=none, above, font=\small] {3} (b);
  \draw (b) -- node[draw=none, fill=none, right, font=\small] {7} (c);
  \draw (c) -- node[draw=none, fill=none, below, font=\small] {2} (d);
  \draw (d) -- node[draw=none, fill=none, below, font=\small] {4} (e);
  \draw (e) -- node[draw=none, fill=none, left, font=\small]  {1} (a);
  \draw (a) -- node[draw=none, fill=none, above right, font=\small] {5} (d);
\end{tikzpicture}
\end{center}
\end{exemple}

%% ──────────────────────────────────────────────────────────────────
\section{Exercices}
\label{sec:graphes-exercices}
%% ──────────────────────────────────────────────────────────────────

\begin{exercice}[Dénombrement d'arêtes]
\label{exo:aretes}
\begin{enumerate}[label=(\alph*)]
  \item Combien d'arêtes possède le graphe complet $K_n$ ?
  \item Combien d'arêtes possède le graphe biparti complet $K_{m,n}$ ?
  \item Montrer qu'un graphe simple d'ordre~$n$ possède au plus
    $\binom{n}{2}$ arêtes.
\end{enumerate}
\end{exercice}

\begin{exercice}[Suites de degrés]
\label{exo:suites-degres}
\index{suite de degres@suite de degrés}
Parmi les suites suivantes, lesquelles sont des suites de degrés d'un graphe
simple ? Pour celles qui le sont, construire un tel graphe.
\begin{enumerate}[label=(\alph*)]
  \item $(3, 3, 3, 3)$
  \item $(1, 1, 2, 3, 3)$
  \item $(0, 1, 2, 3)$
  \item $(2, 2, 2, 2, 2)$
  \item $(1, 2, 3, 4, 5)$
\end{enumerate}
\textit{Indication :} utiliser le théorème d'Erd\H{o}s--Gallai\index{Erdos-Gallai@Erd\H{o}s--Gallai (théorème)} ou la condition de la somme paire.
\end{exercice}

\begin{exercice}[Matrice d'adjacence et chemins]
\label{exo:matrice-chemins}
Soit $A$ la matrice d'adjacence d'un graphe $G$.
\begin{enumerate}[label=(\alph*)]
  \item Montrer que $(A^k)_{ij}$ est le nombre de marches de longueur~$k$
    entre $v_i$ et $v_j$.
  \item En déduire que $\mathrm{tr}(A^2) = 2m$ et que $\mathrm{tr}(A^3)/6$
    est le nombre de triangles dans~$G$.
\end{enumerate}
\end{exercice}

\begin{exercice}[Graphes bipartis]
\label{exo:biparti}
\begin{enumerate}[label=(\alph*)]
  \item Le graphe cycle $C_n$ est-il biparti ? Distinguer les cas $n$ pair
    et $n$ impair.
  \item Montrer que tout arbre est biparti.
  \item Montrer qu'un graphe biparti d'ordre~$n$ possède au plus
    $\lfloor n^2/4 \rfloor$ arêtes.
\end{enumerate}
\end{exercice}

\begin{exercice}[Graphe de Petersen]
\label{exo:petersen}
\begin{enumerate}[label=(\alph*)]
  \item Vérifier que le graphe de Petersen a bien $10$~sommets et $15$~arêtes.
  \item Montrer que le graphe de Petersen est $3$-régulier.
  \item Montrer que le diamètre du graphe de Petersen est~$2$.
  \item Trouver un cycle de longueur~$5$ dans le graphe de Petersen et montrer
    qu'il n'en existe pas de plus court.
\end{enumerate}
\end{exercice}

\begin{exercice}[Connexité et suppression d'arêtes]
\label{exo:connexite}
\begin{enumerate}[label=(\alph*)]
  \item Montrer qu'un graphe connexe d'ordre $n$ et de taille $n - 1$ est
    un arbre.
  \item Montrer que si $G$ est connexe et $e$ est une arête appartenant à un
    cycle de $G$, alors $G - e$ est encore connexe.
  \item Une arête $e$ d'un graphe connexe $G$ est un \textbf{pont}\index{pont}
    si $G - e$ n'est pas connexe. Montrer qu'une arête est un pont si et
    seulement si elle n'appartient à aucun cycle.
\end{enumerate}
\end{exercice}

\begin{exercice}[Isomorphisme]
\label{exo:isomorphisme}
Les paires de graphes suivantes sont-elles isomorphes ? Justifier.
\begin{enumerate}[label=(\alph*)]
  \item $C_6$ et $K_{3,3}$.
  \item Le graphe cycle $C_5$ et le graphe complet $K_5$ privé d'un couplage
    parfait \emph{(impossible car $K_5$ a un nombre impair de sommets ---
    interpréter la question comme $K_5$ privé de deux arêtes non adjacentes)}.
  \item Le graphe cube $Q_3$ (sommets = mots binaires de longueur~3, arêtes
    entre mots différant d'un bit) et le graphe biparti complet $K_{3,3}$ plus
    deux arêtes indépendantes.
\end{enumerate}
\end{exercice}

\begin{exercice}[Sous-graphes induits]
\label{exo:sous-graphe-induit}
Soit $G = K_5$ et $S = \{1, 2, 3\}$.
\begin{enumerate}[label=(\alph*)]
  \item Décrire le sous-graphe induit $G[S]$.
  \item Combien de sous-graphes induits (à isomorphisme près) le graphe $K_4$
    possède-t-il ?
  \item Montrer que tout graphe d'ordre $n \geq 3$ contient un sous-graphe
    induit d'ordre~$3$ qui est soit $K_3$, soit le graphe vide sur $3$~sommets
    (théorème de Ramsey pour $R(3,3)$).
\end{enumerate}
\end{exercice}

\begin{exercice}[Complément d'un graphe]
\label{exo:complement}
\index{complement d'un graphe@complément d'un graphe}
Le \textbf{complément} d'un graphe $G = (V, E)$ est le graphe
$\overline{G} = (V, \binom{V}{2} \setminus E)$.
\begin{enumerate}[label=(\alph*)]
  \item Montrer que $\abs{E(G)} + \abs{E(\overline{G})} = \binom{n}{2}$.
  \item Montrer que $G$ ou $\overline{G}$ est connexe.
  \item Quel est le complément de $C_5$ ?
\end{enumerate}
\end{exercice}

%% ──────────────────────────────────────────────────────────────────
\section*{Résumé du chapitre}
\addcontentsline{toc}{section}{Résumé du chapitre}
%% ──────────────────────────────────────────────────────────────────

\begin{itemize}
  \item Un \textbf{graphe simple} $G = (V, E)$ est un ensemble de sommets~$V$
    relié par des arêtes~$E \subseteq \binom{V}{2}$. Les variantes incluent les
    multigraphes, les digraphes et les graphes pondérés.
  \item Les principales \textbf{représentations} sont la matrice d'adjacence
    ($\Theta(n^2)$ en espace) et les listes d'adjacence ($\Theta(n+m)$).
  \item Le \textbf{lemme des poignées de main} affirme que
    $\sum_{v} \deg(v) = 2m$, ce qui implique que le nombre de sommets de degré
    impair est pair.
  \item Un graphe est \textbf{connexe} si toute paire de sommets est reliée
    par un chemin. Les classes d'équivalence de la relation de connexité
    forment les \textbf{composantes connexes}.
  \item Deux graphes sont \textbf{isomorphes} s'il existe une bijection entre
    leurs sommets préservant l'adjacence.
  \item Les \textbf{familles classiques} incluent le graphe complet~$K_n$
    ($\binom{n}{2}$~arêtes), le biparti complet~$K_{m,n}$ ($mn$~arêtes),
    le cycle~$C_n$ et le chemin~$P_n$.
  \item Un graphe est \textbf{biparti} si et seulement s'il ne contient aucun
    cycle de longueur impaire.
  \item Le \textbf{graphe de Petersen} est un graphe $3$-régulier à
    $10$~sommets et $15$~arêtes, possédant de nombreuses propriétés
    remarquables et servant souvent de contre-exemple.
\end{itemize}

% === FIN DU CHUNK ch7-8 ===
% === DÉBUT DU CHUNK ch9-11+end ===

%%%%%%%%%%%%%%%%%%%%%%%%%%%%%%%%%%%%%%%%%%%%%%%%%%%%%%%%%%%%%%%%%%%%%
%%  CHAPITRE 9 — Arbres, Graphes Eulériens et Hamiltoniens
%%%%%%%%%%%%%%%%%%%%%%%%%%%%%%%%%%%%%%%%%%%%%%%%%%%%%%%%%%%%%%%%%%%%%
\chapter{Arbres, Graphes Eulériens et Hamiltoniens}
\label{ch:arbres-euler-hamilton}
\index{arbre}\index{graphe eulérien}\index{graphe hamiltonien}

\section{Arbres}
\label{sec:arbres}

\begin{definition}[Arbre]
\label{def:arbre}
\index{arbre!définition}
Un \emph{arbre} est un graphe connexe et acyclique. Un \emph{forêt} est
un graphe acyclique (dont chaque composante connexe est un arbre).
\end{definition}

\begin{remarque}
Un arbre à $n$ sommets possède exactement $n-1$ arêtes. La réciproque
est également vraie sous des hypothèses adéquates, comme le montre le
théorème de caractérisation suivant.
\end{remarque}

\begin{theoreme}[Caractérisation des arbres]
\label{thm:caract-arbre}
\index{arbre!caractérisation}
Soit $G$ un graphe à $n \geq 1$ sommets. Les assertions suivantes sont
équivalentes :
\begin{enumerate}[label=(\roman*)]
  \item $G$ est un arbre (connexe et acyclique).
  \item $G$ est connexe et possède exactement $n-1$ arêtes.
  \item $G$ est acyclique et possède exactement $n-1$ arêtes.
  \item Pour tous sommets $u,v$ de $G$, il existe un unique chemin de $u$
        à $v$.
  \item $G$ est connexe et la suppression de toute arête le déconnecte.
  \item $G$ est acyclique et l'ajout de toute arête crée exactement un
        cycle.
\end{enumerate}
\end{theoreme}

\begin{proof}
Nous démontrons les implications principales.

\medskip\noindent
\textbf{(i) $\Rightarrow$ (ii).}\;
Soit $G$ un arbre à $n$ sommets. Procédons par récurrence sur $n$.
Pour $n=1$, le graphe possède $0 = n-1$ arêtes. Supposons le résultat
vrai pour tout arbre à $n-1$ sommets. Puisque $G$ est un arbre à $n\geq 2$
sommets, il possède au moins un sommet pendant $v$ (de degré~$1$), car
sinon tout sommet aurait degré $\geq 2$, ce qui impliquerait l'existence
d'un cycle (contradiction avec l'acyclicité). Le graphe $G - v$ est encore
connexe et acyclique, donc un arbre à $n-1$ sommets ; par hypothèse de
récurrence, il possède $n-2$ arêtes. Ainsi $G$ possède $n-2+1 = n-1$ arêtes.

\medskip\noindent
\textbf{(ii) $\Rightarrow$ (iii).}\;
Supposons $G$ connexe avec $n-1$ arêtes. S'il contenait un cycle $C$, la
suppression d'une arête de $C$ conserverait la connexité, produisant un
graphe connexe à $n-1-1 = n-2$ arêtes. On pourrait répéter jusqu'à obtenir
un graphe connexe acyclique (un arbre) à $n$ sommets avec strictement moins
de $n-1$ arêtes. Mais un arbre à $n$ sommets a $n-1$ arêtes (par (i)
$\Rightarrow$ (ii)), contradiction.

\medskip\noindent
\textbf{(iii) $\Rightarrow$ (i).}\;
Soit $G$ acyclique avec $n-1$ arêtes. Si $G$ a $k$ composantes connexes,
celles-ci sont des arbres de tailles $n_1,\dots,n_k$ avec $\sum n_i = n$.
Chaque composante possède $n_i - 1$ arêtes, d'où un total de $n - k$
arêtes. Comme ce total vaut $n-1$, on obtient $k=1$ et $G$ est connexe.

\medskip\noindent
\textbf{(i) $\Rightarrow$ (iv).}\;
La connexité assure l'existence d'un chemin entre toute paire de sommets.
S'il existait deux chemins distincts de $u$ à $v$, leur concaténation
contiendrait un cycle, contredisant l'acyclicité.

\medskip
Les autres implications se démontrent de manière analogue.
\end{proof}

\begin{exemple}[Quelques arbres]
\label{ex:arbres-petits}

\begin{center}
\begin{tikzpicture}[
  every node/.style={circle, draw, fill=blue!15, inner sep=2pt,
                     minimum size=16pt, font=\footnotesize},
  level distance=12mm, sibling distance=15mm
]
% Arbre chemin P_4
\begin{scope}[xshift=0cm]
  \node (a1) {1}
    child { node (a2) {2}
      child { node (a3) {3}
        child { node (a4) {4} }
      }
    };
  \node[draw=none, fill=none] at (0,-4.2) {$P_4$};
\end{scope}

% Étoile S_4
\begin{scope}[xshift=4.5cm]
  \node (s0) {0}
    child { node {1} }
    child { node {2} }
    child { node {3} }
    child { node {4} };
  \node[draw=none, fill=none] at (0,-2.7) {$S_4$};
\end{scope}

% Arbre binaire complet
\begin{scope}[xshift=9.5cm]
  \node (r) {$r$}
    child { node {$a$}
      child { node {$c$} }
      child { node {$d$} }
    }
    child { node {$b$}
      child { node {$e$} }
      child { node {$f$} }
    };
  \node[draw=none, fill=none] at (0,-4.2) {Arbre binaire complet};
\end{scope}
\end{tikzpicture}
\end{center}
\end{exemple}

\section{Arbres couvrants}
\label{sec:arbres-couvrants}
\index{arbre couvrant}

\begin{definition}[Arbre couvrant]
\label{def:arbre-couvrant}
Soit $G=(V,E)$ un graphe connexe. Un \emph{arbre couvrant} de $G$ est un
sous-graphe couvrant $T=(V,E')$ avec $E' \subseteq E$ qui est un arbre.
\end{definition}

\begin{proposition}
Tout graphe connexe possède au moins un arbre couvrant.
\end{proposition}

\begin{proof}
Soit $G$ connexe. Si $G$ est acyclique, c'est un arbre et il est son
propre arbre couvrant. Sinon, $G$ contient un cycle ; on supprime une
arête du cycle, ce qui préserve la connexité. On itère jusqu'à obtenir
un sous-graphe connexe acyclique couvrant, c'est-à-dire un arbre
couvrant.
\end{proof}

\begin{theoreme}[Formule de Cayley]
\label{thm:cayley}
\index{formule de Cayley}
Le nombre d'arbres couvrants étiquetés du graphe complet $K_n$ est
$n^{n-2}$.
\end{theoreme}

\begin{exemple}
Pour $K_3$, la formule donne $3^1 = 3$ arbres étiquetés, ce qui correspond
aux trois chemins possibles parmi les sommets $\{1,2,3\}$. Pour $K_4$, on
obtient $4^2 = 16$ arbres couvrants étiquetés.
\end{exemple}

\section{Arbres couvrants de poids minimum}
\label{sec:acm}
\index{arbre couvrant!de poids minimum}

\begin{definition}[Graphe pondéré et arbre couvrant de poids minimum]
\label{def:acm}
Un \emph{graphe pondéré} est un graphe $G=(V,E)$ muni d'une fonction de
poids $w : E \to \R$. Un \emph{arbre couvrant de poids minimum} (ACM)
de $G$ est un arbre couvrant dont la somme des poids des arêtes est
minimale.
\end{definition}

\begin{theoreme}[Algorithme de Kruskal]
\label{thm:kruskal}
\index{algorithme de Kruskal}
L'algorithme suivant produit un ACM d'un graphe connexe pondéré
$G = (V,E,w)$ à $n$ sommets :
\begin{enumerate}
  \item Trier les arêtes par poids croissant : $w(e_1) \leq w(e_2) \leq
        \cdots \leq w(e_m)$.
  \item Initialiser $T \leftarrow (V, \varnothing)$.
  \item Pour $i = 1, 2, \dots, m$ : si l'ajout de $e_i$ à $T$ ne crée
        pas de cycle, ajouter $e_i$ à $T$.
  \item Retourner $T$.
\end{enumerate}
L'algorithme se termine lorsque $T$ possède $n-1$ arêtes.
\end{theoreme}

\begin{theoreme}[Algorithme de Prim]
\label{thm:prim}
\index{algorithme de Prim}
L'algorithme suivant produit un ACM d'un graphe connexe pondéré
$G = (V,E,w)$ :
\begin{enumerate}
  \item Choisir un sommet initial $s \in V$. Poser $S \leftarrow \{s\}$,
        $T \leftarrow (S, \varnothing)$.
  \item Tant que $S \neq V$ : parmi les arêtes $\{u,v\}$ avec $u \in S$
        et $v \notin S$, choisir celle de poids minimal ; ajouter $v$ à
        $S$ et l'arête à $T$.
  \item Retourner $T$.
\end{enumerate}
\end{theoreme}

\begin{exemple}[Arbre couvrant de poids minimum]
\label{ex:acm}

\begin{center}
\begin{tikzpicture}[
  every node/.style={circle, draw, fill=blue!15, inner sep=2pt,
                     minimum size=18pt, font=\footnotesize},
  edge label/.style={draw=none, fill=none, font=\scriptsize,
                     inner sep=1pt}
]
  \node (a) at (0,2)   {$a$};
  \node (b) at (2,3)   {$b$};
  \node (c) at (4,2)   {$c$};
  \node (d) at (0,0)   {$d$};
  \node (e) at (2,0.8) {$e$};
  \node (f) at (4,0)   {$f$};

  % Arêtes de l'ACM (en gras)
  \draw[very thick, red!70!black] (a) -- node[edge label, above left] {1} (b);
  \draw[very thick, red!70!black] (b) -- node[edge label, above right] {2} (c);
  \draw[very thick, red!70!black] (a) -- node[edge label, left] {3} (d);
  \draw[very thick, red!70!black] (d) -- node[edge label, below] {2} (e);
  \draw[very thick, red!70!black] (e) -- node[edge label, below right] {1} (f);

  % Arêtes hors ACM
  \draw[thin, gray, dashed] (b) -- node[edge label, right] {5} (e);
  \draw[thin, gray, dashed] (c) -- node[edge label, right] {4} (f);
  \draw[thin, gray, dashed] (a) -- node[edge label, below left] {6} (e);
  \draw[thin, gray, dashed] (d) -- node[edge label, below] {7} (f);
\end{tikzpicture}
\end{center}

\noindent
Les arêtes en trait épais forment un ACM de poids total $1+2+3+2+1=9$.
\end{exemple}

\section{Circuits et chemins d'Euler}
\label{sec:euler}
\index{circuit eulérien}\index{chemin eulérien}

\begin{definition}[Chemin et circuit d'Euler]
\label{def:euler}
Soit $G$ un graphe (simple ou multi\-graphe). Un \emph{chemin d'Euler}
est un chemin qui emprunte chaque arête exactement une fois. Un
\emph{circuit d'Euler} est un chemin d'Euler fermé (le sommet de départ
coïncide avec le sommet d'arrivée). Un graphe possédant un circuit
d'Euler est dit \emph{eulérien}.
\end{definition}

\begin{theoreme}[Condition nécessaire et suffisante pour un circuit d'Euler]
\label{thm:euler-circuit}
\index{circuit eulérien!condition}
Un multigraphe connexe $G$ possède un circuit d'Euler si et seulement si
chaque sommet de $G$ est de degré pair.
\end{theoreme}

\begin{proof}
\textbf{Condition nécessaire.}\;
Supposons que $G$ possède un circuit d'Euler $C$. Chaque fois que $C$
traverse un sommet $v$ (distinct du sommet de départ/arrivée), il y entre
par une arête et en sort par une autre, contribuant $2$ au degré de $v$.
Le sommet de départ/arrivée voit aussi son degré augmenté de $2$ à chaque
passage. Comme chaque arête est parcourue exactement une fois, le degré
de chaque sommet est pair.

\medskip\noindent
\textbf{Condition suffisante.}\;
Supposons que tout sommet de $G$ est de degré pair. On procède par
récurrence forte sur le nombre d'arêtes $m$.

Si $m=0$, le circuit trivial convient. Supposons $m \geq 1$. Puisque
chaque sommet a degré pair $\geq 2$ (le graphe étant connexe avec $m\geq 1$),
on peut construire une marche fermée $W$ partant d'un sommet quelconque $v_0$
en suivant des arêtes non encore visitées, sans jamais rester bloqué sur
un sommet différent de $v_0$ (car les degrés pairs garantissent une arête
de sortie tant que l'on n'est pas revenu à $v_0$).

Si $W$ parcourt toutes les arêtes, c'est un circuit d'Euler. Sinon, le
graphe $G' = G - E(W)$ a tous ses sommets de degré pair. Chaque
composante connexe $G'_i$ de $G'$ ayant un sommet en commun avec $W$
(car $G$ est connexe) possède, par hypothèse de récurrence, un circuit
d'Euler $C_i$. On insère chaque $C_i$ dans $W$ au sommet commun pour
obtenir un circuit d'Euler de $G$.
\end{proof}

\begin{corollaire}[Condition pour un chemin d'Euler]
\label{cor:euler-chemin}
Un multigraphe connexe $G$ possède un chemin d'Euler (non fermé) si et
seulement si $G$ possède exactement deux sommets de degré impair. Dans ce
cas, le chemin relie ces deux sommets.
\end{corollaire}

\begin{theoreme}[Algorithme de Fleury]
\label{thm:fleury}
\index{algorithme de Fleury}
Soit $G$ un multigraphe connexe dont tous les sommets sont de degré pair.
L'algorithme suivant produit un circuit d'Euler :
\begin{enumerate}
  \item Partir d'un sommet quelconque $v$.
  \item À chaque étape, choisir une arête incidente au sommet courant, en
        évitant les \emph{ponts} (arêtes dont la suppression déconnecte le
        graphe restant) sauf s'il n'y a pas d'autre choix.
  \item Supprimer l'arête choisie et passer au sommet suivant.
  \item Répéter jusqu'à avoir parcouru toutes les arêtes.
\end{enumerate}
\end{theoreme}

\begin{exemple}[Circuit d'Euler]
\label{ex:circuit-euler}

\begin{center}
\begin{tikzpicture}[
  every node/.style={circle, draw, fill=green!15, inner sep=2pt,
                     minimum size=18pt, font=\footnotesize},
  decoration={markings, mark=at position 0.55 with {\arrow{>}}},
]
  \node (a) at (0,0)   {$a$};
  \node (b) at (2,0)   {$b$};
  \node (c) at (3,1.5) {$c$};
  \node (d) at (2,3)   {$d$};
  \node (e) at (0,3)   {$e$};
  \node (f) at (-1,1.5){$f$};

  \draw[postaction={decorate}, thick, blue!70!black]
    (a) -- (b) -- (c) -- (d) -- (e) -- (f) -- (a);
  \draw[postaction={decorate}, thick, blue!70!black]
    (a) -- (d);
  \draw[postaction={decorate}, thick, blue!70!black]
    (b) -- (e);
  \draw[postaction={decorate}, thick, blue!70!black]
    (c) -- (f);
\end{tikzpicture}
\end{center}

\noindent
Ce graphe est eulérien : chaque sommet a degré~$4$. Un circuit d'Euler
possible est $a \to b \to e \to d \to c \to f \to a \to d \to \cdots$
(en parcourant toutes les arêtes exactement une fois).
\end{exemple}

\section{Circuits et chemins hamiltoniens}
\label{sec:hamilton}
\index{circuit hamiltonien}\index{chemin hamiltonien}

\begin{definition}[Chemin et circuit hamiltonien]
\label{def:hamilton}
Un \emph{chemin hamiltonien} d'un graphe $G$ est un chemin qui visite
chaque sommet exactement une fois. Un \emph{circuit hamiltonien} est un
cycle qui visite chaque sommet exactement une fois. Un graphe possédant
un circuit hamiltonien est dit \emph{hamiltonien}.
\end{definition}

\begin{remarque}
Contrairement au problème eulérien, il n'existe pas de caractérisation
simple des graphes hamiltoniens. Déterminer si un graphe est hamiltonien
est un problème NP-complet.
\end{remarque}

\begin{theoreme}[Théorème de Dirac]
\label{thm:dirac}
\index{théorème de Dirac}
Soit $G$ un graphe simple à $n \geq 3$ sommets. Si $\deg(v) \geq n/2$
pour tout sommet $v$, alors $G$ est hamiltonien.
\end{theoreme}

\begin{proof}
Supposons par l'absurde que $G$ satisfait l'hypothèse mais n'est pas
hamiltonien. Ajoutons des arêtes à $G$ de sorte que le graphe résultant
$G^*$ ne soit toujours pas hamiltonien, mais que l'ajout de toute arête
supplémentaire crée un circuit hamiltonien. Un tel graphe maximal $G^*$
existe.

Soient $u$ et $v$ deux sommets non adjacents dans $G^*$. L'ajout de
$\{u,v\}$ à $G^*$ crée un circuit hamiltonien, donc $G^*$ possède un
chemin hamiltonien $v_1 = u, v_2, \dots, v_n = v$.

Considérons les ensembles :
\[
  S = \{i : \{u, v_{i+1}\} \in E(G^*)\}, \quad
  T = \{i : \{v, v_i\} \in E(G^*)\}.
\]
On a $S,T \subseteq \{1,\dots,n-1\}$, $\card{S} = \deg_{G^*}(u) \geq n/2$
et $\card{T} = \deg_{G^*}(v) \geq n/2$. Par le principe des tiroirs,
$S \cap T \neq \varnothing$, c'est-à-dire il existe un indice $i$ tel que
$\{u, v_{i+1}\} \in E(G^*)$ et $\{v, v_i\} \in E(G^*)$.

On obtient alors le circuit hamiltonien :
\[
  u = v_1, v_2, \dots, v_i, v, v_{n-1}, v_{n-2}, \dots, v_{i+1}, u,
\]
ce qui contredit le choix de $G^*$.
\end{proof}

\begin{theoreme}[Théorème d'Ore]
\label{thm:ore}
\index{théorème d'Ore}
Soit $G$ un graphe simple à $n \geq 3$ sommets. Si pour toute paire de
sommets non adjacents $u,v$ on a $\deg(u) + \deg(v) \geq n$, alors $G$
est hamiltonien.
\end{theoreme}

\begin{proof}
La preuve est identique à celle du théorème de Dirac, en remplaçant la
condition $\deg(v) \geq n/2$ par $\deg(u)+\deg(v) \geq n$ pour la paire
non adjacente $\{u,v\}$. On obtient de même $\card{S} + \card{T} \geq n$
avec $S, T \subseteq \{1,\dots,n-1\}$, d'où $S \cap T \neq \varnothing$,
et la même construction de circuit hamiltonien mène à la contradiction.
\end{proof}

\begin{corollaire}
Le théorème de Dirac est un cas particulier du théorème d'Ore.
\end{corollaire}

\begin{exemple}[Graphe hamiltonien vs non hamiltonien]
\label{ex:hamilton}

\begin{center}
\begin{tikzpicture}[
  every node/.style={circle, draw, fill=orange!15, inner sep=2pt,
                     minimum size=18pt, font=\footnotesize}
]
% Graphe hamiltonien (Petersen extérieur)
\begin{scope}[xshift=0cm]
  \foreach \i in {1,...,5} {
    \node (v\i) at ({90+72*(\i-1)}:1.5) {$v_\i$};
  }
  \draw[thick] (v1)--(v2)--(v3)--(v4)--(v5)--(v1);
  \draw[thick, red, dashed] (v1)--(v3)--(v5)--(v2)--(v4)--(v1);
  \node[draw=none, fill=none] at (0,-2.3) {$C_5$ : hamiltonien};
\end{scope}

% Graphe de Petersen (non hamiltonien)
\begin{scope}[xshift=6cm]
  \foreach \i in {1,...,5} {
    \node (a\i) at ({90+72*(\i-1)}:2) {};
    \node (b\i) at ({90+72*(\i-1)}:1) {};
  }
  \draw (a1)--(a2)--(a3)--(a4)--(a5)--(a1);
  \draw (b1)--(b3)--(b5)--(b2)--(b4)--(b1);
  \draw (a1)--(b1) (a2)--(b2) (a3)--(b3) (a4)--(b4) (a5)--(b5);
  \node[draw=none, fill=none] at (0,-2.8) {Graphe de Petersen};
  \node[draw=none, fill=none] at (0,-3.4) {(non hamiltonien)};
\end{scope}
\end{tikzpicture}
\end{center}
\end{exemple}

\section{Exercices}
\label{sec:exo-ch9}

\begin{exercice}
Montrer qu'un arbre à $n \geq 2$ sommets possède au moins deux sommets
pendants (de degré~$1$).
\end{exercice}

\begin{exercice}
Déterminer le nombre d'arbres couvrants de $K_4$ à l'aide de la formule de
Cayley, puis les lister explicitement.
\end{exercice}

\begin{exercice}
Appliquer l'algorithme de Kruskal au graphe complet $K_4$ avec les poids
$w(\{1,2\})=3$, $w(\{1,3\})=1$, $w(\{1,4\})=4$, $w(\{2,3\})=5$,
$w(\{2,4\})=2$, $w(\{3,4\})=6$.
\end{exercice}

\begin{exercice}
Le graphe de Königsberg (les sept ponts) admet-il un chemin d'Euler ?
Justifier en calculant les degrés des sommets.
\end{exercice}

\begin{exercice}
Soit $G$ un graphe simple à $n=10$ sommets tel que $\deg(v) \geq 5$ pour
tout sommet $v$. Montrer que $G$ est hamiltonien.
\end{exercice}

\begin{exercice}
Montrer que le graphe de Petersen n'est pas hamiltonien. \textit{Indication :
raisonner par l'absurde en considérant un cycle hamiltonien supposé et les
arêtes restantes.}
\end{exercice}

\begin{exercice}
Montrer que tout graphe connexe possède un parcours d'Euler (chemin ou
circuit) si et seulement s'il possède au plus deux sommets de degré impair.
\end{exercice}

\begin{exercice}
Soit $T$ un arbre à $n$ sommets dont les degrés sont $d_1, d_2, \dots, d_n$.
Montrer que $\sum_{i=1}^n d_i = 2(n-1)$.
\end{exercice}

\bigskip
\begin{center}
\fbox{\parbox{0.85\textwidth}{%
\textbf{Résumé du chapitre~9.}\;
Un arbre est un graphe connexe acyclique ; il possède $n-1$ arêtes et
chaque paire de sommets est reliée par un unique chemin. La formule de
Cayley dénombre les arbres étiquetés. Les algorithmes de Kruskal et Prim
construisent des arbres couvrants de poids minimum. Un graphe connexe est
eulérien si et seulement si tous ses sommets sont de degré pair. Les
théorèmes de Dirac et d'Ore fournissent des conditions suffisantes pour
l'existence de circuits hamiltoniens.
}}
\end{center}


%%%%%%%%%%%%%%%%%%%%%%%%%%%%%%%%%%%%%%%%%%%%%%%%%%%%%%%%%%%%%%%%%%%%%
%%  CHAPITRE 10 — Coloration, Planéité et Graphes Bipartis
%%%%%%%%%%%%%%%%%%%%%%%%%%%%%%%%%%%%%%%%%%%%%%%%%%%%%%%%%%%%%%%%%%%%%
\chapter{Coloration, Planéité et Graphes Bipartis}
\label{ch:coloration-planarite}
\index{coloration}\index{graphe planaire}\index{graphe biparti}

\section{Coloration de sommets}
\label{sec:coloration-sommets}

\begin{definition}[Coloration propre et nombre chromatique]
\label{def:coloration}
\index{coloration!propre}\index{nombre chromatique}
Soit $G=(V,E)$ un graphe. Une \emph{coloration propre} de $G$ en $k$
couleurs est une fonction $c : V \to \{1, 2, \dots, k\}$ telle que
$c(u) \neq c(v)$ pour toute arête $\{u,v\} \in E$. Le \emph{nombre
chromatique} $\chi(G)$ est le plus petit entier $k$ tel que $G$ admet
une coloration propre en $k$ couleurs.
\end{definition}

\begin{exemple}
$\chi(K_n) = n$, $\chi(C_{2k}) = 2$, $\chi(C_{2k+1}) = 3$ pour
$k \geq 1$.
\end{exemple}

\begin{proposition}
\label{prop:chi-borne}
Pour tout graphe $G$, $\chi(G) \leq \Delta(G) + 1$, où
$\Delta(G) = \max_{v \in V} \deg(v)$.
\end{proposition}

\begin{proof}
L'algorithme glouton suivant réalise cette borne.
\end{proof}

\subsection{Algorithme de coloration gloutonne}
\label{subsec:coloration-gloutonne}
\index{coloration!algorithme glouton}

\begin{theoreme}[Coloration gloutonne]
\label{thm:coloration-gloutonne}
Soit $G=(V,E)$ un graphe à $n$ sommets. L'algorithme suivant produit une
coloration propre en au plus $\Delta(G)+1$ couleurs :
\begin{enumerate}
  \item Numéroter les sommets $v_1, v_2, \dots, v_n$ dans un ordre
        quelconque.
  \item Pour $i = 1, 2, \dots, n$ : attribuer à $v_i$ la plus petite
        couleur non utilisée par ses voisins déjà colorés.
\end{enumerate}
\end{theoreme}

\begin{proof}
Lorsqu'on colore $v_i$, celui-ci a au plus $\deg(v_i) \leq \Delta(G)$
voisins déjà colorés, occupant au plus $\Delta(G)$ couleurs distinctes.
Il reste donc au moins une couleur disponible parmi
$\{1, 2, \dots, \Delta(G)+1\}$.
\end{proof}

\begin{theoreme}[Théorème de Brooks]
\label{thm:brooks}
\index{théorème de Brooks}
Soit $G$ un graphe connexe qui n'est ni un cycle impair ni un graphe
complet. Alors $\chi(G) \leq \Delta(G)$.
\end{theoreme}

\begin{theoreme}[Théorème des quatre couleurs]
\label{thm:quatre-couleurs}
\index{théorème des quatre couleurs}
Tout graphe planaire est $4$-colorable, c'est-à-dire $\chi(G) \leq 4$
pour tout graphe planaire $G$.
\end{theoreme}

\begin{remarque}
Le théorème des quatre couleurs, conjecturé en 1852 par Francis Guthrie,
a été démontré en 1976 par Appel et Haken à l'aide d'un programme
informatique vérifiant un nombre fini de cas. C'est l'un des premiers
théorèmes majeurs dont la preuve repose sur l'assistance d'un ordinateur.
\end{remarque}

\section{Polynôme chromatique}
\label{sec:polynome-chromatique}
\index{polynôme chromatique}

\begin{definition}[Polynôme chromatique]
\label{def:polynome-chromatique}
Le \emph{polynôme chromatique} $P(G, k)$ d'un graphe $G$ est le nombre
de colorations propres de $G$ en $k$ couleurs.
\end{definition}

\begin{proposition}[Exemples fondamentaux]
\label{prop:polynome-exemples}
\begin{enumerate}[label=(\roman*)]
  \item $P(K_n, k) = k(k-1)(k-2)\cdots(k-n+1) = k^{(n)}$ (factorielle
        descendante).
  \item $P(\overline{K_n}, k) = k^n$ (graphe sans arêtes).
  \item $P(T, k) = k(k-1)^{n-1}$ pour tout arbre $T$ à $n$ sommets.
\end{enumerate}
\end{proposition}

\begin{theoreme}[Relation de suppression-contraction]
\label{thm:suppression-contraction}
\index{suppression-contraction}
Pour toute arête $e = \{u,v\}$ d'un graphe $G$ :
\[
  P(G, k) = P(G - e, k) - P(G / e, k),
\]
où $G - e$ est le graphe obtenu par suppression de $e$ et $G / e$ le
graphe obtenu par contraction de $e$ (identification de $u$ et $v$).
\end{theoreme}

\begin{proof}
Les colorations propres de $G-e$ se répartissent en deux catégories :
celles où $c(u) \neq c(v)$ (qui sont exactement les colorations propres
de $G$) et celles où $c(u) = c(v)$ (qui correspondent bijectivement aux
colorations propres de $G/e$). D'où
$P(G-e, k) = P(G, k) + P(G/e, k)$.
\end{proof}

\begin{exemple}[Calcul de $P(C_4, k)$]
\label{ex:poly-chrom-C4}
En appliquant la relation de suppression-contraction à une arête de $C_4$ :
\begin{align*}
  P(C_4, k) &= P(P_4, k) - P(C_3, k) \\
    &= k(k-1)^3 - k(k-1)(k-2) \\
    &= k(k-1)\bigl[(k-1)^2 - (k-2)\bigr] \\
    &= k(k-1)(k^2 - 3k + 3).
\end{align*}
On retrouve $P(C_4, 2) = 2 \cdot 1 \cdot 1 = 2$ et
$P(C_4, 3) = 3 \cdot 2 \cdot 3 = 18$.
\end{exemple}

\section{Graphes planaires}
\label{sec:planaires}
\index{graphe planaire}

\begin{definition}[Graphe planaire]
\label{def:planaire}
Un graphe est \emph{planaire} s'il peut être dessiné dans le plan sans
croisement d'arêtes. Un tel dessin est appelé une \emph{représentation
planaire} ou un \emph{plongement planaire}. Les régions délimitées par les
arêtes dans un plongement planaire sont appelées \emph{faces}.
\end{definition}

\begin{theoreme}[Formule d'Euler]
\label{thm:euler-formule}
\index{formule d'Euler (graphes planaires)}
Soit $G$ un graphe planaire connexe à $V$ sommets, $E$ arêtes et $F$
faces (y compris la face extérieure). Alors :
\[
  V - E + F = 2.
\]
\end{theoreme}

\begin{proof}
Procédons par récurrence sur le nombre d'arêtes $E$.

\medskip\noindent
\textbf{Base.}\; Si $E = 0$, alors $G$ est réduit à un sommet ($V=1$)
et possède une seule face ($F=1$). On a $1 - 0 + 1 = 2$.

\medskip\noindent
\textbf{Pas de récurrence.}\; Supposons le résultat vrai pour tout graphe
planaire connexe à moins de $E$ arêtes, et soit $G$ un graphe planaire
connexe à $E \geq 1$ arêtes.

\emph{Cas 1 :} $G$ contient un cycle. Soit $e$ une arête appartenant à un
cycle. L'arête $e$ sépare deux faces distinctes. Le graphe $G - e$ est
encore planaire et connexe (car $e$ appartient à un cycle), avec $V$
sommets, $E-1$ arêtes et $F-1$ faces (la suppression de $e$ fusionne les
deux faces adjacentes). Par hypothèse de récurrence :
\[
  V - (E-1) + (F-1) = 2,
\]
d'où $V - E + F = 2$.

\emph{Cas 2 :} $G$ est acyclique, donc un arbre. Alors $E = V - 1$ et
$F = 1$ (la seule face est la face extérieure). On obtient
$V - (V-1) + 1 = 2$.
\end{proof}

\begin{corollaire}
\label{cor:euler-arete-borne}
Pour tout graphe planaire simple connexe à $V \geq 3$ sommets :
$E \leq 3V - 6$.
\end{corollaire}

\begin{proof}
Chaque face est délimitée par au moins $3$ arêtes, et chaque arête borde
au plus $2$ faces, d'où $2E \geq 3F$. Combiné avec $F = 2 - V + E$
(formule d'Euler), on obtient $2E \geq 3(2 - V + E) = 6 - 3V + 3E$,
soit $E \leq 3V - 6$.
\end{proof}

\begin{corollaire}
\label{cor:euler-sans-triangle}
Pour tout graphe planaire simple connexe sans triangle à $V \geq 3$ sommets :
$E \leq 2V - 4$.
\end{corollaire}

\begin{proof}
Sans triangle, chaque face est délimitée par au moins $4$ arêtes, d'où
$2E \geq 4F$. On conclut de manière analogue.
\end{proof}

\begin{theoreme}[$K_5$ n'est pas planaire]
\label{thm:K5-non-planaire}
\index{$K_5$!non-planarité}
Le graphe complet $K_5$ n'est pas planaire.
\end{theoreme}

\begin{proof}
$K_5$ possède $V = 5$ sommets et $E = \binom{5}{2} = 10$ arêtes. Or
$3V - 6 = 9 < 10 = E$, contredisant le corollaire~\ref{cor:euler-arete-borne}.
\end{proof}

\begin{theoreme}[$K_{3,3}$ n'est pas planaire]
\label{thm:K33-non-planaire}
\index{$K_{3,3}$!non-planarité}
Le graphe biparti complet $K_{3,3}$ n'est pas planaire.
\end{theoreme}

\begin{proof}
$K_{3,3}$ possède $V = 6$ sommets et $E = 9$ arêtes. Étant biparti, il
ne contient aucun cycle de longueur impaire, en particulier pas de
triangle. Par le corollaire~\ref{cor:euler-sans-triangle},
$E \leq 2V - 4 = 8 < 9 = E$, contradiction.
\end{proof}

\begin{theoreme}[Théorème de Kuratowski]
\label{thm:kuratowski}
\index{théorème de Kuratowski}
Un graphe est planaire si et seulement s'il ne contient aucun sous-graphe
homéomorphe à $K_5$ ou à $K_{3,3}$.
\end{theoreme}

\begin{exemple}[Graphes planaire et non planaire]
\label{ex:planaire}

\begin{center}
\begin{tikzpicture}[
  every node/.style={circle, draw, fill=blue!15, inner sep=2pt,
                     minimum size=16pt, font=\footnotesize}
]
% K_4 planaire
\begin{scope}[xshift=0cm]
  \node (a) at (0,0)   {1};
  \node (b) at (2,0)   {2};
  \node (c) at (1,1.7) {3};
  \node (d) at (1,0.6) {4};
  \draw (a)--(b)--(c)--(a) (a)--(d) (b)--(d) (c)--(d);
  \node[draw=none, fill=none] at (1,-0.7) {$K_4$ (planaire)};
\end{scope}

% K_{3,3} non planaire
\begin{scope}[xshift=6cm]
  \node[fill=red!15] (u1) at (0,2) {$u_1$};
  \node[fill=red!15] (u2) at (1.5,2) {$u_2$};
  \node[fill=red!15] (u3) at (3,2) {$u_3$};
  \node[fill=green!20] (w1) at (0,0) {$w_1$};
  \node[fill=green!20] (w2) at (1.5,0) {$w_2$};
  \node[fill=green!20] (w3) at (3,0) {$w_3$};
  \draw (u1)--(w1) (u1)--(w2) (u1)--(w3);
  \draw (u2)--(w1) (u2)--(w2) (u2)--(w3);
  \draw (u3)--(w1) (u3)--(w2) (u3)--(w3);
  \node[draw=none, fill=none] at (1.5,-0.7) {$K_{3,3}$ (non planaire)};
\end{scope}
\end{tikzpicture}
\end{center}
\end{exemple}

\section{Graphes bipartis}
\label{sec:graphes-bipartis}
\index{graphe biparti}

\begin{definition}[Graphe biparti]
\label{def:biparti}
Un graphe $G=(V,E)$ est \emph{biparti} s'il existe une partition
$V = A \cup B$ ($A \cap B = \varnothing$) telle que toute arête de $G$
relie un sommet de $A$ à un sommet de $B$. La paire $(A,B)$ est appelée
une \emph{bipartition} de $G$.
\end{definition}

\begin{theoreme}[Caractérisation des graphes bipartis]
\label{thm:biparti-caract}
\index{graphe biparti!caractérisation}
Un graphe $G$ est biparti si et seulement s'il ne contient aucun cycle de
longueur impaire.
\end{theoreme}

\begin{proof}
\textbf{($\Rightarrow$)}\;
Soit $G$ biparti de bipartition $(A,B)$ et soit
$C = v_0, v_1, \dots, v_\ell = v_0$ un cycle dans $G$. Les sommets de $C$
alternent entre $A$ et $B$ : si $v_0 \in A$, alors $v_1 \in B$,
$v_2 \in A$, etc. Puisque $v_\ell = v_0 \in A$, on a $\ell$ pair. Donc
tout cycle est de longueur paire.

\medskip\noindent
\textbf{($\Leftarrow$)}\;
Supposons que $G$ ne contient aucun cycle de longueur impaire. On peut
supposer $G$ connexe (sinon, on traite chaque composante séparément).
Fixons un sommet $s$ et posons :
\[
  A = \{v \in V : d(s,v) \text{ est pair}\}, \quad
  B = \{v \in V : d(s,v) \text{ est impair}\},
\]
où $d(s,v)$ est la distance (longueur d'un plus court chemin) de $s$ à $v$.

Montrons que $(A,B)$ est une bipartition. Soit $\{u,w\} \in E$ ; nous
devons montrer que $u$ et $w$ ne sont pas dans la même partie. Par
l'absurde, supposons $u, w \in A$ (le cas $B$ est analogue) avec
$d(s,u) = 2p$ et $d(s,w) = 2q$. Soient $P_u$ un plus court chemin de
$s$ à $u$ et $P_w$ un plus court chemin de $s$ à $w$. La concaténation de
$P_u$, de l'arête $\{u,w\}$ et du renversé de $P_w$ donne une marche
fermée de longueur $2p + 1 + 2q$, qui est impaire. Cette marche fermée
contient un cycle de longueur impaire, contredisant l'hypothèse.
\end{proof}

\begin{exemple}[Coloration et bipartition]
\label{ex:biparti}

\begin{center}
\begin{tikzpicture}[
  every node/.style={circle, draw, inner sep=2pt,
                     minimum size=18pt, font=\footnotesize}
]
  \node[fill=red!20]   (a) at (0,1.5) {$a$};
  \node[fill=blue!20]  (b) at (1.5,2.5) {$b$};
  \node[fill=red!20]   (c) at (3,1.5) {$c$};
  \node[fill=blue!20]  (d) at (1.5,0.5) {$d$};
  \draw (a)--(b)--(c)--(d)--(a);

  \node[draw=none, fill=none] at (1.5,-0.5) {$C_4$ : biparti ($2$-colorable)};
\end{tikzpicture}
\end{center}
\end{exemple}

\section{Couplages et théorèmes de König et Hall}
\label{sec:couplages}
\index{couplage}

\begin{definition}[Couplage]
\label{def:couplage}
Un \emph{couplage} dans un graphe $G=(V,E)$ est un ensemble $M \subseteq E$
d'arêtes deux à deux disjointes (sans sommet commun). Un couplage est
\emph{maximal} s'il ne peut être étendu, et \emph{maximum} si $\card{M}$
est maximale. Un couplage \emph{parfait} est un couplage qui couvre tous
les sommets.
\end{definition}

\begin{definition}[Couverture par sommets]
\label{def:couverture}
\index{couverture par sommets}
Une \emph{couverture par sommets} de $G$ est un ensemble $S \subseteq V$
tel que toute arête de $G$ est incidente à au moins un sommet de $S$.
\end{definition}

\begin{theoreme}[Théorème de König]
\label{thm:konig}
\index{théorème de König}
Dans un graphe biparti, la taille d'un couplage maximum est égale à la
taille d'une couverture minimale par sommets.
\end{theoreme}

\begin{theoreme}[Théorème du mariage de Hall]
\label{thm:hall}
\index{théorème de Hall}
Soit $G$ un graphe biparti de bipartition $(A,B)$. Il existe un couplage
couvrant tous les sommets de $A$ si et seulement si la \emph{condition de
Hall} est satisfaite :
\[
  \card{N(S)} \geq \card{S} \quad \text{pour tout } S \subseteq A,
\]
où $N(S) = \{b \in B : \exists\, a \in S,\; \{a,b\} \in E\}$ désigne le
voisinage de $S$.
\end{theoreme}

\begin{proof}
\textbf{($\Rightarrow$)}\;
Si un couplage $M$ couvre $A$, alors pour tout $S \subseteq A$, les
sommets de $S$ sont appariés à $\card{S}$ sommets distincts de $B$
appartenant à $N(S)$, d'où $\card{N(S)} \geq \card{S}$.

\medskip\noindent
\textbf{($\Leftarrow$)}\;
Procédons par récurrence sur $\card{A}$.

\emph{Base.} Si $\card{A} = 0$ ou $\card{A}= 1$, le résultat est immédiat.

\emph{Pas de récurrence.} Supposons $\card{A} \geq 2$.

\emph{Cas 1 :} $\card{N(S)} \geq \card{S} + 1$ pour tout
$\varnothing \neq S \subsetneq A$. Choisissons un sommet $a \in A$ et un
voisin $b \in N(a)$. Le graphe $G' = G - \{a,b\}$ est biparti de
bipartition $(A' = A \setminus\{a\},\, B' = B \setminus\{b\})$. Pour
tout $S \subseteq A'$ non vide, $\card{N_{G'}(S)} \geq \card{N_G(S)} - 1
\geq (\card{S}+1) - 1 = \card{S}$. Par récurrence, $G'$ possède un
couplage couvrant $A'$, que l'on complète par l'arête $\{a,b\}$.

\emph{Cas 2 :} il existe $\varnothing \neq S_0 \subsetneq A$ tel que
$\card{N(S_0)} = \card{S_0}$. Par récurrence, le sous-graphe induit par
$S_0 \cup N(S_0)$ possède un couplage $M_1$ couvrant $S_0$. Posons
$A'' = A \setminus S_0$, $B'' = B \setminus N(S_0)$, et considérons le
sous-graphe $G''$ induit par $A'' \cup B''$. Pour tout $S \subseteq A''$,
on a :
\[
  \card{N_{G''}(S)} = \card{N_G(S \cup S_0)} - \card{N(S_0)}
  \geq \card{S \cup S_0} - \card{S_0} = \card{S}.
\]
(Si $\card{N_{G''}(S)} < \card{S}$, alors
$\card{N_G(S \cup S_0)} < \card{S} + \card{S_0} = \card{S \cup S_0}$,
contradiction avec la condition de Hall.) Par récurrence, $G''$ possède
un couplage $M_2$ couvrant $A''$. L'union $M_1 \cup M_2$ couvre $A$.
\end{proof}

\begin{exemple}[Couplage dans un graphe biparti]
\label{ex:couplage}

\begin{center}
\begin{tikzpicture}[
  every node/.style={circle, draw, inner sep=2pt,
                     minimum size=18pt, font=\footnotesize}
]
  \node[fill=red!20]   (a1) at (0,3) {$a_1$};
  \node[fill=red!20]   (a2) at (0,2) {$a_2$};
  \node[fill=red!20]   (a3) at (0,1) {$a_3$};
  \node[fill=red!20]   (a4) at (0,0) {$a_4$};

  \node[fill=blue!20]  (b1) at (3,3) {$b_1$};
  \node[fill=blue!20]  (b2) at (3,2) {$b_2$};
  \node[fill=blue!20]  (b3) at (3,1) {$b_3$};
  \node[fill=blue!20]  (b4) at (3,0) {$b_4$};

  % Arêtes du couplage
  \draw[very thick, red!70!black] (a1)--(b2);
  \draw[very thick, red!70!black] (a2)--(b1);
  \draw[very thick, red!70!black] (a3)--(b3);
  \draw[very thick, red!70!black] (a4)--(b4);

  % Autres arêtes
  \draw[thin, gray] (a1)--(b1) (a1)--(b3);
  \draw[thin, gray] (a2)--(b2);
  \draw[thin, gray] (a3)--(b4);
  \draw[thin, gray] (a4)--(b3);

  \node[draw=none, fill=none] at (1.5,-0.8) {Couplage parfait en rouge};
\end{tikzpicture}
\end{center}
\end{exemple}

\section{Exercices}
\label{sec:exo-ch10}

\begin{exercice}
Calculer le nombre chromatique des graphes suivants : $K_6$, $C_7$,
$K_{3,3}$, le graphe de Petersen.
\end{exercice}

\begin{exercice}
Calculer le polynôme chromatique de $C_5$ par suppression-contraction.
\end{exercice}

\begin{exercice}
Vérifier la formule d'Euler pour le cube ($V=8$, $E=12$, $F=6$).
\end{exercice}

\begin{exercice}
Montrer que tout graphe planaire possède un sommet de degré au plus $5$.
\textit{Indication : utiliser $E \leq 3V - 6$ et la formule des degrés.}
\end{exercice}

\begin{exercice}
Montrer que le graphe de Petersen n'est pas planaire.
\textit{Indication : il contient un sous-graphe homéomorphe à $K_{3,3}$.}
\end{exercice}

\begin{exercice}
Déterminer le polynôme chromatique du graphe complet biparti $K_{2,3}$.
\end{exercice}

\begin{exercice}
On considère $n$ filles, chacune connaissant un certain nombre de garçons.
Chaque fille connaît au moins $k$ garçons, et chaque garçon est connu
d'au plus $k$ filles. Montrer qu'il existe un couplage complet (chaque
fille est appariée à un garçon qu'elle connaît).
\textit{Indication : vérifier la condition de Hall.}
\end{exercice}

\begin{exercice}
Un graphe planaire biparti à $V \geq 3$ sommets satisfait $E \leq 2V - 4$.
Montrer ce résultat et en déduire que $K_{3,3}$ n'est pas planaire.
\end{exercice}

\begin{exercice}
Montrer que le polynôme chromatique d'un arbre à $n$ sommets est
$P(T, k) = k(k-1)^{n-1}$.
\textit{Indication : récurrence en utilisant un sommet pendant.}
\end{exercice}

\bigskip
\begin{center}
\fbox{\parbox{0.85\textwidth}{%
\textbf{Résumé du chapitre~10.}\;
Le nombre chromatique $\chi(G)$ est le nombre minimal de couleurs pour une
coloration propre. Le théorème de Brooks affine la borne $\Delta(G)+1$.
Le polynôme chromatique se calcule par suppression-contraction. La formule
d'Euler $V-E+F=2$ caractérise les plongements planaires et entraîne
$E \leq 3V-6$, ce qui permet de démontrer la non-planarité de $K_5$ et
$K_{3,3}$. Un graphe est biparti si et seulement s'il ne contient pas de
cycle impair. Le théorème de Hall donne une condition nécessaire et
suffisante pour l'existence d'un couplage complet dans un graphe biparti.
}}
\end{center}


%%%%%%%%%%%%%%%%%%%%%%%%%%%%%%%%%%%%%%%%%%%%%%%%%%%%%%%%%%%%%%%%%%%%%
%%  CHAPITRE 11 — Introduction à la Théorie des Codes
%%%%%%%%%%%%%%%%%%%%%%%%%%%%%%%%%%%%%%%%%%%%%%%%%%%%%%%%%%%%%%%%%%%%%
\chapter{Introduction à la Théorie des Codes}
\label{ch:codes}
\index{théorie des codes}

\section{Motivation}
\label{sec:codes-motivation}

Les données numériques --- texte, images, sons --- sont transmises sous
forme de suites de symboles (typiquement des bits $0$ et $1$) à travers
des canaux imparfaits : câbles, ondes radio, supports de stockage. Des
erreurs de transmission peuvent altérer les données. La \emph{théorie des
codes correcteurs} vise à ajouter de la redondance aux messages de manière
structurée, afin de \emph{détecter} et, si possible, \emph{corriger} les
erreurs introduites lors de la transmission.

\begin{remarque}
La théorie des codes se situe à l'intersection de l'algèbre linéaire,
de la combinatoire et de la théorie de l'information. Elle constitue une
application majeure des mathématiques discrètes.
\end{remarque}

\section{Alphabet, mots et distance de Hamming}
\label{sec:hamming-distance}
\index{distance de Hamming}\index{poids de Hamming}

\begin{definition}[Alphabet et mots]
\label{def:alphabet}
Un \emph{alphabet} est un ensemble fini $\mathcal{A}$ de symboles. Un
\emph{mot} de longueur $n$ sur $\mathcal{A}$ est un élément de
$\mathcal{A}^n$. Dans ce chapitre, nous travaillons principalement avec
l'alphabet binaire $\mathcal{A} = \mathbb{F}_2 = \{0,1\}$, le corps à
deux éléments.
\end{definition}

\begin{definition}[Distance de Hamming]
\label{def:distance-hamming}
\index{distance de Hamming!définition}
Soient $\mathbf{x} = (x_1, \dots, x_n)$ et
$\mathbf{y} = (y_1, \dots, y_n)$ deux mots de $\mathcal{A}^n$. La
\emph{distance de Hamming} entre $\mathbf{x}$ et $\mathbf{y}$ est :
\[
  d(\mathbf{x}, \mathbf{y}) = \card{\{i : x_i \neq y_i\}}.
\]
\end{definition}

\begin{definition}[Poids de Hamming]
\label{def:poids-hamming}
Le \emph{poids de Hamming} d'un mot $\mathbf{x} \in \mathbb{F}_2^n$ est
le nombre de composantes non nulles :
\[
  w(\mathbf{x}) = d(\mathbf{x}, \mathbf{0}) = \card{\{i : x_i \neq 0\}}.
\]
\end{definition}

\begin{proposition}
La distance de Hamming est une métrique sur $\mathcal{A}^n$ :
\begin{enumerate}[label=(\roman*)]
  \item $d(\mathbf{x}, \mathbf{y}) \geq 0$, avec égalité si et seulement
        si $\mathbf{x} = \mathbf{y}$.
  \item $d(\mathbf{x}, \mathbf{y}) = d(\mathbf{y}, \mathbf{x})$
        (symétrie).
  \item $d(\mathbf{x}, \mathbf{z}) \leq d(\mathbf{x}, \mathbf{y}) +
        d(\mathbf{y}, \mathbf{z})$ (inégalité triangulaire).
\end{enumerate}
\end{proposition}

\begin{proof}
Les propriétés (i) et (ii) sont immédiates. Pour (iii), si
$x_i \neq z_i$, alors nécessairement $x_i \neq y_i$ ou $y_i \neq z_i$
(ou les deux). Donc chaque position contribuant à $d(\mathbf{x},
\mathbf{z})$ contribue à $d(\mathbf{x}, \mathbf{y})$ ou
$d(\mathbf{y}, \mathbf{z})$.
\end{proof}

\section{Codes détecteurs et correcteurs d'erreurs}
\label{sec:detection-correction}
\index{code détecteur}\index{code correcteur}

\begin{definition}[Code en bloc]
\label{def:code-bloc}
Un \emph{code en bloc} de longueur $n$ et de taille $M$ sur un alphabet
$\mathcal{A}$ est un sous-ensemble $C \subseteq \mathcal{A}^n$ avec
$\card{C} = M$. Les éléments de $C$ sont les \emph{mots de code}. La
\emph{distance minimale} du code est :
\[
  d(C) = \min_{\substack{\mathbf{x}, \mathbf{y} \in C \\
         \mathbf{x} \neq \mathbf{y}}} d(\mathbf{x}, \mathbf{y}).
\]
On note un tel code un $(n, M, d)$-code.
\end{definition}

\begin{theoreme}[Capacités de détection et de correction]
\label{thm:detection-correction}
Un code de distance minimale $d$ peut :
\begin{enumerate}[label=(\roman*)]
  \item détecter jusqu'à $d-1$ erreurs ;
  \item corriger jusqu'à $\lfloor (d-1)/2 \rfloor$ erreurs.
\end{enumerate}
\end{theoreme}

\begin{proof}
(i) Si au plus $d-1$ erreurs transforment un mot de code $\mathbf{c}$ en
un mot $\mathbf{r}$, alors $d(\mathbf{c}, \mathbf{r}) \leq d-1 < d$,
donc $\mathbf{r}$ n'est pas un mot de code (sinon la distance minimale
serait $< d$). L'erreur est donc détectée.

(ii) Posons $t = \lfloor (d-1)/2 \rfloor$. Si $\mathbf{c}$ est envoyé
et $\mathbf{r}$ est reçu avec au plus $t$ erreurs, alors
$d(\mathbf{c}, \mathbf{r}) \leq t$. Pour tout autre mot de code
$\mathbf{c'} \neq \mathbf{c}$, par l'inégalité triangulaire :
$d \leq d(\mathbf{c}, \mathbf{c'}) \leq d(\mathbf{c}, \mathbf{r}) +
d(\mathbf{r}, \mathbf{c'})$, d'où
$d(\mathbf{r}, \mathbf{c'}) \geq d - t > t$. Le mot de code le plus
proche de $\mathbf{r}$ est donc $\mathbf{c}$.
\end{proof}

\begin{theoreme}[Borne de l'empilement de sphères (borne de Hamming)]
\label{thm:hamming-bound}
\index{borne de Hamming}
Si $C$ est un code binaire $(n, M, d)$ avec $d = 2t+1$, alors :
\[
  M \cdot \sum_{i=0}^{t} \binom{n}{i} \leq 2^n.
\]
\end{theoreme}

\begin{proof}
Les boules de rayon $t$ centrées sur les mots de code sont deux à deux
disjointes (car la distance entre deux mots de code est $\geq 2t+1$).
Chaque boule contient $\sum_{i=0}^t \binom{n}{i}$ mots de
$\mathbb{F}_2^n$, et l'espace total contient $2^n$ mots.
\end{proof}

\begin{definition}[Code parfait]
\label{def:code-parfait}
\index{code parfait}
Un code atteignant l'égalité dans la borne de Hamming est appelé
\emph{code parfait} : les boules de rayon $t$ centrées sur les mots de
code forment une partition de $\mathbb{F}_2^n$.
\end{definition}

\section{Codes linéaires}
\label{sec:codes-lineaires}
\index{code linéaire}

\begin{definition}[Code linéaire]
\label{def:code-lineaire}
Un \emph{code linéaire} binaire $[n, k, d]$ est un sous-espace vectoriel
$C$ de dimension $k$ de $\mathbb{F}_2^n$ dont la distance minimale est
$d$. Le \emph{rendement} (ou \emph{taux}) du code est $R = k/n$.
\end{definition}

\begin{proposition}
Pour un code linéaire, la distance minimale est égale au poids minimal
d'un mot de code non nul :
\[
  d(C) = \min_{\mathbf{c} \in C \setminus \{\mathbf{0}\}} w(\mathbf{c}).
\]
\end{proposition}

\begin{proof}
Soient $\mathbf{x}, \mathbf{y} \in C$ distincts. Alors
$\mathbf{x} - \mathbf{y} = \mathbf{x} + \mathbf{y} \in C$ (dans
$\mathbb{F}_2$, $- = +$) et
$d(\mathbf{x}, \mathbf{y}) = w(\mathbf{x}+\mathbf{y})$. Comme
$\mathbf{x}+\mathbf{y}$ parcourt tous les mots de code non nuls lorsque
$(\mathbf{x},\mathbf{y})$ parcourt les paires distinctes, le minimum est
atteint.
\end{proof}

\begin{definition}[Matrice génératrice et matrice de contrôle]
\label{def:matrice-generatrice}
\index{matrice génératrice}\index{matrice de contrôle}
Soit $C$ un code linéaire $[n,k]$.
\begin{itemize}
  \item Une \emph{matrice génératrice} de $C$ est une matrice
        $G \in \mathbb{F}_2^{k \times n}$ dont les lignes forment une
        base de $C$. On a $C = \{\mathbf{u} G : \mathbf{u} \in
        \mathbb{F}_2^k\}$.
  \item Une \emph{matrice de contrôle} (ou de parité) est une matrice
        $H \in \mathbb{F}_2^{(n-k) \times n}$ telle que
        $C = \ker(H) = \{\mathbf{x} \in \mathbb{F}_2^n : H\mathbf{x}^T
        = \mathbf{0}\}$.
\end{itemize}
Si $G = [I_k \mid A]$ (forme systématique), alors $H = [-A^T \mid
I_{n-k}] = [A^T \mid I_{n-k}]$ (dans $\mathbb{F}_2$).
\end{definition}

\begin{theoreme}[Distance minimale et matrice de contrôle]
\label{thm:distance-H}
La distance minimale $d$ d'un code linéaire $C$ de matrice de contrôle $H$
est le plus petit nombre de colonnes de $H$ linéairement dépendantes.
\end{theoreme}

\begin{proof}
Un mot $\mathbf{c}$ de poids $w$ appartient à $C$ si et seulement si
$H\mathbf{c}^T = \mathbf{0}$, c'est-à-dire si et seulement si la somme
des $w$ colonnes de $H$ correspondant aux positions non nulles de
$\mathbf{c}$ est nulle. Le poids minimal non nul correspond donc au plus
petit nombre de colonnes linéairement dépendantes.
\end{proof}

\subsection{Décodage par syndrome}
\label{subsec:syndrome}
\index{syndrome}

\begin{definition}[Syndrome]
\label{def:syndrome}
Soit $H$ la matrice de contrôle d'un code $[n,k]$ et soit $\mathbf{r}$
un mot reçu. Le \emph{syndrome} de $\mathbf{r}$ est
$\mathbf{s} = H\mathbf{r}^T \in \mathbb{F}_2^{n-k}$.
\end{definition}

\begin{proposition}
Le syndrome ne dépend que de l'erreur, pas du mot de code envoyé. Si
$\mathbf{r} = \mathbf{c} + \mathbf{e}$ (mot reçu = mot envoyé + erreur),
alors $\mathbf{s} = H\mathbf{e}^T$.
\end{proposition}

\begin{proof}
$H\mathbf{r}^T = H(\mathbf{c}+\mathbf{e})^T = H\mathbf{c}^T +
H\mathbf{e}^T = \mathbf{0} + H\mathbf{e}^T = H\mathbf{e}^T$.
\end{proof}

\begin{remarque}
Le décodage par syndrome consiste à :
\begin{enumerate}
  \item calculer $\mathbf{s} = H\mathbf{r}^T$ ;
  \item trouver le vecteur d'erreur $\mathbf{e}$ de poids minimal tel que
        $H\mathbf{e}^T = \mathbf{s}$ (le \emph{chef de classe}) ;
  \item décoder $\hat{\mathbf{c}} = \mathbf{r} - \mathbf{e}$.
\end{enumerate}
\end{remarque}

\section{Codes de Hamming}
\label{sec:codes-hamming}
\index{code de Hamming}

\begin{definition}[Code de Hamming $\mathrm{Ham}(r, 2)$]
\label{def:hamming-code}
Pour un entier $r \geq 2$, le \emph{code de Hamming} binaire
$\mathrm{Ham}(r, 2)$ est le code linéaire $[n, k, 3]$ avec $n = 2^r - 1$
et $k = 2^r - 1 - r$, dont la matrice de contrôle $H$ de taille
$r \times n$ a pour colonnes tous les vecteurs non nuls de
$\mathbb{F}_2^r$.
\end{definition}

\begin{theoreme}[Propriétés du code de Hamming]
\label{thm:hamming-proprietes}
Le code $\mathrm{Ham}(r, 2)$ possède les propriétés suivantes :
\begin{enumerate}[label=(\roman*)]
  \item Sa distance minimale est $d = 3$.
  \item Il corrige toute erreur simple (affectant au plus un bit).
  \item C'est un code parfait.
\end{enumerate}
\end{theoreme}

\begin{proof}
(i) Les colonnes de $H$ sont toutes distinctes et non nulles, donc deux
colonnes ne sont jamais proportionnelles (dans $\mathbb{F}_2$, il n'y a
que $0$ et $1$). Ainsi $d \geq 3$ par le théorème~\ref{thm:distance-H}.
De plus, les trois premières colonnes de $H$ (les vecteurs de la base
canonique) satisfont une relation de dépendance si $r \geq 2$, par
exemple les colonnes correspondant à $(1,0,\dots,0)^T$, $(0,1,0,\dots)^T$
et $(1,1,0,\dots)^T$ sont liées. Donc $d = 3$.

(ii) Découle de $t = \lfloor(3-1)/2\rfloor = 1$.

(iii) On vérifie la borne de Hamming :
$M \cdot (1 + n) = 2^k \cdot 2^r = 2^{k+r} = 2^n$, donc l'égalité est
atteinte.
\end{proof}

\begin{exemple}[{Code de Hamming $\mathrm{Ham}(3, 2)$ : le code $[7,4,3]$}]
\label{ex:hamming-7-4}
\index{code de Hamming!$[7{,}4{,}3]$}

La matrice de contrôle est :
\[
  H = \begin{pmatrix}
    0 & 0 & 0 & 1 & 1 & 1 & 1 \\
    0 & 1 & 1 & 0 & 0 & 1 & 1 \\
    1 & 0 & 1 & 0 & 1 & 0 & 1
  \end{pmatrix},
\]
dont les colonnes sont les représentations binaires des entiers
$1, 2, \dots, 7$. La matrice génératrice sous forme systématique est :
\[
  G = \begin{pmatrix}
    1 & 0 & 0 & 0 & 0 & 1 & 1 \\
    0 & 1 & 0 & 0 & 1 & 0 & 1 \\
    0 & 0 & 1 & 0 & 1 & 1 & 0 \\
    0 & 0 & 0 & 1 & 1 & 1 & 1
  \end{pmatrix}.
\]

Si le mot reçu est $\mathbf{r} = (1,0,1,1,0,1,0)$, le syndrome est
$\mathbf{s} = H\mathbf{r}^T = (1,1,0)^T$, qui est la colonne~$6$ de $H$.
L'erreur est donc en position~$6$ et le mot corrigé est
$(1,0,1,1,0,0,0)$.
\end{exemple}

\section{Codes de répétition et codes de parité}
\label{sec:repetition-parite}
\index{code de répétition}\index{code de parité}

\begin{definition}[Code de répétition]
\label{def:code-repetition}
Le \emph{code de répétition} de longueur $n$ est le code $[n, 1, n]$
défini par $C = \{(0,0,\dots,0), (1,1,\dots,1)\}$. Chaque bit
d'information est répété $n$ fois.
\end{definition}

\begin{remarque}
Le code de répétition corrige $\lfloor (n-1)/2 \rfloor$ erreurs mais
a un rendement très faible $R = 1/n$.
\end{remarque}

\begin{definition}[Code de parité]
\label{def:code-parite}
Le \emph{code de contrôle de parité} de longueur $n$ est le code
$[n, n-1, 2]$ défini par :
\[
  C = \{(x_1, \dots, x_n) \in \mathbb{F}_2^n : x_1 + x_2 + \cdots +
  x_n = 0\}.
\]
La matrice de contrôle est $H = (1\; 1\; \cdots\; 1)$.
\end{definition}

\begin{remarque}
Le code de parité détecte toute erreur simple ($d=2$) mais ne corrige
aucune erreur. Son rendement $R = (n-1)/n$ est élevé.
\end{remarque}

\begin{proposition}
Les codes de répétition et de parité sont duaux l'un de l'autre :
le dual du code de répétition $[n,1,n]$ est le code de parité $[n,n-1,2]$,
et réciproquement.
\end{proposition}

\section{Codes de Reed-Solomon}
\label{sec:reed-solomon}
\index{code de Reed-Solomon}

\begin{remarque}
Les \emph{codes de Reed-Solomon} sont des codes linéaires définis sur des
corps finis $\mathbb{F}_q$ (avec $q > 2$ en général). Un code de
Reed-Solomon $\mathrm{RS}(n, k)$ sur $\mathbb{F}_q$ a pour paramètres
$[n, k, n-k+1]$ avec $n = q-1$. Il atteint la \emph{borne de Singleton}
$d \leq n - k + 1$ et est donc un code MDS (\emph{maximum distance
separable}).

Ces codes sont utilisés dans de nombreuses applications pratiques : CD,
DVD, codes QR, transmissions spatiales, etc. Leur étude détaillée dépasse
le cadre de ce cours d'introduction mais illustre la puissance des
méthodes algébriques (polynômes sur les corps finis) en théorie des codes.
\end{remarque}

\section{Lien avec les mathématiques discrètes}
\label{sec:codes-maths-discretes}

La théorie des codes fait appel de manière essentielle à plusieurs
branches des mathématiques discrètes :

\begin{itemize}
  \item \textbf{Algèbre linéaire sur les corps finis.}\;
  Les codes linéaires sont des sous-espaces vectoriels de
  $\mathbb{F}_q^n$. Les opérations de codage, décodage et détection
  d'erreurs se ramènent à des calculs matriciels sur $\mathbb{F}_q$.

  \item \textbf{Combinatoire.}\;
  Le dénombrement des mots de code de poids donné (distribution des poids)
  utilise les coefficients binomiaux et les identités combinatoires. Les
  bornes fondamentales (Hamming, Singleton, Plotkin, Gilbert-Varshamov)
  sont de nature combinatoire.

  \item \textbf{Théorie des graphes.}\;
  Certains codes sont définis à partir de graphes (codes LDPC, codes
  de Tanner). Les graphes de Cayley interviennent dans l'analyse de
  codes algébriques.

  \item \textbf{Théorie des nombres.}\;
  Les corps finis $\mathbb{F}_{p^r}$ sont construits à l'aide de
  polynômes irréductibles sur $\mathbb{F}_p$. Les racines primitives
  et la structure multiplicative de $\mathbb{F}_q^*$ jouent un rôle
  central dans la construction des codes BCH et Reed-Solomon.
\end{itemize}

\section{Exercices}
\label{sec:exo-ch11}

\begin{exercice}
Calculer la distance de Hamming entre les mots $\mathbf{x} = (1,0,1,1,0)$
et $\mathbf{y} = (0,1,1,0,0)$.
\end{exercice}

\begin{exercice}
Un code binaire $C$ a pour mots de code :
$\{00000, 01101, 10110, 11011\}$.
\begin{enumerate}[label=(\alph*)]
  \item Calculer la distance minimale de $C$.
  \item Combien d'erreurs $C$ peut-il détecter ? Corriger ?
  \item $C$ est-il linéaire ?
\end{enumerate}
\end{exercice}

\begin{exercice}
Construire la matrice de contrôle du code de Hamming $\mathrm{Ham}(4, 2)$,
c'est-à-dire le code $[15, 11, 3]$. Vérifier que c'est un code parfait.
\end{exercice}

\begin{exercice}
Soit $C$ le code de Hamming $[7,4,3]$ de l'exemple~\ref{ex:hamming-7-4}.
\begin{enumerate}[label=(\alph*)]
  \item Coder le message $\mathbf{u} = (1,0,1,1)$ en calculant
        $\mathbf{c} = \mathbf{u}G$.
  \item On reçoit $\mathbf{r} = (1,1,1,1,0,0,1)$. Calculer le syndrome
        et corriger l'erreur.
\end{enumerate}
\end{exercice}

\begin{exercice}
Montrer que le code de parité $[n, n-1, 2]$ détecte exactement une erreur
mais ne peut en corriger aucune.
\end{exercice}

\begin{exercice}
Montrer que la borne de Singleton s'écrit $M \leq q^{n-d+1}$ pour un
code $(n, M, d)$ sur un alphabet de taille $q$, c'est-à-dire $d \leq n-k+1$
pour un code linéaire $[n,k,d]$.
\textit{Indication : projeter les mots de code sur $n-d+1$ coordonnées.}
\end{exercice}

\begin{exercice}
On souhaite transmettre des messages de $4$ bits avec la capacité de
corriger $1$ erreur. Quel est le nombre minimal de bits de redondance
nécessaires ? En déduire que le code de Hamming $[7,4,3]$ est optimal
pour ces paramètres.
\end{exercice}

\begin{exercice}
Montrer que le dual d'un code de Hamming $[2^r - 1, 2^r - 1 - r, 3]$ est
un code $[2^r - 1, r, 2^{r-1}]$ (code simplexe). Calculer la distribution
des poids du code simplexe pour $r=3$.
\end{exercice}

\bigskip
\begin{center}
\fbox{\parbox{0.85\textwidth}{%
\textbf{Résumé du chapitre~11.}\;
La théorie des codes correcteurs d'erreurs associe aux messages une
redondance structurée. La distance de Hamming mesure la différence entre
mots et détermine les capacités de détection ($d-1$ erreurs) et de
correction ($\lfloor(d-1)/2\rfloor$ erreurs). Les codes linéaires, décrits
par des matrices génératrices et de contrôle, permettent un codage et un
décodage efficaces. Les codes de Hamming sont des codes parfaits
$1$-correcteurs. Cette théorie illustre l'interaction profonde entre
algèbre linéaire, combinatoire et théorie des graphes.
}}
\end{center}


%%%%%%%%%%%%%%%%%%%%%%%%%%%%%%%%%%%%%%%%%%%%%%%%%%%%%%%%%%%%%%%%%%%%%
%%  ANNEXE — Table d'identités combinatoires
%%%%%%%%%%%%%%%%%%%%%%%%%%%%%%%%%%%%%%%%%%%%%%%%%%%%%%%%%%%%%%%%%%%%%
\appendix
\chapter{Table d'identités combinatoires}
\label{app:identites}
\index{identités combinatoires}

Cette annexe rassemble les principales identités combinatoires utilisées
dans ce cours et au-delà. Les preuves de la plupart de ces identités
figurent dans les chapitres précédents ou peuvent être établies par les
méthodes qui y sont développées (récurrence, double comptage, fonctions
génératrices).

\section{Identités sur les coefficients binomiaux}
\label{sec:identites-binomiales}
\index{coefficient binomial!identités}

\renewcommand{\arraystretch}{1.8}
\begin{longtable}{@{}>{\raggedright\arraybackslash}p{5cm}
                    >{$\displaystyle}l<{$}@{}}
\toprule
\textbf{Nom} & \textbf{Identité} \\
\midrule
\endfirsthead
\toprule
\textbf{Nom} & \textbf{Identité} \\
\midrule
\endhead
\bottomrule
\endfoot

Symétrie &
  \binom{n}{k} = \binom{n}{n-k} \\

Absorption &
  k\binom{n}{k} = n\binom{n-1}{k-1} \\

Formule de Pascal\index{formule de Pascal} &
  \binom{n}{k} = \binom{n-1}{k-1} + \binom{n-1}{k} \\

Théorème binomial\index{théorème binomial} &
  (x+y)^n = \sum_{k=0}^{n} \binom{n}{k} x^k y^{n-k} \\

Somme d'une ligne &
  \sum_{k=0}^{n} \binom{n}{k} = 2^n \\

Somme alternée &
  \sum_{k=0}^{n} (-1)^k \binom{n}{k} = 0 \quad (n \geq 1) \\

Identité de Vandermonde\index{identité de Vandermonde} &
  \sum_{k=0}^{r} \binom{m}{k}\binom{n}{r-k} = \binom{m+n}{r} \\

Convolution de Vandermonde &
  \sum_{k} \binom{r}{k}\binom{s}{n-k} = \binom{r+s}{n} \\

Identité du bâton de hockey\index{identité du bâton de hockey} &
  \sum_{i=r}^{n} \binom{i}{r} = \binom{n+1}{r+1} \\

Somme des carrés &
  \sum_{k=0}^{n} \binom{n}{k}^2 = \binom{2n}{n} \\

Binomiale supérieure &
  \sum_{k=0}^{n} \binom{k}{r} = \binom{n+1}{r+1} \\

Formule de Chu-Vandermonde &
  \sum_{k=0}^{r} \binom{-n}{k}(-1)^k \binom{m}{r-k} = \binom{m+n}{r} \\

Double comptage &
  \binom{n}{k}\binom{k}{j} = \binom{n}{j}\binom{n-j}{k-j} \\

\end{longtable}

\section{Nombres de Stirling}
\label{sec:stirling}
\index{nombre de Stirling!de seconde espèce}
\index{nombre de Stirling!de première espèce}

\renewcommand{\arraystretch}{1.8}
\begin{longtable}{@{}>{\raggedright\arraybackslash}p{5cm}
                    >{$\displaystyle}l<{$}@{}}
\toprule
\textbf{Nom / Identité} & \textbf{Formule} \\
\midrule
\endfirsthead
\toprule
\textbf{Nom / Identité} & \textbf{Formule} \\
\midrule
\endhead
\bottomrule
\endfoot

Stirling 2\textsuperscript{e} espèce --- récurrence &
  S(n,k) = k\, S(n-1,k) + S(n-1, k-1) \\

Stirling 2\textsuperscript{e} espèce --- formule close &
  S(n,k) = \frac{1}{k!}\sum_{j=0}^{k} (-1)^{k-j}\binom{k}{j} j^n \\

Nombre de Bell\index{nombre de Bell} &
  B_n = \sum_{k=0}^{n} S(n,k) \\

Relation avec les puissances &
  x^n = \sum_{k=0}^{n} S(n,k)\, x^{(k)} \\
  & \text{où } x^{(k)} = x(x-1)\cdots(x-k+1) \\

Stirling 1\textsuperscript{re} espèce --- récurrence &
  \left[\!\begin{smallmatrix} n \\ k \end{smallmatrix}\!\right]
  = (n-1)\left[\!\begin{smallmatrix} n-1 \\ k
  \end{smallmatrix}\!\right]
  + \left[\!\begin{smallmatrix} n-1 \\ k-1
  \end{smallmatrix}\!\right] \\

Factorielle montante &
  x^{(n)} = \sum_{k=0}^{n}
  \left[\!\begin{smallmatrix} n \\ k \end{smallmatrix}\!\right]
  (-1)^{n-k} x^k \\

\end{longtable}

\section{Autres identités remarquables}
\label{sec:autres-identites}

\renewcommand{\arraystretch}{1.8}
\begin{longtable}{@{}>{\raggedright\arraybackslash}p{5cm}
                    >{$\displaystyle}l<{$}@{}}
\toprule
\textbf{Nom / Identité} & \textbf{Formule} \\
\midrule
\endfirsthead
\toprule
\textbf{Nom / Identité} & \textbf{Formule} \\
\midrule
\endhead
\bottomrule
\endfoot

Inclusion-exclusion\index{inclusion-exclusion} &
  \card{A_1 \cup \cdots \cup A_n}
  = \sum_{k=1}^{n}(-1)^{k-1}\!\!\sum_{I \subseteq [n],\, \card{I}=k}
  \!\!\card{\bigcap_{i\in I} A_i} \\

Nombre de surjections\index{surjection} &
  S(n,k) \cdot k! = \sum_{j=0}^{k}(-1)^{k-j}\binom{k}{j} j^n \\

Dérangements\index{dérangement} &
  D_n = n!\sum_{k=0}^{n} \frac{(-1)^k}{k!} \\

Nombres de Catalan\index{nombre de Catalan} &
  C_n = \frac{1}{n+1}\binom{2n}{n}
  = \binom{2n}{n} - \binom{2n}{n+1} \\

Récurrence de Catalan &
  C_{n+1} = \sum_{i=0}^{n} C_i \, C_{n-i} \\

Formule de Cayley\index{formule de Cayley} &
  \text{Nombre d'arbres étiquetés sur } [n] = n^{n-2} \\

Multinomiale\index{coefficient multinomial} &
  (x_1+\cdots+x_k)^n = \sum_{\substack{n_1+\cdots+n_k=n \\ n_i \geq 0}}
  \binom{n}{n_1,\dots,n_k} \prod_{i=1}^k x_i^{n_i} \\

\end{longtable}

\bigskip
\begin{remarque}
Pour une table plus complète et de nombreuses preuves, on consultera avec
profit l'ouvrage de Graham, Knuth et Patashnik~\cite{GKP} ainsi que celui
de Cameron~\cite{Cameron}.
\end{remarque}


%%%%%%%%%%%%%%%%%%%%%%%%%%%%%%%%%%%%%%%%%%%%%%%%%%%%%%%%%%%%%%%%%%%%%
%%  BIBLIOGRAPHIE
%%%%%%%%%%%%%%%%%%%%%%%%%%%%%%%%%%%%%%%%%%%%%%%%%%%%%%%%%%%%%%%%%%%%%
\begin{thebibliography}{99}

\bibitem{Rosen}
\textsc{Rosen}, K.~H.,
\textit{Discrete Mathematics and Its Applications},
8\textsuperscript{e} édition, McGraw-Hill, 2019.
\index{Rosen}

\bibitem{GKP}
\textsc{Graham}, R.~L., \textsc{Knuth}, D.~E. et \textsc{Patashnik}, O.,
\textit{Concrete Mathematics: A Foundation for Computer Science},
2\textsuperscript{e} édition, Addison-Wesley, 1994.
\index{Graham, Knuth, Patashnik}

\bibitem{Lovasz}
\textsc{Lov\'asz}, L., \textsc{Pelik\'an}, J. et \textsc{Vesztergombi}, K.,
\textit{Discrete Mathematics: Elementary and Beyond},
Springer, 2003.
\index{Lovász}

\bibitem{Diestel}
\textsc{Diestel}, R.,
\textit{Graph Theory},
5\textsuperscript{e} édition, Graduate Texts in Mathematics~173,
Springer, 2017.
\index{Diestel}

\bibitem{Cameron}
\textsc{Cameron}, P.~J.,
\textit{Combinatorics: Topics, Techniques, Algorithms},
Cambridge University Press, 1994.
\index{Cameron}

\bibitem{vanLintWilson}
\textsc{van Lint}, J.~H. et \textsc{Wilson}, R.~M.,
\textit{A Course in Combinatorics},
2\textsuperscript{e} édition, Cambridge University Press, 2001.
\index{van Lint}\index{Wilson}

\bibitem{Grimaldi}
\textsc{Grimaldi}, R.~P.,
\textit{Discrete and Combinatorial Mathematics: An Applied Introduction},
5\textsuperscript{e} édition, Pearson, 2004.

\bibitem{Bondy}
\textsc{Bondy}, J.~A. et \textsc{Murty}, U.~S.~R.,
\textit{Graph Theory},
Graduate Texts in Mathematics~244, Springer, 2008.
\index{Bondy}\index{Murty}

\bibitem{HillCoding}
\textsc{Hill}, R.,
\textit{A First Course in Coding Theory},
Oxford University Press, 1986.

\bibitem{Roman}
\textsc{Roman}, S.,
\textit{Introduction to Coding and Information Theory},
Springer, 1997.

\end{thebibliography}


%%%%%%%%%%%%%%%%%%%%%%%%%%%%%%%%%%%%%%%%%%%%%%%%%%%%%%%%%%%%%%%%%%%%%
\printindex
\end{document}
